%% GEB LaTeX v2 — Part 02 / 7
%% 自动生成  split_v2.py
%% 独立编译（从 GEB_LaTeX_v2/ 目录）：
%%   xelatex -output-directory=split part02.tex

%% GEB 精排模板 ── XeLaTeX
\PassOptionsToPackage{unicode}{hyperref}
\PassOptionsToPackage{hyphens}{url}

%% ─────────────────────────────────────────────
%%  文档类：memoir（精排书籍首选）
%% ─────────────────────────────────────────────
\documentclass[
  12pt,
  twoside,
  openright,
  oldfontcommands,
]{memoir}

%% memoir 纸张尺寸（必须在字体等 usepackage 之前）
\setstocksize{260mm}{185mm}   % 大16开
\settrimmedsize{\stockheight}{\stockwidth}{*}
\settrims{0pt}{0pt}

%% ─────────────────────────────────────────────
%%  字体
%% ─────────────────────────────────────────────
\usepackage{fontspec}
\usepackage{xeCJK}

% ※（U+203B）Georgia 无此字符，强制用 CJK 字体渲染
\xeCJKDeclareCharClass{CJK}{"203B}

% geb.ttf：EPUB 自定义字体，含 Tai Viet 占位字形（U+AA9E–U+AAAF）
% 以及 U+31CF、U+39DF 等特殊字形
\newfontfamily\gebfont{geb.ttf}[Path=./fonts/]

% 中文正文：思源宋体（Noto Serif CJK SC）
\setCJKmainfont[
  BoldFont    = {Noto Serif CJK SC SemiBold},
  ItalicFont  = {Kaiti SC},          % 楷体作斜体
]{Noto Serif CJK SC}

% 中文无衬线（标题等）
\setCJKsansfont[
  BoldFont = {Heiti SC},
]{Heiti SC}

% 中文等宽
\setCJKmonofont{PingFang SC}

% 西文正文：Palatino（等高数字，基线与思源宋体对齐更佳）
\setmainfont[
  BoldFont       = {Palatino Bold},
  ItalicFont     = {Palatino Italic},
  BoldItalicFont = {Palatino Bold Italic},
]{Palatino}

% 西文无衬线
\setsansfont{Helvetica Neue}

% 数学字体
\usepackage{amsmath,amssymb}
\usepackage{unicode-math}

%% ─────────────────────────────────────────────
%%  版面：大16开（185×260mm）页边距（memoir 原生布局）
%% ─────────────────────────────────────────────
\setlrmarginsandblock{25mm}{20mm}{*}
\setulmarginsandblock{22mm}{25mm}{*}
\setheaderspaces{*}{7mm}{*}
\setheadfoot{\baselineskip}{12mm}
\checkandfixthelayout

%% ─────────────────────────────────────────────
%%  颜色
%% ─────────────────────────────────────────────
\usepackage{xcolor}
\definecolor{accent}{HTML}{1A3A5C}     % 深蓝（章号、页眉线）
\definecolor{accentlight}{HTML}{4A7DAA}% 中蓝（节标题）
\definecolor{rulegray}{HTML}{BBBBBB}   % 分隔线

%% ─────────────────────────────────────────────
%%  章节样式
%% ─────────────────────────────────────────────
\usepackage{titlesec}

% 章：大号宋体居中，上方装饰线，下方细实线
\titleformat{\chapter}[display]
  {\normalfont\centering}
  {\color{accent}\small\CJKfamily{}\sffamily 第\ \thechapter\ 章}
  {0pt}
  {\color{accent}\LARGE\bfseries\CJKfamily{}\rmfamily}
  [\vspace{4pt}\textcolor{rulegray}{\titlerule[0.4pt]}]

\titlespacing*{\chapter}{0pt}{30pt}{40pt}

% 节：宋体加粗左对齐，左侧蓝色竖线装饰
\titleformat{\section}
  {\normalfont\large\bfseries\color{accentlight}}
  {}
  {0pt}
  {\colorbox{white}{\raisebox{0pt}{\rule[-1pt]{3pt}{14pt}}\hspace{6pt}}}
  []

\titlespacing*{\section}{0pt}{18pt}{8pt}

% 小节
\titleformat{\subsection}
  {\normalfont\normalsize\bfseries}
  {}
  {0pt}
  {}

%% ─────────────────────────────────────────────
%%  页眉页脚（memoir 内建）
%% ─────────────────────────────────────────────
\makepagestyle{geb}
\makeevenhead{geb}{\small\thepage}{}{\small\itshape\leftmark}
\makeoddhead {geb}{\small\itshape\rightmark}{}{\small\thepage}
\makeheadrule {geb}{\textwidth}{0.4pt}
\pagestyle{geb}
\renewcommand{\chaptermark}[1]{\markboth{第\ \thechapter\ 章\enspace #1}{第\ \thechapter\ 章\enspace #1}}
\renewcommand{\sectionmark}[1]{\markright{\thesection\enspace #1}}

%% ─────────────────────────────────────────────
%%  中文标准名称
%% ─────────────────────────────────────────────
\renewcommand{\contentsname}{目录}
\renewcommand{\figurename}{图}
\renewcommand{\tablename}{表}
\renewcommand{\appendixname}{附录}
\renewcommand{\bibname}{参考文献}
\renewcommand{\indexname}{索引}

%% ─────────────────────────────────────────────
%%  段落间距与行距（memoir 内建命令）
%% ─────────────────────────────────────────────
\setSpacing{1.3}
\setlength{\parindent}{2em}
\setlength{\parskip}{3pt plus 1pt minus 1pt}
\raggedbottom

%% ─────────────────────────────────────────────
%%  图片与浮动体
%% ─────────────────────────────────────────────
\usepackage{graphicx}
\usepackage{float}
\usepackage{caption}
\captionsetup{
  font=small,
  labelfont={bf,color=accent},
  labelsep=quad,
  width=0.88\linewidth,
  justification=centering,
  singlelinecheck=false,
}

%% pandocbounded：memoir 兼容版，限制图片最大宽度和高度
\newcommand*\pandocbounded[1]{%
  \begingroup
  \setkeys{Gin}{width=\linewidth,height=0.8\textheight,keepaspectratio}%
  #1%
  \endgroup
}
\def\fps@figure{htbp}

%% ─────────────────────────────────────────────
%%  引用块样式
%% ─────────────────────────────────────────────
\usepackage{mdframed}
\mdfdefinestyle{gebquote}{
  leftline=true,
  rightline=false,
  topline=false,
  bottomline=false,
  linewidth=2pt,
  linecolor=accentlight,
  backgroundcolor=black!3,
  innerleftmargin=12pt,
  innerrightmargin=8pt,
  innertopmargin=6pt,
  innerbottommargin=6pt,
}
\renewenvironment{quote}
  {\begin{mdframed}[style=gebquote]\small}
  {\end{mdframed}}

%% ─────────────────────────────────────────────
%%  表格
%% ─────────────────────────────────────────────
\usepackage{longtable,booktabs,array,multirow,calc}
\usepackage{etoolbox}
\makeatletter
\patchcmd\longtable{\par}{\if@noskipsec\mbox{}\fi\par}{}{}
\makeatother
\IfFileExists{footnotehyper.sty}{\usepackage{footnotehyper}}{\usepackage{footnote}}
\makesavenoteenv{longtable}

% 美化 booktabs 规则
\setlength{\aboverulesep}{0pt}
\setlength{\belowrulesep}{0pt}
\setlength{\heavyrulewidth}{0.8pt}
\setlength{\lightrulewidth}{0.4pt}

%% ─────────────────────────────────────────────
%%  脚注
%% ─────────────────────────────────────────────
\usepackage[
  bottom,
  hang,
  flushmargin,
]{footmisc}
\setlength{\footnotemargin}{1.5em}
\renewcommand{\footnoterule}{%
  \vspace{4pt}\noindent\textcolor{rulegray}{\rule{0.3\linewidth}{0.4pt}}\vspace{4pt}%
}

%% ─────────────────────────────────────────────
%%  目录样式
%% ─────────────────────────────────────────────
\usepackage{titletoc}

%% memoir 的 chapternumberline/partnumberline 需映射到 \numberline
%% 才能被 titletoc 识别为「有编号」条目
\makeatletter
\renewcommand{\chapternumberline}[1]{\numberline{#1}}
\renewcommand{\partnumberline}[1]{\numberline{#1}}
\makeatother

\titlecontents{chapter}[0em]
  {\addvspace{6pt}\bfseries\color{accent}}
  {\makebox[3em][l]{\thecontentslabel}}
  {}
  {\hfill\contentspage}

\titlecontents{section}[3em]
  {\addvspace{2pt}\small}
  {}
  {}
  {\titlerule*[0.5em]{.}\contentspage}

%% ─────────────────────────────────────────────
%%  超链接
%% ─────────────────────────────────────────────
\usepackage{bookmark}
\IfFileExists{xurl.sty}{\usepackage{xurl}}{}
\urlstyle{same}
\usepackage{hyperref}
\hypersetup{
  pdftitle={哥德尔、艾舍尔、巴赫——集异璧之大成},
  pdfauthor={〔美〕侯世达},
  pdflang={zh-CN},
  colorlinks=true,
  linkcolor={accent},
  urlcolor={accentlight},
  citecolor={accent},
  pdfcreator={XeLaTeX + GEB Template},
}

%% ─────────────────────────────────────────────
%%  杂项
%% ─────────────────────────────────────────────
\usepackage{microtype}
\setlength{\emergencystretch}{3em}
\providecommand{\tightlist}{%
  \setlength{\itemsep}{0pt}\setlength{\parskip}{0pt}}

% 对话角色加粗
\newcommand{\speaker}[1]{\textbf{#1}}

%% ─────────────────────────────────────────────
%%  GEB v2 字体命令
%% ─────────────────────────────────────────────
\newCJKfontfamily\kaitifamily[BoldFont={Kaiti SC Bold}]{Kaiti SC}
\newcommand{\kaititext}[1]{{\kaitifamily #1}}
\newcommand{\kaitismall}[1]{{\kaitifamily\small #1}}
\newcommand{\kaitibold}[1]{{\kaitifamily\bfseries #1}}
\newcommand{\heititext}[1]{{\sffamily\bfseries #1}}
\newcommand{\heitismall}[1]{{\sffamily\bfseries\small #1}}
\newcommand{\dialogname}[1]{{\sffamily\bfseries #1}}
\newcommand{\westerntext}[1]{{\rmfamily #1}}
\newcommand{\westernital}[1]{{\rmfamily\itshape #1}}

%% ─────────────────────────────────────────────
%%  GEB v2 排版环境（memoir 内建 hangparas，不需要 hanging 宏包）
%% ─────────────────────────────────────────────

% 对话悬挂缩进行
\newenvironment{hangdialog}{%
  \begin{hangparas}{2em}{1}%
}{%
  \par\end{hangparas}%
}

% 楷体舞台指示（居中小楷）
\newenvironment{stagedir}{%
  \begin{center}\kaitifamily\small%
}{%
  \end{center}%
}

% 对话引导语（楷体斜体缩进）
\newenvironment{dialogguide}{%
  \begin{adjustwidth}{1em}{0em}\kaitifamily\itshape\small%
}{%
  \end{adjustwidth}%
}

% 仿宋/楷体引用块（带左竖线）
\newenvironment{fsquote}{%
  \begin{mdframed}[style=gebquote]\kaitifamily\small%
}{%
  \end{mdframed}%
}

% 楷体加粗引用框
\newenvironment{ktquote}{%
  \begin{mdframed}[style=gebquote]\kaitifamily\small%
}{%
  \end{mdframed}%
}

% 禅宗公案块
\newenvironment{zenkoan}{%
  \begin{mdframed}[style=gebquote]\kaitifamily\small%
}{%
  \end{mdframed}%
}

% 禅宗公案诗
\newenvironment{zenkoanpoem}{%
  \begin{center}\kaitifamily\small%
}{%
  \end{center}%
}

% 图片副说明（小楷居中，灰色）
\newenvironment{imgsub}{%
  \begin{center}\kaitifamily\small\color{rulegray}%
}{%
  \end{center}%
}

% 对话诗歌
\newenvironment{dialogpoem}{%
  \begin{adjustwidth}{3em}{0em}\kaitifamily%
}{%
  \end{adjustwidth}%
}

% 章节前置数字大标题（如"第一章"）
\newenvironment{chappreheader}{%
  \begin{center}\large\sffamily\bfseries\color{accent}%
}{%
  \end{center}\vspace{-10pt}%
}

% 部分前置大标题
\newenvironment{partpreheader}{%
  \begin{center}\Large\sffamily\bfseries\color{accent}%
}{%
  \end{center}\vspace{-10pt}%
}

%% ─────────────────────────────────────────────
%%  扉页（自定义，替换 maketitle）
%% ─────────────────────────────────────────────
\newcommand{\gebmaketitle}{%
  \begin{titlingpage}
    \vspace*{6\baselineskip}
    \begin{center}
      {\fontsize{28}{34}\selectfont\bfseries\color{accent}
        哥德尔、艾舍尔、巴赫}\\[12pt]
      {\Large\color{rulegray}——集异璧之大成}\\[40pt]
      \textcolor{rulegray}{\rule{0.4\linewidth}{0.4pt}}\\[20pt]
      {\large \textit{Gödel, Escher, Bach: An Eternal Golden Braid}}\\[6pt]
      {\normalsize \color{accentlight} Douglas R.\ Hofstadter}\\[40pt]
      {\normalsize 〔美〕侯世达　著\\[4pt]
        本书翻译组　译}\\[60pt]
      {\small\color{rulegray} 商务印书馆}
    \end{center}
    \vfill
  \end{titlingpage}
}

%% pandoc 在 body 中会注入 \babel@toc，提供空定义避免报错
\makeatletter
\providecommand{\babel@toc}[2]{}
\makeatother

%% memoir + XeLaTeX 下需手动告知 xdvipdfmx 纸张尺寸
\AtBeginDocument{%
  \special{papersize=\the\paperwidth,\the\paperheight}%
}

%% ═══════════════════════════════════════════
%%  正文开始
%% ═══════════════════════════════════════════

\begin{document}
\mainmatter
\setcounter{chapter}{8}
\markboth{}{}

\chapter{WU谜题}

形式系统

本书中心的概念之一是形式系统。我用的这种形式系统是美国逻辑学家艾米尔·波斯特在二十世纪二十年代发明的，通常被称为``波斯特产生式系统''。本章将向读者介绍一个形式系统，并且我希望读者愿意------至少是有一点点愿意------去探索这个形式系统。为唤起读者的好奇心，我设计了一个小小的谜题。

这个谜题就是``你能产生WU吗？''。一开始的时候，我提供给你一个符号串（即一串字母）。为了解除你的悬念，可以告诉你那个符号串就是WJ。然后，你会得知一些规则，运用这些规则你可以将一个符号串变成另一个。如果某条规则在某处是适用的，并且你也愿意用它，你就可以用。但是，如果同时有几条可以适用的规则，又没有什么指导来说明你该用哪一条，这将留给你去选择。正是这一点使得做形式系统游戏成为一种技巧。重要的、几乎用不着说明的一点是，你决不能做任何背离规则的事情。我们也许可以把这个限制称做``形式化要求''。在本章中，这一点也许根本无需强调。但是我预言------这说起来挺怪的------当你在下面的章节中与形式系统周旋时，你会发现，你将一次又一次地违反形式化要求，除非你以前用过形式系统。

关于我们的形式系统------WJU系统------要说的第一件事是它只用了三个字母：W、J、U。就是说WJU系统的所有符号串都是由这三个字母组成的。以下是WJU系统的一些符号串。

\begin{quote}
WU

UJW

WUUWUU

UJJUWJUUJWUJJUWJUUJWUJJU
\end{quote}

虽然所有这些符号串都是合法的，但它们还不算``归你所有''。事实上，到现在为止归你所有的符号串只有WJ。要增加你自己的储备，只能靠使用规则，我们现在就介绍这些规则。下面是第一条规则：

规则Ⅰ：如果一个归你所有的符号串结尾是J，则可以在其后面再加上一个U。

这里顺便指出一个事实（假如你到目前还未猜出``符号串''的意思）：符号串中的字母是有固定顺序的。比如WJ和JW是两个不同的符号串。一个符号串不是其字母的顺序无关紧要的``一堆符号''。

以下是第二条规则：

规则Ⅱ：如果你有Wx，那么Wxx也归你所有。

我的意思可在下面几个例子里表示出来：

\begin{quote}
从WJU，可以得出WJUJU

从WUW，可以得出WUWUW

从WU，可以得出WUU
\end{quote}

就是说，在这条规则中字母``x''可以代表任意的符号串。但是，一旦决定了它代表哪一个符号串，就得遵循这个选择（直到你再一次使用这条规则，才可以做一个新的选择）。请注意上面第三个例子。它表明如何在有了WU之后，为你的储备增加另一个符号串。但是得先有WU！关于字母``x''，我想再加上最后一条说明：它不像``W''、``J''和``U''这三个字母一样，是形式系统的一部分。无论如何，如果我们能有某种符号的方式一般性地谈论系统中的符号串，那将是很有用的------而``x''正是起了这个作用：代表任意一个符号串。假如在你的符号串``储备''中加入一个含有``x''的符号串，那就错了，因为WJU系统中的符号串决不会带有``x''！

以下是第三条规则：

规则Ⅲ：如果JJJ出现在你的储备中的一个符号串里，那么你可以用U代替这个JJJ而得到一个新的符号串。

例子：

\begin{quote}
从UWJJJWU，可以得到UWUWU

从WJJJJ，可以得到WJU（同样可以得到WUJ）

从JJWJJ，用这条规则什么也得不到（三个J必须是连续的）

从WJJJ，可以得到WU
\end{quote}

在任何情况下，这条规则都不能倒过来使用，像下面的例子里这样：

\begin{quote}
从WU，得到WJJJ。　\gebfont{⇐}这是错误的。
\end{quote}

规则是单向的。

下面是最后一条规则：

规则Ⅳ：如果UU出现在你的一个符号串中，你可以去掉它。

\begin{quote}
从UUU，可以得U

从WUUUJJJ，可以得到WUJJJ
\end{quote}

规则就是这些。现在可以试试找出WU了。假如得不出来不要担心。只是试一下------主要是体会一下WU谜题的味道。愿读者从中获得享受。

定理、公理、规则

WU谜题的答案后面会给出来。至于现在，重要的不是找出答案，而是寻找它。读者也许已经做了一些产生WU的尝试。在做这些尝试的过程中，读者建立了自己的符号串储备。我们把这种可以从规则中产生的符号串叫做定理。当然，``定理''这个词在数学中也使用，但与这里的意思很不一样。那里的意思是普通语言中的某个陈述，已通过严格的论证被证明是真的（这时我们称之为``真理''），例如，关于``运动无有''的\heititext{芝诺定理}，或者关于素数有无穷多个的\heititext{欧几里德定理}。但是，在形式系统中，不必把定理看做是陈述------它们仅仅是一些由符号组成的串。它们也不是证明出来的，而是产生出来的，就像是按照一定的印刷规则通过机器产生出来的一样。为了强调``定理''这个词在意义上的重要区别，在本书中我将采用以下约定：当``定理''这个词用黑体字写出来时，它的意思是我们通常所说的意思------一个\heititext{定理}是日常语言中的一个陈述，这个陈述已被某人用某种逻辑的论证证明是真理。如果不用黑体则``定理''这个词有它的技术上的意义：可以在某个形式系统中产生出来的符号串。在这样的情况下，WU谜题是问WU是不是WJU系统中的一个定理。

开始的时候我无偿地提供了一个定理，即WJ。这个``无偿提供的''定理叫做``公理''------同样的，这是技术上的意义，它与通常的意义很不一样。一个形式系统中可以有零个、一个、几个、甚至无穷多个公理。本书中将会出现所有这些不同类型的形式系统。

每一个形式系统中都有符号变换规则，如WJU系统中的四条规则。这些规则既可以叫做生成规则，也可以叫做推理规则。这两种叫法我都用。

我在这里要介绍的最后一个术语是推导。以下是一个定理WUJJU的推导：

\begin{longtable}[]{@{}lll@{}}
\toprule\noalign{}
\endhead
\bottomrule\noalign{}
\endlastfoot
⑴ & WJ & 公理 \\
⑵ & WJJ & 使用规则Ⅱ，从⑴中得出 \\
⑶ & WJJJJ & 使用规则Ⅱ，从⑵中得出 \\
⑷ & WJJJJU & 使用规则Ⅰ，从⑶中得出 \\
⑸ & WUJU & 使用规则Ⅲ，从⑷中得出 \\
⑹ & WUJUUJU & 使用规则Ⅱ，从⑸中得出 \\
⑺ & WUJJU & 使用规则Ⅳ，从⑹中得出 \\
\end{longtable}

一个定理的推导说明如何按照这个形式系统中的规则逐行生成该定理的一个明显的论证。``推导''这一概念是向``证明''概念看齐的，但比后者要朴素一些。如果说证明出了WUJJU，听起来就很奇怪，但是如果说推导出了WUJJU，就不那么奇怪了。

系统内外

多数人解WU谜题的办法是：先相当盲目地推出一些定理，看一看得到的会是什么。很快地，他们就开始注意到他们产生出的定理的一些性质，人的智能就在此处起作用了。例如，也许你在试开头的几次时，并不清楚所有的定理都是以W开头的。推了一阵以后，这个模式浮现在你脑海中。你不仅可以看出这个模式，而且通过检查那些规则，还可以理解这个模式。那些规则有一个性质，即：新产生的定理都继承了原先的那条定理的第一个字母。归根到底，所有的定理的第一个字母都可以追溯到唯一的公理WJ的第一个字母------这是一个证明，指出WJU系统中所有的定理都必须以W开头。

这里发生的情况是有着重要意义的。它展示了人与机器之间的一种区别。非常可能------事实上很容易------给一部计算机编一个程序，使它可以生成WJU系统的一条又一条定理，并且我们可以在这个程序中加进一条命令，使之在生成U之后才停止。现在你知道了：运行这个程序的计算机会永不停止，而这并不使你感到惊异。但是，假如你要让你的一位朋友试图生成U，将会怎么样呢？要是他做了一会儿就罢手不干了，跑来向你抱怨他无法摆脱第一个字母W，而想生成U完全是徒劳的，你同样不会吃惊。即使是一个不很聪明的人，也会禁不住地观察他自己在做什么，这些观察给予他关于所做的这项工作的洞察力。而计算机程序------正如我们所描述过的------则缺乏这种洞察力。

现在让我更清楚地谈谈我说``这展示出了人与机器之间的区别''是什么意思。我的意思是：有可能给机器编一个以上述方式来做机械工作的程序，机器永远不会发现（即便这是非常明显的）它自己在做什么。但是人的意识有一个固有的特点，他总是能看出关于他正在做的事情的某些事实。关于这一点，你是早已知道的。假如你在加法机上按一下``1''，然后再加1，再加1，再加1，再加1\ldots\ldots 如此一直加上几个小时又几个小时，虽然任何人都会很快地懂得自己在重复一个动作，而那台机器却永远不懂得在你按1之前，自己去加。或者，另举一个愚蠢点儿的例子：无论用了多久，无论多好的一部汽车，也永远不会有躲避道路上的其它汽车或障碍物的能力，永远不会知道它的主人常走的路线。

因此，区别在于：机器有可能在做某件事情时不去观察，而人不可能不去观察。请注意，我不是说所有的机器都不具有复杂的观察能力，只是有些机器不具有。我也不是说所有的人都总是进行复杂的观察，事实上，人常常是不注意观察的。但是机器可以被制造成是完全不去观察的，而人则不能。并且事实上，到目前为止，多数机器是非常接近于不去观察的。也许因为这个原因，对于大多数人来说，具有``不去观察''这一特点，似乎是机器的特征。例如，如果有人说某项工作是``机械的''，这指的不是人没有做这项工作的能力，而是隐含着这样的意思：只有机器可以永不抱怨或者不知疲倦地反复做这项工作。

跳出系统

能够跳出正在进行的工作并且看一下已经做了些什么，这是智能固有的特点。它总是寻找并且常常能找到模式。我说的是智能可以跳出它所做的工作，这不意味着它总是这样做。然而，常常是略微推动它一下就足够了。就好比一个正在读书而感觉困倦的人，他不是连续地把这本书看完，而可能把书搁在一边儿，把灯关掉。他``出了系统''，而这对我们来说似乎是再自然不过的了。或者，假设一个人，A，正在看电视，B进了屋子，并且明显地表示出对于当时的状况不喜欢。A可能认为他理解了问题的所在，并且试图通过从当前系统（那时的电视节目）退出来改变一下现状，于是他轻轻按了一下频道钮，找一个好一点儿的节目。不过B可能对于什么是``退出系统''有更极端的概念------也就是，把电视机关掉！当然也有这种情况：只有极少数的人有那种眼光，看出一个支配着许多人生活的系统，而以前却从没人认为这是一个系统。然后这些人常常就投入毕生的精力去使其他人相信那个系统确实存在，而且应该从中退出！

若想教给计算机跳出系统这件事，能教会它到什么程度呢？我举一个使一些人吃惊的实例：不久以前，在加拿大举行的一次计算机国际象棋邀请赛上，一个程序------所有竞争者中最弱的一个------有一个非同一般的特性：它可以远在棋下完之前早早退出。它的棋下得不太好，但它至少有预先估计到没有希望的棋局的能力，然后就此停下，而不是等着另一个程序走完将死它的乏味的过程。虽然每次下棋它都输，但是它做得很有风度。它给很多当地的下棋能手留下了深刻的印象。假如把``系统''定义成``在下棋的游戏中走棋的系统''，那么，很清楚，这个程序有一个很复杂的事先编好的有能力从系统中退出的程序。另一方面，如果把``系统''定义成``这部计算机按照程序所做的一切''，那么无疑地，这部计算机没有什么从那个系统中退出的能力。

在研究形式系统时，很重要的一点是要区分在系统之内的工作，和做出关于系统的判断或说明。我假定读者像大多数人一样，开始时，想通过在系统内部工作解决WU谜题；逐渐地，你开始着急了，这种焦急心情增长到一定程度时就不再作任何考虑地退出了这个系统，然后盘算一下产生了一些什么，想想为什么没有取得成功------即产生出WU。也许你发现了没能产生出WU的原因所在，这是对于系统进行思索。也许你在中途什么时候产生了WJU，这是在系统之内进行工作。我不想把事情弄得好像两种方式完全不相容，我敢肯定任何一个人都有某种程度的在系统之内进行工作，同时思考他正在做什么的能力。实际上，在人类事务中，通常是几乎不可能将``在系统之内''和``在系统之外''清楚地区别开的，生活是由许多连结并交织又常常是不协调的``系统''组成的，用``系统内''、``系统外''这类词汇来思考似乎是过于简单化了。但是，非常清楚地形成一些简单的概念，以便人们可以在思考更复杂的概念时把它们作为模型，这常常是重要的。这便是我为什么要向读者展示形式系统的原因。现在是回头讨论WJU系统的时候了。

W方式、J方式、U方式

WU谜题的陈述方式鼓励人们在WJU系统内部进行探索------做一些定理推导。但这个问题的陈述方式并不暗示在系统之内必然会产生结果。因此它使得人们往返于两种方式之间。一个把两种方式区分开的方法可以是：准备两张纸，在一张纸上，你``像机器那样''地工作，这样在纸上写满了的都是W、J、U这三个字母。在另一张纸上，你``像有思维能力的人那样''地工作，你可以做你的智能要你做的任何事情------可以使用汉语、记下要点、进行逆推、使用缩写（比如使用字母``X''）、把几个步骤缩成一步、把规则加以修改看看可以得出什么东西，以及做你所能想象到的任何其它事情。要注意3和2这两个数起着很重要的作用，因为，三个J是可以扔掉的，两个U也可以去掉。另外根据规则Ⅱ，可以把长度延长一倍（除了那个W）。所以第二张纸上也可以有一些计算。我们会时常重新提到处理形式系统的这两种方式，我们分别称之为机方式（J方式）和惟方式（W方式）。为了完满，使WJU系统的每一个字母都与一种方式相关联，我再提出最后一种方式------无方式（U方式），它是禅宗的处世态度。关于这点在后面几章中会更多地讲到。

判定过程

对于这个谜题的观察表明，它涉及到两种具有不同趋向的规则------加长的规则和缩短的规则。其中两条规则（规则Ⅰ和Ⅱ）允许你加长符号串（当然只能是非常刻板地按预先规定的方式来加长），另外两条规则允许多少缩短一点符号串（同样是以非常刻板的方式）。看起来应该是有无限多种不同的顺序来应用这些不同种类的规则，这就使得用这种或那种方式来产生WU有了希望。可能是使符号串加长成一个庞然大物，然后再一段一段地抽掉，直到只剩下两个符号为止；或者更糟糕，使符号串连续地加长、缩短、再加长、再缩短，如此下去，而结果如何却没有保证。事实上，我们已经看出U是不可能产生的，即使是加长或缩短到世界末日也无济于事。

然而，U和WU似乎很不相同。通过U的非常表面的特征，我们可以辨别出不可能产生它：它不是以W开头的（而所有的定理都必须是）。用这样一个简单的方法去发现非定理是很方便的。然而，能说这个测试可以发现出所有的非定理吗？可能会有许多以W开头但无法产生的符号串。也许WU就是其中之一。这也就是说``字头测试''的用途是有限的，它只能发现一部分非定理，而漏掉其它的。但是仍然有可能有一些更精细的测试，这些测试完全能区分哪些符号串通过规则可以产生，哪些不能产生。这里我们得面对一个问题：``我们的测试是什么意思？''目前这个问题的意义或者重要性也许不明显。下面我举一个关于``测试''的例子，某种意义上似乎有点违反这个词的精神：

设想在这个世界上有一个怪物，它有足够的时间，并且喜欢把时间用来以有秩序的方法去产生WJU系统的定理。例如，下面就是这个怪物可能采取的一种方式：

步骤1：对于公理WJ应用每一条可以应用的规则。这产生出两个新的定理：WJU、WJJ。

步骤2：对于从步骤1中产生的定理应用所有可以应用的规则。这产生出三个新的定理：WJJU、WJUJU和WJJJJ。

步骤3：对于步骤2中产生的定理应用所有可以应用的规则。这产生出五个新的定理：WJJJJU、WJJUJJU、WJUJUJUJU、WJJJJJJJJ和WUJ。

\gebfont{\gebfont{}

这种方法早晚会产生出所有的定理，因为每条规则都在任意能设想出的次序下应用过。（参看图11）我们以上提到的所有那些加长------缩短的交替用法最终都用上了。然而对于一个给定的符号串，就不知道要经过多少步它才能出现在这张表里，因为定理是按照它们的推导长度来列出的。这不是一个很有用的次序，假如你对一个特殊的符号串（像WU）感兴趣的话，你甚至不知道是否有个推导能得出它，更不知道那个推导会有多长了。

\pandocbounded{\includegraphics[keepaspectratio]{./media/OEBPS/Images/img_011.png}}

图11．WJU系统中所有定理的``树''

在这个系统中向下到第N步，它包含了恰恰需要N个步骤推导的那些定理。圆圈中的数字表示使用了哪条规则。WU是不是在这棵树上的什么地方？

现在我们给出这个所谓的``定理性测试''：

等待产生出所要求的符号串，当它出现了的时候，你知道它是一个定理------假如它永远不出现，它就不是一个定理。

这看起来是荒唐的，因为它预先假设了我们为了得到一个答案，并不在乎等待哪怕是无限长的时间。这就是该把什么算作``测试''这个问题的关键。最首要的问题是保证我们将在一个有限长的时间内得到答案。如果有一个检验定理的测试，一个总是在有限长的时间内终结的测试，那么，这个测试就叫做给定形式系统的判定过程。

当你有了一个判定过程，那么你就有了对系统中所有定理的性质的非常具体的刻划。从表面上看，似乎形式系统中的规则与公理同判定过程一样完全地刻划了这个系统中的定理。问题在于``刻划''这个词。当然那些推理规则与WJU系统中的公理是隐含地刻划了那些是定理的符号串。不仅如此，它们还以更隐蔽的方式刻划了那些非定理的符号串。但是，在许多情况下，隐含地刻划还不够。如果有人声称有了一个对所有定理的刻划，但是若要推断某一个特别的符号串不是一个定理则需要无限长的时间，你大概会倾向于说，这个刻划中缺少点什么------它不够具体。这就是为什么发现判定过程的存在是非常重要的一步。所谓发现，从效果上看就是：你可以为检验一个符号串是不是一个定理进行一种测试，即便这个测试很复杂，它也是保证有结果的。原则上，这个测试与检查符号串的第一个字母是不是W同样地容易，同样地机械，同样地有限，同样是十分确定的。一个判定过程是测试一个符号串是否是定理的``试纸''。

顺便说一下，对形式系统的一个要求是，其中的公理必须具有判定过程------必须存在有关于一个符号串是不是一条公理的试纸。这就保证了至少在一开始进行的时候没有问题。这就是一组公理与一组定理之间的区别：前者总是有一个判定过程的，后者则可以没有。

我敢保证你会同意，当你第一次看见WJU系统时，你必须得面对这个问题。单独一个公理是已知的，推理规则是简单的，因此定理被隐含地刻划了------然而它仍然不很清楚，你不知道这个刻划的结果是什么。尤其是，WU究竟是不是一个定理，这一点仍然不清楚。

\pandocbounded{\includegraphics[keepaspectratio]{./media/OEBPS/Images/img_012.jpg}}

图12．空中城堡

艾舍尔作（木刻，1928）

\phantomsection\label{Dialog02.xhtml}{}

\chapter{二部创意曲}

或，

乌龟说给阿基里斯的话

刘易斯·卡罗尔{作}\footnote{}

阿基里斯追上乌龟之后，舒舒服服地坐在了龟背上。

``这么说，你已经到达我们这场赛跑的终点了，是吗？''乌龟说。``虽然比赛的路程是由无数段路程组成的，你还是跑到啦？我记得有个自作聪明的人证明过这是不可能做到的。''

``这是可以做到的，''阿基里斯说。``这已经做到了！马到成功。你瞧，那些路程在不断缩短，因此------''

``可要是它们一直在不断增长呢？''乌龟打断他的话说。``那会怎么样？''

``那样一来我现在就不会在这儿了，''阿基里斯老老实实地说。``而你到这时应该已经绕地球好几圈了！''

``你对我真是过奖了------哦，我是说你真是过重了，''乌龟说。``因为你的确太沉了！哎，你想听听有关另一场比赛的事吗？在这场比赛里大多数人都以为他们在两三步之内就能达到终点，实际上这场比赛也是由无数段路程组成的，其中每一段都比它前面那段要长些。''

``太想知道了！''那位希腊勇士一边说，一边从他的头盔里（那时几乎没有哪个希腊勇士身上有口袋）拿出一个硕大的笔记本和一支铅笔。``开始吧！请慢点说！速记现在还没发明出来呢！''\,``那个漂亮的欧几里德第一命题！''完全沉浸进去的乌龟喃喃地说。``你钦佩欧几里德吗？''

``钦佩得五体投地！至少，对于一部未来几个世纪之后才将问世的著作来说，我钦佩它的程度是无与伦比的了！''

``那好，让我们先来看看第一命题的论证中一个小小的片断------仅仅两步------以及由此得出的结论。请把它们记在你的本上。为了谈到它们时方便起见，我们将它们称作A、B、Z：------

（A）同等于一物的彼此亦相等。

（B）这个三角形的两条边同等于一物。

（Z）这个三角形的两条边彼此相等。

我想欧几里德的读者们会认为Z是A和B的合乎逻辑的推论，所以任何人只要认为A和B为真，则必定认为Z也真，不是吗？''

``毫无疑问！连上初中的毛孩子------等到两千多年以后发明出初中来------也会对此表示同意的。''

``如果有某个读者不认为A和B为真，我想他也还会认为这一推论是有效的，对吗？''

``这种读者肯定有。他会说，`我同意下述假言命题是真的：如果A和B为真，则Z必为真。但是我不接受A和B是真的。'这种读者应该明智点儿，放弃欧几里德，改踢足球去。''

``还可能有读者会说`我同意A和B为真，但不接受那个假言命题'，这也有可能吧？''

``当然有这种可能。他最好也去踢足球。''

``迄今为止，''乌龟继续说，``这两种读者都没有接受Z为真的逻辑必然性，对吗？''

``没错，''阿基里斯赞同地说。

``嗯，那么，我想让你把我当作是那第二种读者，用逻辑来迫使我接受Z为真。''

``一只会踢足球的乌龟会是------''阿基里斯有点吞吞吐吐。

``------当然是挺反常的，''乌龟急忙抢过话头。``你先别往旁边岔。咱们先谈Z，后谈球！''

``我得迫使你接受Z，是吗？''阿基里斯若有所思地说。``你现在的观点是只接受A和B，而不接受假言------''

``把它称作C吧，''乌龟说。

``------而不接受

（C）如果A和B为真，Z必为真。''

``这就是我目前的观点，''乌龟说。

``那么我得说服你必须要接受C。''

``只要你一把它记在你的本子上，我马上就接受它。''乌龟说，``你在上面还记了些什么？''

``只是一些摘要，''阿基里斯一边说，一边紧张地翻着那个笔记本：``是些摘要------是关于使我出了名的那些场战斗的！''

``我看到有好多空白页！''乌龟叫道，这一发现使他很高兴。``它们都会用得上的！''\kaitismall{（阿基里斯颤了一下。）}``现在，我说你写：------

（A）同等于一物的彼此亦相等。

（B）这个三角形的两条边同等于一物。

（C）如果A和B为真，则Z必为真。

（Z）这个三角形的两条边彼此相等。''

``你应该把最后这个命题Z改称为D，''阿基里斯说。``它紧承着上面那三个命题而来。如果你接受A、B、C，你就必须接受Z。''

``为什么我必须接受？''

``因为它是前三个的合乎逻辑的推论。如果A、B、C为真，则Z必为真。我看这是无可争辩的吧？''

``如果A、B、C为真，则Z必为真，''乌龟若有所思地重复着。``这又是一个假言判断，对不对？如果我不觉得这一假言判断是真的，那么我可以接受A、B、C，而仍然不接受Z，对吗？''

``你可以，''这位正直的英雄承认道。``可头脑这么迟钝的人也太罕见了。不过，这事儿还是可能的。所以我必须再使你接受一个假言判断。''

``好极了，我很愿意接受，只要你把它写下来。我们可把它称作

（D）如果A和B和C为真，则Z必为真。

你把它记到你的本上了吗？''

``已经记上了！''阿基里斯快活地宣布，同时把铅笔插进笔帽里。``我们终于达到了这场思想竞赛的终点！你现在接受了A、B、C、D，你当然要接受Z。''

``是吗？''乌龟装傻充楞地说。``让我们搞搞清楚。我接受了A、B、C、D，请设想一下，我依然拒绝接受Z，那会怎么样？''

``那样逻辑就会掐着你的脖子，迫使你接受！''阿基里斯得意洋洋地回答道。``逻辑会告诉你，`这事你作不了主。你既然已经接受了A、B、C、D，你就必须接受Z！'因此你别无选择，明白了吗？''

``这美妙的逻辑提供给我们的一切，都值得记下来，''乌龟说。``请把它记在你的本上吧。我们将把它称作

（E）如果A、B、C、D为真，则Z必为真。

在我接受它之前，我当然不一定要接受Z。因此它是很必要的一步，你明白吗？''

``我明白，''阿基里斯说。他的声音里带着点悲哀。

到这里，叙述者由于家里有急事，不得不离开这可爱的一对儿了。直到几个月以后，他才又路过这个地方。这时阿基里斯还坐在耐心的乌龟的背壳上，往他那几乎写满了的本子上写个不停。乌龟这时说：``你把刚才那步记下了吗？我要是没数错，已经有一千零一步了。将来还会有亿万步呢。算是帮我个忙，你不会介意考虑考虑我们两人的这篇对话究竟给了十九世纪的逻辑学家多少教益吧------你不会介意利用一下我表兄假海龟到那时将会发明出来的那个双关语，把你的名字改成`悟诡'吧？''

``随你便，''这位精疲力竭的勇士把脸埋在双手里回答说，他的声音因绝望而变得很空洞。``只要你愿意利用假海龟从未发明过的双关语，把你的名字改成`厌极易死'！''

\begin{enumerate}
\item
  \phantomsection\label{Dialog02.xhtml_B_1}
\end{enumerate}

\phantomsection\label{Chapter02.xhtml}{}

第二章

\chapter{数学中的意义与形式}

就是这篇《二部创意曲》启发了我，使我构想出本书对话中那两个角色。就像刘易斯·卡罗尔信手借来芝诺的乌龟和阿基里斯一样，我也信手借来刘易斯·卡罗尔的乌龟和阿基里斯。在卡罗尔的对话中，同样的事件一次又一次地发生，只是每一次都发生在更高的层次上，它与巴赫的``无穷升高的卡农''绝妙地相似。即便不谈其中的妙语，卡罗尔的对话也仍然包含有深刻的哲学问题，即：文字和思维是否遵循形式规则？这个问题也是本书的问题。

在本章与下一章中，我们将研究一些新的形式系统。这将使我们对形式系统这一概念有一个更为广阔的视野。读完这两章之后，对于形式系统的力量，以及为什么它们引起了数学家和逻辑学家的兴趣，读者就会有一个相当清楚的认识了。

pq系统

本章中的形式系统叫做pq系统。它对于数学家和逻辑学家并不重要------事实上，它只是我的一个简单发明。它的重要性，实际上仅在于它为在本书中起着相当重要作用的许多概念提供了一个非常好的例子。pq系统有三个不同的符号：

\textbf{p　q　-}

------字母p、q和短杠。

pq系统有无穷多条公理。由于不能把它们全部写出来，我们只能用另外的方法描述它们。实际上我们不仅需要对于公理的描述，我们还需要一种辨别某个给出的符号串是否是一条公理的方法。对于公理仅仅做一个描述也许能充分地刻划它们，但是这种刻划是很弱的------与WJU系统中对于定理的刻划所存在的问题是一样的。我们并不想仅仅为确定某个符号串是否是公理而花费一段不知多长的------甚至可能是无限长的------时间。所以，我们将给公理这样下定义，使得对于由p、q和短杠所组成的符号串是否是一条公理，有一个明显的判定过程。

定义：只要x仅由一串短杠组成，那么x-qxp-就是一条公理。

请注意``x''在两次出现时必须是代表同一串短杠。比如：-\/-\/-q-\/-p-是一条公理。当然``x-qxp-''这个表达式本身并不是一条公理（因为``x''不属于pq系统）。它更像是一个铸出所有公理的模子------它被称做公理模式。

pq系统只有一条生成规则：

规则：假设x、y和z都代表只包含短杠的特定的符号串，并且假设xqypz是一条已知的定理，那么x-qypz-就是一条定理。

比如让x是``-''，y是``-\/-''，z是``-\/-\/-''，这条规则就是：

如果-q-\/-p-\/-\/-是一条定理，则-\/-q-\/-p-\/-\/-\/-也是一条定理。

正像生成规则通常的形式那样，这个陈述在一个符号串是否是定理与另一个符号串是否是定理这两者之间建立了因果关系，但并不断定这些符号串本身是否是定理。

对于读者最有用的练习就是为pq系统的定理找出一个判定过程。这并不困难，找一会儿就能找出来。不妨试一试。

判定过程

我假定读者已经试过了。首先我想指出------这也许是太显而易见而不必提起的------pq系统中的每一条定理都有三组分离的短杠，并且起分离作用的成分依次为q，p（这可以由基于``继承性''的一个论证来证明，就是证明WJU系统中的定理都是以W开头的那种方法）。这就是说，仅仅从形式上，我们就可以排除像-\/-\/-\/-\/-\/-\/-\/-q-\/-p-\/-p-\/-p-\/-这样的符号串，指出它不是一条定理。

强调``仅仅从形式上''这几个字似乎显得是笨拙了：除了形式，一个符号串还有什么？有什么其它的东西可能在确定它的性质时起作用？很显然，没有。但是在继续讨论形式系统时，请把这一点记在心里：``形式''的概念会变得更复杂、更抽象。所以我们恐怕得多想一下``形式''这个词的意义。不管情况怎么样，让我们把任何一个以一组短杠开头，然后有一个q，接着是第二组短杠，然后是p，最后是另一组短杠这样的符号串都叫作``构形良好的''（简称``良构''）符号串。

让我们回到判定过程上来。是否是定理的标准，是后两组短杠加起来是否等于第一组短杠。例如：-\/-\/-\/-q-\/-p-\/-是一条定理，因为4等于2加2，而-q-\/-p-\/-不是一条定理，因为1不等于2加2。要知道为什么这是正确的标准，首先要看一下公理模式。显然，它只制造满足这个加法标准的公理。其次再看产生规则。假如第一个符号串满足加法标准，那么第二个一定也满足------反之亦然：假如第一个符号串不满足加法标准，则第二个符号串也不满足。这条规则使加法标准成了定理的一个遗传特性：任何一条定理都把这个特性传给由它产生的定理。这就证明了为什么加法标准是正确的。

顺便提一下，我们有一个关于pq系统的事实，这一事实使我们可以有把握地说pq系统有一个判定过程，甚至在找到加法标准之前就可以这么说。这个事实是：pq系统没有因为加长规则和缩短规则这两种相反的力量而复杂化。它只有加长规则。对于任何一个只告诉我们如何从较短的定理得到较长的定理，而永远不会反过来的形式系统，其定理都有一个判定过程。比如，假定给了一个符号串，首先，检查它是不是一条公理（我假设有一个判定公理的判定过程------否则一切就都没有希望了）。如果它是一条公理，那么它也可以称作是一条定理，这个测试就完毕。于是再假定它不是一条公理，那么，它要是一条定理，就一定是从一个较短的符号串得出来的。在逐条地试用各条规则时，不仅能准确地找出可能产生那个符号串的规则，而且还能准确地找出哪一个较短的符号串会是``家谱''中它的上一代。用这种方法，把问题``归约''成为去确定几个新的、但是较短的符号串是否都是定理，然后对它们依次地进行同样的测试。可能出现的最坏情形是产生出大量的、越来越多的、但是越来越短的需要测试的符号串。当以这种方式一步一步地回溯时，就一定会距离一切定理的源泉------公理模式越来越近。越来越短不会是无限的，因此，最终或是发现那些短符号串的其中一个是一条公理，或是到达了某一个无法再往回退的点------也就是没有一个短符号串是公理，并且没有一个短符号串可以通过某条规则或其它倒退步骤进一步缩短。这表明了在形式系统中，只有加长的规则确实是没有多大意思的。是加长与缩短的规则交互作用使形式系统具有了某种魅力。

自底向上之别于自顶向下

上述方法可以称为自顶向下的判定过程，相比之下，现在要给出的判定过程可以称为自底向上。它使人想起那个怪物用来在WJU系统中系统地生成定理的方法。但是，由于有公理模式，现在的情形要复杂些。我们将准备一只桶，当定理生成了以后，我们就把它们丢进桶里。以下便是具体做法：

（1a）将最为简单的公理（-\/-q-p-）丢进桶里。

（1b）对于桶里的每个公理应用推理规则，将推得的结果丢进桶里。

（2a）将第二简单的公理丢进桶里。

（2b）对桶里的每个符号串应用规则，并将结果都丢进桶里。

（3a）将第三简单的公理丢入桶里。

（3b）对桶里的每个符号串应用规则，并将结果都丢进桶里。

等等，等等

略微想一想之后，可以看到，你用这种方法可以把pq系统里的每一条定理都产生出来。同时，随着时间的延续，桶里会装入越来越长的定理。这又是由于没有缩短规则的结果。因此，假如你有一个特定的，像-\/-\/-\/-\/-q-p-\/-\/-这样的符号串，并且想测试一下它是不是一条定理，只要按照上述标了号码的步骤，一直不断地检查是否出现了这个符号串。如果它出现了------它就是定理。假如到了某一个时候，放进桶里的每个东西都比要测试的这个符号串长，那么，就可以丢开它------它不是一条定理。这个判定过程是自底向上的，因为它是从基础，也就是说是从公理开始往上进行的。前面的那个判定过程是自顶向下的，因为它恰恰与此相反：它是向着基础往回进行的。

同构产生意义

现在我们涉及到了本章------实际上是本书的------中心问题。也许读者已经想到，pq定理很像是加法。-\/-\/-\/-\/-q-\/-p-\/-\/-这个符号串是一条定理，因为5等于2加3。读者甚至可以认为-\/-\/-\/-\/-q-\/-p-\/-\/-这条定理是一个用奇怪的记法记下来的陈述，它的意义是5等于2加3。这样看问题合乎道理吗？嗯，我是有意地选择``q''，因为英语里``等于''这个词是equal，而选择``p''是因为英语里``加''这个词是plus。那么，-\/-\/-\/-\/-q-\/-p-\/-\/-这个符号串真的意味着``5等于2加3''吗？

是什么东西使我们那样想的呢？我的回答是：我们在pq定理与加法之间看到了同构。在导言中，``同构''这个词的定义是保存信息的变换。现在我们可以深入到那个概念中去，从另一个角度看一看。``同构''这个词的适用情景是：两个复杂结构可以互相映射，并且每一个结构的每一部分在另一个结构中都有一个相应的部分。这里``相应''的意思是：在各自的结构中，相应的两个部分起着相类似的作用。``同构''这个词的这种用法是从数学中一个更精确的概念那里得来的。

当一个数学家发现在他所知道的两个结构之间有同构关系时，就会感到很愉快。这种关系常常``从天而降''，让人惊喜不已。认识到两个已知结构有同构关系，这是知识的一个重要发展------并且我认为，正是这种对于同构的认识在人们的头脑中创造了意义。关于对同构的认识要说的最后一点是：形象地说，由于它们可以是各种各样千姿百态的，人们并非总能搞清何时才是真正地发现了同构。因此，``同构''具有通常的词所具有的模糊性。这一点是它的缺陷，但也是长处。

对于这种情况，我们有一个关于同构概念的极好范例。我们的同构有一个``较低的层次''------也就是，在两个结构中各个部分之间的映射关系：

\begin{longtable}[]{@{}lll@{}}
\toprule\noalign{}
\endhead
\bottomrule\noalign{}
\endlastfoot
q & 　\gebfont{⇔}　 & 等于 \\
p & \gebfont{⇔} & 加 \\
- & \gebfont{⇔} & 1 \\
-\/- & \gebfont{⇔} & 2 \\
-\/-\/- & \gebfont{⇔} & 3 \\
\multicolumn{3}{@{}l@{}}{%
等等} \\
\end{longtable}

这种符号与词之间的对应关系有一个名称：解释。

其次，在一个高一点的层次上，真陈述（也称``真理''）和定理之间存在着对应。但是，请注意，如果不在事先给这些符号选定一种解释，这种高一层次的对应是看不出来的。因此，更准确地说，对应是存在于真陈述和经过解释的定理之间的。不管怎么说，我们是展现了一个两层的对应，这在所有的同构中是典型的。

当遇到一个你一无所知的形式系统，并且假如你希望去发现它某种隐藏的含义时，你的问题就在于如何给它的符号赋予一种有意义的解释------也就是，通过某种方式，使得在真陈述和定理之间出现一个高层次的对应。在找到一组联系于这些符号的合适的词之前，你可能要在黑暗里进行一番摸索。这与破译密码或者释读用一种失传了的文字写成的铭文（比如克里特岛的线形文字B）非常相似。释读的唯一方法就是建立在以知识为基础的猜想上的试错法。当你发现一个好的选择，一个``有意义''的选择时，突然间就觉得顺当了，并且工作的速度大大加快了。不久，件件事情就都各就各位了。约翰·查德威克的《克里特岛线形文字的释读》一书中就包含有这种令人激动的经历。

但是，对挖掘一个已经毁灭了的文明时发现的形式系统进行释读，这种机会实在太少了。数学家（近来还有语言学家、哲学家和其它人）是仅有的使用形式系统的人。而且他们心里对于他们所使用和出版的形式系统总是有个解释的。这些人的想法是：建立一个形式系统，它的定理同构地反映一部分现实世界。在这种情况下，对于符号的选择是有高度目的性的，就像是对于符号串的生成规则的选择那样。我设计pq系统时就是这样的。你可以看出来，我为什么选用了我所选的符号。定理与加法同构这一点不是偶然的，同构之所以发生，是因为我有意地找到了一种通过符号形式反映加法的方法。

有意义的解释与无意义的解释

你可以选择与我不同的解释。你无需使每条定理都成为真的。但要是建立一个解释，使每条定理都成为假的，这也没什么意思。当然，要是作一个解释，其中的定理和真理之间没有一点正的或反的相关性，这就更没有意思了。因此，让我们区别开对于形式系统的两类解释。首先，我们可以有一些无意义的解释，在这种解释下，看不到系统中的定理与现实世界之间有什么同构关系。这样的解释多得很------任何随意的选择，都是这样的解释。例如下面这个：

\begin{longtable}[]{@{}ll@{}}
\toprule\noalign{}
\endhead
\bottomrule\noalign{}
\endlastfoot
q　 & 马 \\
p　 & 幸福 \\
-　 & 苹果 \\
\end{longtable}

那么-q-p-\/-就获得了一个新的解释：``苹果马苹果幸福苹果苹果''------很有诗意的，它可能会打动马，甚至可能使它们喜欢这种对于pq符号串的解释！这种解释没什么``意义''。在这种解释中，定理听起来既不比非定理真一些，也不比它们好些。一匹马可以欣赏``幸福幸福幸福苹果马''（映射成ppp-q），就像它欣赏任何一条经过解释的定理一样。

另一种解释叫做有意义的。在这种解释下，定理和真理是相对应的------也就是，同构存在于定理与现实的某一部分中。这就是区别解释与意义的好处。任何已有的词都可以用来解释``p''，但是``加''是我们所碰到的唯一有意义的一个。总的来说，``p''的意思看来应该是``加''，虽然它还可以有成千上万个不同的解释。

主动意义之别于被动意义

如果深一步地理解，这一章的最重要的事实也许是：pq系统似乎在迫使我们认识到，形式系统中的符号虽然一开始没有意义，但至少在发现了同构关系时，不可避免地会带上某种``意义''。然而，形式系统中的意义与语言中的意义的区别是非常重要的。这区别是：在语言中，当我们知道了一个记号的意义，我们就能基于这个记号做出新的陈述。从某种意义上说，语言中的意义是``主动的''，因为它为创造句子带来了一条新规则。就是说，我们对于语言的掌握，不像是对于一件完成了的产品那样。当我们知道了新的意义的时候，造句子的规则也增加了。另一方面，在形式系统中，定理都是运用产生式规则事先定义了的，我们可以选择以定理与真陈述之间的同构（假如可以找到的话）为基础的``意义''。但是，不是说这样就准许我们在已建立的定理之外增加新的定理。这就是第一章中形式化要求所警告过的。

在WJU系统中，当然没有什么东西诱惑我们去超出那四条规则，因为还没有找到或发现一种解释。但是在这里，在我们的新系统中，人们可能受到新发现的每一个符号的``意义''的影响而认为符号串

-\/-\/-\/-\/-\/-\/-\/-q-\/-p-\/-p-\/-p-\/-

是一条定理。至少，人们可能希望这个符号串是一条定理。但是愿望并不能改变它并不是一条定理这一事实。仅仅因为8等于2加2加2加2，就认为它``一定''是一条定理，这是极大的错误。甚至赋予它意义就是错误的，因为它不是良构的，而我们的有意义的解释整个都是从良构符号串中得出的。

在一个形式系统中，意义一定是``被动的''。我们可以根据其组成符号的意义来读每个符号串，但是我们没有权力只根据我们给符号指定的意义而创造出新的定理。经过解释的形式系统是横跨在没有意义的系统与有意义的系统之间的。可以认为，它们的符号串是``表达''事物的，但这也必定是系统的形式性质的产物。

双重意义！

现在我要破除一种错觉，即认为已经找到了pq系统中符号的唯一意义。考虑下列联想：

\begin{longtable}[]{@{}lll@{}}
\toprule\noalign{}
\endhead
\bottomrule\noalign{}
\endlastfoot
q & 　\gebfont{⇔}　 & 减 \\
p & \gebfont{⇔} & 等于 \\
- & \gebfont{⇔} & 1 \\
-\/- & \gebfont{⇔} & 2 \\
\multicolumn{3}{@{}l@{}}{%
等等} \\
\end{longtable}

那么-\/-\/-\/-\/-q-\/-p-\/-\/-就有了一个新的解释：``5减2等于3''。当然这是一个真陈述。在这种新的解释之下，所有的定理都将是真理。这就像原来那个解释一样有意义。很显然，``哪一个是这个符号串的唯一意义？''就成了一个很愚蠢的问题。解释只要精确地反映现实世界的某种同构，就是有意义的。当现实世界中的不同方面彼此同构时（这里是说加法和减法），一个形式系统可以与这两者都同构，因此可以有两种被动意义。这种符号和符号串的二值性是一个极其重要的现象。在这里，它似乎是不足道的、奇怪的、令人恼火的。但是，在以后更深入的讨论中，它还会再次出现，并带来极丰富的内容。

现在来把我们对于pq系统所做的观察作一个总结。给定两个有意义的解释，那么用其中的任何一个，每个良构的符号串都对应于一个合法的断言------有些是真的，有些是假的。在一个形式系统中，良构符号串的含义就是：当对它一个符号一个符号地去解释时，就产生合语法的句子（当然，这取决于如何解释。但是，通常在人们脑子里有一种解释）。定理出现在良构符号串中。这些定理是由一个公理模式和一条生成规则来定义的。我发明pq系统的目的是要模仿加法：我想让每一条定理在解释之下都表达一个真正的加法；反之，我想让每一个真正的恰好是两个整数相加的加法都能翻译成一个是定理的符号串。我的目的达到了。请注意，这样，所有的错的加法，像``6等于2加3''，将映射成一些不是定理的良构符号串。

形式系统和现实

在那些基于现实的一部分的形式系统中，pq系统是我们遇到的第一个实例，并且看起来模仿得很好。在这种模仿中，它的定理和有关那部分现实中的真理同构。然而，现实和形式系统是互相独立的，并不需要意识到在两者之间存在着同构关系。每一方面都独立存在------不论我们知道不知道2等于1加1，-\/-q-p-都是一条定理；并且不论我们是否把它与加法相联系，-\/-q-p-总是一条定理。

读者可能想知道，构造这样或者那样的形式系统，是否能就其论域（它的解释所在的领域）中的真理给我们一些启发。通过构造pq定理，我们知道了什么新的加法吗？当然没有，但是我们知道了关于加法作为一种操作的性质------也就是说，它很容易用支配无意义符号的印刷符号（以下简称印符）规则加以模仿。这仍然不应使人惊异，因为加法是那么一个简单的概念。加法可以很简单地体现在一部机器的旋转的齿轮之中，就像在现金出纳机中那样。

但是很清楚，就形式系统而言，我们几乎只触及了表面。人们很自然地想知道，现实的哪一部分能够用一组支配无意义符号的形式规则来加以模仿。现实世界的一切都可以变为形式系统吗？从一个很广的意义上说，回答可能是肯定的。比如，人们可以设想，现实世界本身只不过是一个非常复杂的形式系统。它的符号不是在纸上移动，而是在一个三维空间里运动，它们是组成一切事物的基本粒子（这里包含着一个隐含的假定，即：物质的可分性是有尽头的，因此，基本粒子这个词是有意义的）。``印符规则''是物理法则，它告诉人们如何根据给定的时刻从给出的所有粒子的位置和速度，得出属于``下一个''瞬间的一组新的位置和速度。所以这个宏伟的形式系统的定理，就是粒子在宇宙历史中不同时间的可能布局。唯一的公理是所有粒子在``时间的开始''时的原始布局。这是一个非常宏伟的设想，然而它只有纯粹的理论意义。另外，量子力学（及物理学的其它部分）也使人们对于上述概念所含的理论价值产生了某种怀疑。从根本上说，我们是在问宇宙是否是以完全确定的方式运动，而这是个尚未解决的问题。

数学与符号处理

我们将不去讨论那么大的一幅图景，而只以数学做为我们的``现实世界''。这里出现了一个严重的问题：假如我们试着按照某数学分支构造一个形式系统，我们如何确定我们是否已把这项工作做得很准确了呢？尤其是假如我们不是已经百分之百地熟悉了那部分数学的话。假定形式系统的目的是要给我们带来那个学科的新知识。那么，在我们证明了同构是完全的之前，怎么知道对于每一条定理的解释都是真理？并且，如果我们在开始时不知道那个学科中的所有真理，我们又怎样证明同构是完全的呢？

假定在某个遗址的发掘中，我们确实发现了某种神秘的形式系统。我们可以试验各种不同的解释方法，并且也许最终碰上了一种看起来使每一个定理成为真理，每一个非定理成为``假理''的解释。但这是一种只能在有穷多个情形里进行直接检验的东西，而定理的集合多半可能是无穷的。我们怎么会知道所有的定理在这种解释方法之下表达了真理？除非我们知道关于形式系统及相应论域的每一个细节。

当我们企图将自然数世界（也就是非负整数：0，1，2，3\ldots\ldots）和一个形式系统的印刷符号相匹配时，我们就会发现自己是处于这样一种奇怪的境地了。我们将努力去理解，我们在数论中称作``真理''的，和我们借符号处理所能做到的，这两者之间是什么关系。

因此，让我们简略地看看，当我们称数论中的某些陈述为真，而另一些为假时，其根据是什么。12乘以12得多少？人人皆知，得144。但是有多少给出这个答案的人在他们的生活中曾经画过一个12×12的长方形，然后数它的小方格？多数人认为画格和数它们是不必要的。他们会代之以在纸上写下如下的式子作为证明：

\begin{center}

\pandocbounded{\includegraphics[keepaspectratio]{./media/OEBPS/Images/Formula01.png}}

\end{center}

这就是``证明''了。几乎所有的人都相信，如果去数小方格，一定是144个，几乎没有人会对此得数感到怀疑。

当考虑确定987654321×123456789的值时，这两种观点的差别就会看得更清楚了。首先，根本不可能构造一个合适的长方形；更糟糕的是，即便造出来了，也得要大批的人用几个世纪的时间，不停地数那些格子。而只有极易轻信的人，才愿意相信他们最后数出的结果。某个人、在某个地方、以某种方式出个小小的失误，这简直是太可能发生的事了。因此，究竟是否有可能知道它的得数？如果信任按照某种简单规则处理数字的符号操作过程，回答是肯定的。这个过程作为一种得到正确答案的方法提供给孩子们，而对于许多孩子来说，这个过程的道理由于在这么多的符号中摆弄来摆弄去，往往就显得不清楚了。乘法的数字变换法则多半是基于加法、乘法的少数几条性质，而这些性质被假定是对于所有的数都有效的。

算术的基本法则

对我所说的那种假定，现在就可以说明如下。设划了几个道道：

/ // // // / /

现在数一数它们。同时另外一个人从另一头数。两个人会得到同一个答案，这一点是否清楚？数的结果与计数的方法是无关的。这就是对于什么是计数的一种假定。要试图证明这一点将是毫无意义的，因为它太基本了，不管是看得出来还是看不出来。如果看不出来，就是证明也不会有一点儿帮助。

从这样的假定中，可以得出加法的交换律和结合律（也就是说：第一：永远有b+c=c+b，第二：永远有b+(c+d)=(b+c)+d）。同样的假定也可以引出乘法交换律和结合律，这只要想想许多小方块汇集在一起形成一个大的长方体的情形。乘法的交换律和结合律就是这样一种假定：以各种方式旋转那个长方体时，小立方体的数量不会改变。这些假定不是在所有可能的情况下都可以验证的，因为，这类的情况是无穷的。我们认为这是当然的，我们像相信其它事物一样深信它们（假如我们想到它们）。我们从街上走过时，口袋里的钱上上下下地相撞，但是数量不会改变。如果我们把我们的书装入盒子中，放进汽车，开一百里，把盒子拿下汽车，打开，把书放上一个新的书架，书的数量也不会改变。所有这些都是数的性质的一部分。

有这样的人，当某种不可否认的事实一旦写下来时，总是喜欢找点理由，使那个``事实''干脆成了假的。我就是这样的人。在我刚刚写下上面的关于线段、钱和书等例子的时候，我就又设想了一些情境，在其中它们可能并不成立。读者也可以这样做。它说明抽象的数与日常生活中我们使用的数确实是很不相同的。

人们喜欢造些违反基本算术的说法，以便说明某种``更深刻的''真理。比如``1+1=1''（对恋人而言），或者``1+1+1=1''（三位一体）。很容易从上述说法里找出漏洞，比如，指出在上述两种情况中，为什么使用加号是不合适的。但这样的情况多得很。两滴雨水从窗户玻璃上流下来合成了一滴，1加1是不是等于1？一朵云彩分成两朵------还要更多类似的证据吗？发生的事情可以称作``加''，或者需要别的什么词，在这两者之间划一条明显的界限可不是轻而易举的。如果你考虑这个问题，你大概就会提出关于空间中物体的分隔之类的标准，并且肯定每一种物体与其它物体有很明显的区别。但是如何去数思想呢？如何去数组成大气的气体的数量？假如查一查的话，大概能在某个地方发现一个像``印度有十七种语言和四百六十二种方言''这样的陈述。这样精确的陈述倒多少令人感到奇怪了，因为``语言''和``方言''这两个概念本身就是模糊的。

理想的数

数作为一种``实在''是不守规矩的。然而，人们有一种古老的内在感觉，即数不应该不守规矩。在数的抽象概念中，就有那么一种纯净的东西，独立于数念珠、方言或云朵。应该有一种谈论数的方式，这种方式不会总是受到现实所具有的那种糟糕状况的侵入和捣乱。支配``理想''数的铁定的规则构成了算术，并且进一步构成了数论。在把作为实际事物的数转变成作为形式事物的数的时候，只有一个有关的问题需要提出来。一旦你决定了要把所有数论装入一个理想的系统，是否真有可能彻底地做到这一点呢？数是否那么整洁，那么有规律，像结晶一样，以至于形式系统的规则能彻底地把握它们的性质？艾舍尔最漂亮的画之一《释放》（图13），就是形式与非形式之间的绝妙的对比，它还有一个迷人的过渡区域。数真的像鸟那么自由吗？它们被囚禁在遵守规则的系统中时，也像那些鸟一样难以忍受吗？现实中的数与写在纸上的数之间是否有一个不可思议的过渡区域？

\pandocbounded{\includegraphics[keepaspectratio]{./media/OEBPS/Images/img_013.jpg}}

图13．释放

艾舍尔作（石版画，1955)

说到自然数的性质时，我不是单指像``特定的两个整数的和''这样的性质。这种性质可以通过计数找出来。每一个在二十世纪里长大的人都不会怀疑像这样的计数以及加法、乘法等等这些过程的可机械化性。我所指的性质是数学家有兴致去探讨的那些问题，而计数过程不足以为那些问题提供答案------甚至在理论上，也不足以提供答案。让我们举一个古老的关于自然数性质的例子：``有无穷多个素数。''首先，没有一个计数过程能够证实或推翻这一断言。我们可以做到的，最多只不过是数上一会儿有多少个素数，然后承认，确实``有不少''。但是，不管你数出了多少素数，你的结果都无法确定素数的个数是有穷的还是无穷的。总是数不完。这个陈述------被称做``\heititext{欧几里德定理}''（请注意是黑体字）------不是显然的。看上去它可能是合理的和令人信服的，但不是显然的。然而，欧几里德及后来的数学家一直说它是真的。什么原因呢？

欧几里德的证明

原因是，推理告诉他们是这样的。让我们来看看这种推理。我们来看一个欧几里德证明的变种。这个证明告诉人们，不论挑出哪一个数，都还有一个比它大的素数。挑出一个数------N，把所有从1到N的正整数相乘，换句话说，也就是做N的阶乘，写作``N!''。得到的数可以被从1到N之间的任何数除尽。然后在N!上加上1，结果这个数就

不能是2的倍数（因为当被2除时，余1）；

不能是3的倍数（因为当被3除时，余1）；

不能是4的倍数（因为当被4除时，余1）；

\ldots\ldots{}

不能是N的倍数（因为当被N除时，余1）；

换句话说，假如N!+1是可以被除尽的（除了被1和它本身除之外），也只能是被比N大的数除尽。因此，或者它本身是素数，或者它的素因数比N大。但是不论是哪种情况，我们已经表明，都一定有比N大的素数存在。不论N是什么数，这个论证过程都能进行。因此，不管N是什么数，都有一个素数大于N。至此，就完成了素数有无穷多个这一证明。

顺便说一下，这最后一步叫做推广。在今后更形式化的上下文中，我们还会遇到它。在那里我们就单独一个数(N)来进行一个论证，然后指出我们用的N是非特定的，因此我们的论证是一般性的。

欧几里德的证明是典型的所谓``真正的数学''。它简洁、优美且令人信服。它表明人们可以通过一些很简短的步骤从起点走出很远。在我们这个例子里，起点是关于乘法和除法等等的基本概念。短小的步骤是指推理的步骤。虽然每一个单个的推理步骤看起来是显然的，但最后的结果并不显然。我们永远不能直接地检查这个陈述是否是真的，然而我们相信它，因为我们相信推理。如果你接受推理，似乎就别无选择。一旦你同意听欧几里德把话讲完，你就得同意他的结论。这是件幸事------因为这意味着数学家对于把什么样的陈述标为``真的''，把什么样的标为``假的''，永远持一致意见。

这个证明是有条理的思维过程的一个例子。每一个陈述都以一种不可抗拒的方式与前一个陈述有着联系。这就是为什么称它为``证明''而不只是``好的证据''。在数学中，目的总是给予某个不那么显然的陈述一个滴水不漏的证明。存在有这种以滴水不漏的方式联系在一起的步骤，这一事实本身就暗示着可能存在一个具有模式的结构把这些陈述联在一起了。展示这个结构的最好办法是找一些新的词汇------有特定的风格、由符号构成的词汇------仅适于表达关于数的陈述的新词汇。然后，我们可以看看经这些新词汇翻译后的证明。它将是一行接一行地以某种可以检查的方式联系在一起的一些陈述。但是，因为这些陈述是由具有特定风格的小符号集合来表示的，所以它们具有模式的样子。换句话说，虽然大声朗读出来时它们好像是关于数以及它们的性质的陈述，但当把它们写在纸上时，它们仍然好像是抽象的模式------并且这个证明的一行接一行的结构可能开始变得像是模式按照某些印符规则在缓慢地变换。

绕过无穷

虽然欧几里德的证明是一个关于所有的数都有某种性质的证明，它却回避了分别讨论无穷多种情况中的每一个。它用像``不管是什么''或者``不论N是一个什么数''这样的话来绕过它。我们也可以使用``所有的N''这样的词再陈述一遍这个证明。知道了合适的上下文及使用这些词语的正确方式，我们就永远不需要处理无穷多个陈述。我们只要处理两三个概念，比如``所有''这个词------虽然它本身是有穷的，但它体现了无穷。通过使用它们，我们避开了要证明无穷多个事实这种表面问题。

我们使用``所有''这个词的可能方式是由推理的思维过程来规定的。也就是说，在我们使用``所有''这个词时，要遵循一些规则。我们也许没有意识到它们，并且倾向于认为，我们是在这个词的意义的基础上操作的。但是，这毕竟只是一种迂回的说法，实际上我们是由一些从未摆出来的规则所指引的。我们一生中总是在某些句型中使用某些词，但我们不把这些句型叫做``规则''，而是把我们的思维运转过程归因于词汇的``意义''。这个发现在通向数论形式化的漫长道路上是一个关键性的认识。

假如我们越来越仔细地研究欧几里德的证明，我们就会看到，它是由许许多多小的------几乎是无穷小的------步骤组成的。假如把所有这些步骤都一行一行地都写出来，这个证明会难以置信地复杂。对于我们的头脑来说，把许多步骤压缩在一起形成单个句子是最为清晰的。假如我们用慢动作来观察这个证明，我们就会开始察觉到一个个的框架。换句话说，这个解剖过程到此为止，这之后我们就遇到推理过程的``原子化''性质。一个证明可以分解成一系列小而不连续的跳跃。但从一个更高更便利的角度来观察它的时候，这些步骤是很平滑地过渡的。在第八章中，我将说明一种把证明分成原子单位的方法，读者会看到它涉及了许多步骤，多得令人难以置信。即使这样，它或许也不应该让你太惊异。当欧几里德发明这个证明时，一定有几百万个神经元（神经细胞）参加了操作，其中有许多仅在一秒钟内就被激活了几百次。仅仅说出一个句子就要涉及成千上万个神经元。假如欧几里德的想法那么复杂的话，说它的证明包含了巨大数目的步骤就有意义了（他脑中神经元的活动与我们的形式系统里的一个证明可能没有什么直接联系。但是两者的复杂程度是可以相比的。这就好像是大自然想要使``素数有无穷多''这个证明的复杂性保持恒定，即使涉及到的系统相互之间极其不同）。

在下面的几章中，我们将构造出一个形式系统：（1）它包含一个具有特殊风格的词汇表，所有的关于自然数的陈述都可以通过它来表达；（2）它有与所有看起来必要的推理类型相对应的规则。一个很重要的问题是：我们那时所给出的、用于符号处理的规则，是否真的与我们通常的心理推理能力相等（就数论而言）------或者更一般地说，用某种形式系统来达到我们思维能力的水平，在理论上是否有可能。

\phantomsection\label{Dialog03.xhtml}{}

\chapter{无伴奏阿基里斯奏鸣曲}

电话响了，阿基里斯拿起话筒。

\dialogname{阿基里斯：}喂，我是阿基里斯。

\dialogname{阿基里斯：}噢，你好，龟兄，你怎么样？

\dialogname{阿基里斯：}得了窝颈症？啊，真不幸。你自己知道是什么引起的吗？

\dialogname{阿基里斯：}你呆在那种姿态下有多久？

\dialogname{阿基里斯：}哦，难怪变得僵硬了。究竟是什么东西居然叫你那么长时间地拧着脖子？

\dialogname{阿基里斯：}其中有很多都挺怪模怪样的，是吗？都是些什么东西啊？举个例子？

\dialogname{阿基里斯：}你说什么------``魔幻般的各种动物''？

\pandocbounded{\includegraphics[keepaspectratio]{./media/OEBPS/Images/img_014.jpg}}

图14．镶嵌画Ⅱ

艾舍尔作（蚀版画，1957）

\dialogname{阿基里斯：}你认出了你的朋友螃蟹？它干嘛要变成那种样子？什么？螃蟹总是爱把自己弄进最不可思议的困境里？这呆瓜！幸亏那对钳子没变样。

\dialogname{阿基里斯：}还有一把吉他？太奇怪啦。说到音乐，你想过来听一听那位最对你口味的作曲家巴赫的无伴奏小提琴奏鸣曲吗？我刚买到一张这些奏鸣曲的唱片，那效果真是太棒啦。我现在还没从那里面醒过劲儿来呢。巴赫只用了一把小提琴就创作出了这么有趣的作品！

\dialogname{阿基里斯：}怎么，还没劲儿？我知道到了你这样的岁数会怎么样。真遗憾，也许你应该上床读会儿书什么的。

\dialogname{阿基里斯：}哦，是吗，没电了？有时会出现这种情况的，不过我看你可以点上蜡烛。

\dialogname{阿基里斯：}哦，哦，原来是因为失眠，那最好就别看书了。可是，究竟是什么叫你睡不着觉呢？

\dialogname{阿基里斯：}哦，是这么回事啊。那好，如果是这个叫你一直神不守舍，那你也许最好说给我听听，让我和你一起来研究研究。

\dialogname{阿基里斯：}一个词，其居中的两个部首依次是``昔''和``火''\ldots\ldots 嗯\ldots\ldots``秋鹊''怎么样？

\dialogname{阿基里斯：}的确，``昔''、``火''在这个词里顺序是颠倒的。

\dialogname{阿基里斯：}一连猜了几个小时？听起来我好像一头栽进了一个非常难解的字谜里了。这个该死的字谜你是打哪儿听来的？

\dialogname{阿基里斯：}你是说他看去像是在思索深奥的佛理，可实际上他只是在力图解开复杂的字谜，是吗？

\dialogname{阿基里斯：}啊哈！------原来是那只蜗牛知道这家伙在干什么。可是你怎么想起跟那只蜗牛说话的呢？

\dialogname{阿基里斯：}听我说，我曾经见过一个字谜，跟这个有点像。你想知道吗？这会不会弄得你更加神不守舍呢？

\dialogname{阿基里斯：}我也这么想------不会有任何害处的。下面就是那个字谜：什么词以部首``虫''开头，又以部首``虫''结尾？

\dialogname{阿基里斯：}真机灵------不过这近乎赖皮了。这当然不是我心里所想的那个答案。

\dialogname{阿基里斯：}你当然很对------它满足了条件，不过这是种``走捷径''的解决办法。我心里另有谜底。

\dialogname{阿基里斯：}叫你猜着了！你怎么这么快就想出来啦？

\dialogname{阿基里斯：}我明白了。事实上停电不仅没有耽误你的事，反而帮了你的忙。太妙了！可是对于你的那个``昔火''谜，我这儿还是一抹黑呢。

\dialogname{阿基里斯：}祝贺你！这样一来你也许能睡着了！告诉我吧，谜底是什么？

\dialogname{阿基里斯：}是的，我一般不喜欢提示，但是没关系，你要给我什么提示呢？

\dialogname{阿基里斯：}我不懂你在这里用``图形''和``衬底''是表示什么。

\dialogname{阿基里斯：}我当然知道《镶嵌画Ⅱ》那幅画！我熟悉艾舍尔的所有作品。毕竟，他是我最喜欢的艺术家。我无论到哪儿总要在墙上挂一幅《镶嵌画Ⅱ》的印刷复制品，我现在一抬眼就能看得见。

\dialogname{阿基里斯：}是的，我看出了所有的黑颜色动物。

\dialogname{阿基里斯：}对，我还看到它们中间的``负空间''------也就是剩下的画面------是些白颜色的动物图案。

\dialogname{阿基里斯：}你所说的``图形''和``衬底''原来是这个意思啊。可是这同那个``昔火''字谜有什么关系呢？

\dialogname{阿基里斯：}哦，对我来说这太复杂了。我真希望能马上来个``肚里点蜡烛------透心明白''。

\dialogname{阿基里斯：}没听说过？这可是我最喜欢的歇后语了。

\dialogname{阿基里斯：}你现在就准备过来？可是我刚才还想------

\dialogname{阿基里斯：}那好吧。也许到那时我会想出你那个字谜的正确答案，对，用你那个``图形-衬底''的暗示，把它同我的那个字谜联系起来。

\dialogname{阿基里斯：}我很愿意把它们放给你听。

\dialogname{阿基里斯：}你发明了一种有关它们的理论？

\dialogname{阿基里斯：}用什么乐器伴奏？

\dialogname{阿基里斯：}嗯，如果真是这样的话，他不愿写出羽管键琴部分并把它发表就有点奇怪了。

\dialogname{阿基里斯：}我懂了------带点儿随意的色彩。听它们可以有两种办法------带伴奏的或不带伴奏的。但是人们怎么会知道那伴奏听起来怎么样呢？

\dialogname{阿基里斯：}啊，对，我想那是最好的，把它留给听众去想象。况且还说不定正像你所说的那样，巴赫甚至从未想到过要有伴奏。这些奏鸣曲现在这样似乎也确实很好。

\dialogname{阿基里斯：}没错儿。好吧，一会儿见。

\dialogname{阿基里斯：}再见，龟兄。

\phantomsection\label{Chapter03.xhtml}{}

第三章

\chapter{图形与衬底}

素数之别于合数

想想简单的印符操作能把握概念，不免有点令人奇怪。迄今为止我们已把握到的概念是加法，这说来还不显得有多怪。可我们现在的目标是做出一个形式系统，使定理都是Px的形状，其中字母x代表一个短杠符号串，并且只有当短杠符号串中的短杠数目是素数时才能成为定理。这样，P-\/-\/-是个定理，而P-\/-\/-\/-则不是。怎么能用印符操作来做这件事？首先，说清楚所谓印符操作是什么意思是很重要的。我们能用的东西都已在WJU系统和pq系统中给出了，所以我们实际上只需列出我们所允许的种种操作：

⑴．读入并认出有限的符号集合中任何一个符号；

⑵．写下属于该集合的任何一个符号；

⑶．把任何一些那样的符号从一处复制到另一处；

⑷．删除任何一些那样的符号；

⑸．检测一个符号是否与另一个符号相同；

⑹．保存并使用一列以前产生的定理。

这张表有点多余，不过没关系。重要的是它很明显只涉及一些极简单的能力，它们远比区分素数与非素数的能力简单得多。那么，我们怎样才能组合这些操作，构造一个形式系统，在其中素数能与合数区别开呢？

tq系统

第一步该是试图去解决一个简单些、但却与此有关的问题。我们可以做一个与pq系统类似的系统，差别是这次我们表示的是乘法。我们称之为tq系统，``t''代表``乘''（英文是times）。进一步还得说明，我们用X、Y、Z分别表示短杠符号串x、y、z中的短杠数目（注意我在此特别用心地区分符号串与它所包含的短杠数目）。现在我们希望符号串xqytz是个定理，当且仅当X等于Y乘以Z。例如，-\/-\/-\/-\/-\/-q-\/-t-\/-\/-应该是个定理，因为2乘以3等于6，但-\/-\/-q-t-\/-不应是个定理。tq系统可以像pq系统一样很容易地刻划------就是说，只用一个公理模式和一条推理规则：

公理模式：xqxt-是一个公理，对任何一个短杠符号串x都是如此。

推理规则：设x、y、z都是短杠符号串。设xqytz是个已有的定理。那么，xyqytz-是个新的定理。

下面是定理-\/-\/-\/-\/-\/-q-\/-t-\/-\/-的推导：

\begin{longtable}[]{@{}lll@{}}
\toprule\noalign{}
\endhead
\bottomrule\noalign{}
\endlastfoot
⑴ & -\/-q-\/-t- & （公理） \\
⑵ & -\/-\/-\/-q-\/-t-\/- & （用推理规则，⑴作为已有定理） \\
⑶ & -\/-\/-\/-\/-\/-q-\/-t-\/-\/- & （用推理规则，⑵作为已有定理） \\
\end{longtable}

若注意每次使用推理规则时最后的那个短杠符号串是怎样增长的，你就可以预言，如果想导出一个最后有十个短杠的定理，就得接连用九次推理规则。

把握合数

乘法是个比加法稍难对付的概念，现在已由印符操作``把握住''了，就像艾舍尔的《释放》里的那些鸟。素数性质怎么样呢？这里有个看来聪明的计划：使用tq系统，定义一集形如Cx的新定理来刻划合数：

规则：设x、y、z是短杠符号串。如果xqy-tz-是定理，则Cx是定理。

这能行得通，因为：说X（x中的短杠数目）是合数，仅当它是两个大于1的数的乘积------即，是Y+1（y-中的短杠数目）和Z+1（z-中的短杠数目）的乘积。我可以提出某种``惟方式''的论证来为这条规则辩护，因为你是人，想知道为什么会有这样一条规则。假如你是完全以``机方式''运转的，你就不会需要任何论证了，因为J方式工作者只是机械地、快乐地遵循规则，从不产生疑问。

由于你以W方式工作，你会趋于把符号串与对符号串的解释相混淆。你知道，一旦你从你所操作的符号里找到``意义''，事情就会变得相当混乱。你不得不和你的自我作斗争，阻止自己把符号串``-\/-\/-''认作数目3。对形式化的要求，在第一章里或许显得令人迷惑不解（因为它看上去如此显然），在这里却变得相当复杂，并且很关键。这里的实质在于防止你自己把W方式与J方式相混淆，或者换一种说法，在于防止你自己把算术事实与印符定理相混淆。

对素数的非法刻划

一个很诱人的做法是从C型定理直接跳到P型定理，就是草拟下面这样一个规则：

草拟的规则：设x是一个短杠符号串。如果Cx不是个定理，那么Px是个定理。

这里致命的缺陷是：Cx是否是个定理并非一个显明地表达出的印符操作。要确定WU不是WJU系统的定理，你必须到系统的外面去\ldots\ldots 对这个草拟的规则也是如此。这是一个完全违反了形式系统观念的规则，它要求你非形式化地进行操作------就是说，在系统之外进行操作。印符操作（6）允许你查看以前所发现的定理的储备，可是这条草拟的规则却要求你去查看假定的一张``非定理表''。但要想产生这样一张表，你必须得在系统之外进行一番推理------这种推理是说明为什么某些符号串不能在系统内生成。当然很可能有那么一个另外的形式系统，能纯粹通过印符手段产生``非定理表''。事实上，我们的目标恰恰就是要找这么一个形式系统。但上面草拟的那个规则不是个印符规则，必须放弃。

这一点非常重要，因此我们还得再稍微多说几句。在我们的C系统里（包括tq系统及定义C型定理的那条规则），我们有形如Cx的定理，其中``x''像通常一样代表一个短杠符号串。也有形如Cx的非定理（这些就是当我说``非定理''时所指的东西，虽然tt-Cqq以及其它乱七八糟不规矩的东西也显然是非定理）。区别在于，定理含有合数数目的短杠，非定理含有素数数目的短杠。现在所有的定理都具有一个共同的``形式''，也就是说，都源自一个共同的印符规则集。那么，在同一意义上，所有的非定理是否也具有一个共同``形式''呢？下面是一列C型定理，不带推导。它们后面括号中的数字只是记录其中的短杠数目。

\begin{longtable}[]{@{}l@{}}
\toprule\noalign{}
\endhead
\bottomrule\noalign{}
\endlastfoot
C-\/-\/-\/-(4) \\
C-\/-\/-\/-\/-\/-(6) \\
C-\/-\/-\/-\/-\/-\/-\/-(8) \\
C-\/-\/-\/-\/-\/-\/-\/-\/-(9) \\
C-\/-\/-\/-\/-\/-\/-\/-\/-\/-(10) \\
C-\/-\/-\/-\/-\/-\/-\/-\/-\/-\/-\/-(12) \\
C-\/-\/-\/-\/-\/-\/-\/-\/-\/-\/-\/-\/-\/-(14) \\
C-\/-\/-\/-\/-\/-\/-\/-\/-\/-\/-\/-\/-\/-\/-(15) \\
C-\/-\/-\/-\/-\/-\/-\/-\/-\/-\/-\/-\/-\/-\/-\/-(16) \\
C-\/-\/-\/-\/-\/-\/-\/-\/-\/-\/-\/-\/-\/-\/-\/-\/-\/-(18) \\
\gebfont{︙} \\
\end{longtable}

这张表中的``洞''是那些非定理。重复那个早先的问题：那些洞也具有某种共同的``形式''吗？对，也不对。无可否认的是它们都具有某种印符特性，但我们是否要称之为``形式''却是不清楚的。理由是那些洞仅仅是以否定的方式定义的------它们是被一列以肯定的方式定义的东西所排除的那些东西。

图形和衬底

这让人想起艺术中对``图形''和``衬底''的著名区分。当一个图形或``正空间''（例如，一个人形、一个字母、一个静物）画在画框里时，不可避免地也就画上了与它互补的形状------也称作``衬底''、``背景''或``负空间''。在多数绘画中，这种图形与衬底的关系不起多少作用。艺术家对衬底远不如对图形那么感兴趣。但有时候，艺术家也会对衬底同样感兴趣。

有一种漂亮的汉字书写方式就是着意摆弄这种图形与衬底的区分。下面的图15就是用这种笔画来传达信息的。猛一看，它像是一团团随便涂上的墨斑，但若你退后一些盯着它看一会，你会突然发现，有四个汉字出现在这里\ldots\ldots{}

\pandocbounded{\includegraphics[keepaspectratio]{./media/OEBPS/Images/img_015.png}}

图15．``以子之矛''

刘皓明绘

类似的效果可以在我画的《以烟为号》（图139）中看到。沿着这条思路下去，你可以考虑一下下面这个谜题：你能否想个什么办法作一幅画，使其中的图形及衬底都是图形？

\pandocbounded{\includegraphics[keepaspectratio]{./media/OEBPS/Images/img_016.jpg}}

图16．以鸟作瓦

艾舍尔作（选自1942年的速写本）

现在我们来正式地区分两类图形：流畅可画出的和倍流畅的（顺便说明，这些词是我自己的用语------并非通常都这么用）。在流畅可画出的图形中，衬底仅仅是绘画过程中顺带的副产品。而在倍流畅的图形中，衬底本身也可视为一个图形。这通常都需要艺术家的精心构思。``倍流畅''的``倍''字表示前景和背景都是流畅地画出来的------图形是``双倍流畅的''。在倍流畅图形中，图形与衬底的每道分界线都是双刃宝剑。艾舍尔是画倍流畅图形的大师------例如，他倍流畅地画出的那些漂亮的鸟（图16）。

我们这种区分并非像数学里那么精确。谁能肯定地说一幅画中的衬底就一定不是个图形？一旦指出这一点，那么几乎任何衬底都有它自己的价值。在这种意义下，每个图形都是倍流畅的。但那不是我这用语的意图。我们对可识别的形状都有种自然直觉的观念。前景与背景是否都是可识别的形状？如果是，那么这幅画是倍流畅的。当你看大多数线条画的衬底时，你会发现它们是极难识别的。这就证明了：

{存在可识别的形状，其负空间是不可识别的形状。}

用更为``技术性''的词语来说，就是：

{存在非倍流畅的流畅可画出图形。}

\pandocbounded{\includegraphics[keepaspectratio]{./media/OEBPS/Images/img_017.png}}

图17．FIGURE-FIGURE图形

斯科特·凯姆作（1975）。你能看出其中的单词吗？

对于前面说到的那个谜题，斯科特·凯姆给出了一个解：他直接在英语单词FIGURE(``图形''的意思)上作文章。图17就是他作的画，我称之为FIGURE-FIGURE图形。如果你看了这幅画中白的部分再看黑的部分，你看到的都会是FIGURE，就是说，你总是看到``图形''，而永远看不到``衬底''。这是个倍流畅图形的典范。在这幅巧妙的画中，可以有两种不等价的方式来刻划黑的部分：

⑴．作为白的部分的负空间；

⑵．作为白的部分的变形副本（通过将白的部分涂黑并平移而得到）。

（在FIGURE-FIGURE图形这种特殊情况下，两种刻划确实是等价的------但在大多数黑白图画中，两者不会等价。）将来在第八章，当我们构思我们的印符数论（即TNT）时，我们将希望能用两种类似的方式来刻划所有的数论假陈述的集合：

⑴．作为TNT的所有定理的集合的负空间；

⑵．作为TNT的所有定理的集合的变形副本（通过否定每个TNT定理而得到）。

但这个希望将破灭，因为：

⑴．在非定理组成的集合中有真理。

⑵．在所有被否定的定理组成的集合之外有``假理''。

在第十四章里你会看到这是为什么，以及是怎样发生的。同时，也请你琢磨琢磨用图形表达出来的这种境况（图18）。

\pandocbounded{\includegraphics[keepaspectratio]{./media/OEBPS/Images/img_018.png}}

图18．TNT系统中各种符号串之间的关系

最大的方框代表所有TNT符号串的集。次一级的方框代表所有良构的TNT符号串的集，在它里面有TNT的所有语句的集。现在事情开始有趣了。所有定理的集合画成了一棵树，树干代表公理的集合。选树来作象征是因为它展现出来的递归生长的模式：新的枝条不断地从旧的里面生长出来。手指一样的枝条探到了约束区域（真理集）的各个角落，然而永远无法完全覆盖这个区域。真理和``假理''之间的分界线有意地提示出一种曲折随机的海岸线，无论你如何仔细地察看，总是有更细微的结构，因而不可能用有穷的方式精确地加以描述（参见孟德尔布劳特所著《碎片》）。对应的树代表定理的否定的集：它们都是假的，但它们整个却无法充满假陈述的空间。{[}作者绘{]}

音乐中的图形与衬底

在音乐中也可以找到图形与衬底。类比之一就是旋律与伴奏之间的区别------因为在某种意义上讲，旋律总是处在我们注意力的前沿，而伴奏是第二位的。因此当我们在一部乐曲的较低声部发现可识别的旋律时会很惊奇。这种乐曲在巴罗克以后的音乐中不太常见。和声通常不被当作是前景。但在巴罗克音乐中------尤其是在巴赫的音乐里------各个声部，不论是高是低或在中间，都是起``图形''作用的。从这个意义上讲，巴赫的曲子可以称为``倍流畅的''。

音乐中还有另一种图形衬底的区分：强半拍与弱半拍之间的区分。如果你为一个小节里的音符打拍子：``一‡嗒，二‡嗒，三‡嗒，四‡嗒''，多数旋律音符会落在数字上，而不是``嗒''上。但有时候，一段旋律会刻意地落在各个``嗒''上，纯粹追求这种效果。例如，肖邦有几首钢琴练习曲就有这种旋律。巴赫的曲子里也有------特别是他的无伴奏小提琴奏鸣曲和帕蒂塔，以及无伴奏大提琴组曲。

在这些曲子里，巴赫用各种办法使两个或更多的曲调同时进行。有时他让独奏乐器用``双弦''奏法------两个音同时奏，有时他让一个声部落在强半拍上，而另一个声部落在弱半拍上，这样听起来就能把它们区分开来，听到两个不同的旋律盘旋交织，互相协调地融合在一起。不用说，巴赫还远不只停留在这种水平的复杂程度上\ldots\ldots{}

递归可枚举集之别于递归集

现在让我们把图形与衬底的概念带回到形式系统这个领域。

在前面的例子里，正空间的角色是由C型定理来扮演，负空间的角色是由带素数个短杠的符号串来扮演的。到目前为止，我们找到的用印符规则表示素数的仅有手段是把它们视为负空间。那么，是否有什么手段------我不在乎多复杂------能把素数作为正空间表示出来------就是说，作为某种形式系统的定理的集合？

这里，不同人的直觉给出不同的回答。我还能生动地记起当初为了体会正的刻划与负的刻划之间的区别，我是怎样地着迷而大伤脑筋的。那时我相当肯定地认为不仅素数，而且任何能被负地表示出来的数的集合都同样能正地表示出来。我这个信念背后的直觉可以用一个问题来表达：``一个图形与其衬底怎么会不带有完全相同的信息？''在我看来，它们体现了同样的信息，只是以互补的两种方式对信息编码而已。在你看来正确的回答是什么？

结果就素数而言我是对的，但对一般情形来说我错了。这使我很吃惊，而且就是到了今天我也仍继续感到吃惊。那是这样一个事实：

{存在一个形式系统，其负空间（非定理集）不是任何一个形式系统的正空间（定理集）。}

后来证明了，这个结果与哥德尔定理具有同样的深度------因此我的直觉出错是不奇怪的。我，正像二十世纪初的数学家们一样，期望着形式系统与自然数的世界比实际上更可预测。用更为技术性的术语来说，上面的结果就是：

{存在非递归的递归可枚举集。}

``递归可枚举''这个词（通常缩写为``r.e.''）是我们那个艺术上的``流畅可画出''概念的数学对应物------而``递归''则是``倍流畅''的对应物。一集符号串是``r.e.''的，意思是它能按照印符规则来生成------例如，C型定理的集合、WJU系统的定理集合------其实，任何形式系统的定理集都是这样。这可以和``图形''概念，作为``能按照艺术规则生成的线条的集合''（且不管那是些什么样的规则），作个比较。而一个``递归集''就像是个其衬底也是个图形的图形------不仅它r.e.，它的补集也r.e.。

上述结果导致：

{存在一些形式系统，它们没有用印符规则表述的判定过程。}

怎么会得出这样的结论呢？很简单。一个用印符规则表述的判定过程是一种分辨定理与非定理的测试方法。有了这样的测试，我们就能系统地生成所有的非定理。方法很简单：把所有的符号串排成一列，对它们逐个施用这个测试，除掉那些非良构串及那些定理。这就相当于一个生成非定理集合的印符方法了。但按照前面的那个论断（在这里我们当作信念而接受），对于某些系统这是不可能的。因此我们必须下结论：并非所有的形式系统都有用印符规则表述的判定过程。

假定我们找到了一个自然数的集合T（相应于``图形''），能用某种形式的方法生成------就像合数那样。假定它的补集是C（相应于``衬底''）------就像素数那样。T和C加在一起就是全部自然数了。我们知道一个产生T集中所有的数的规则，但我们知道没有同样的产生C集中所有的数的规则。了解下面这一点是很重要的：如果T集中的元素总能按长度递增的方式生成，那么我们总是能够刻划C。问题在于许多r.e.集的生成方式是以任意的次序添加元素的，因此你永远无法知道某个数被跳过之后过了很长时间会不会被收进来，如果你只要稍稍再等一会的话。

对于那个艺术问题：``所有的图形都是倍流畅的吗？''我们回答了``否''。现在我们看到了，对于数学中的类似问题：``所有的集合都是递归的吗？''我们必须同样地回答否。带着这样的观察，让我们回到那个费解的词``形式''。让我们再摆出我们的图形集T和衬底集C。我们可以承认集T中所有的数都有某种共同的``形式''------但对于集C中的数也能这么说吗？这是个奇怪的问题。当我们着手对付一个无限的集合------自然数------的时候，通过去掉某个子集而产生的洞大概是很难用什么显明的方式来定义的。因此它们或许不带有任何共同的特性或``形式''。分析到最后，用不用``形式''这个词成了口味问题------但思考这件事是很有意思的。也许最好是不去定义``形式''，而把它留给灵活的直觉。

这里有个供思考的谜题，它与上面谈的事情紧密相关：你能刻划出下面这样的整数集合（或者它的负空间）吗？

1 3 7 12 18 26 35 45 56 69\ldots\ldots{}

这个序列与FIGURE-FIGURE图形有何相似？

素数作为图形而非衬底

最后，让我们来看看生成素数的形式系统。它是怎么做到的？技巧是跳过乘法，直接把不可整除性当作要正面表示的东西。下面是一个公理模式和一条规则，它们产生的定理可以准确地表达一个数不整除(BZC)另一个数这个概念：

公理模式：xyBZCx，其中x和y是短杠符号串。

例如，-\/-\/-\/-\/-BZC-\/-，这里``-\/-''代替了x而``-\/-\/-''代替了y。

规则：若xBZCy是个定理，则xBZCxy是个定理。

如果你把这个规则用两次，你就能生成下面这个定理：

-\/-\/-\/-\/-BZC-\/-\/-\/-\/-\/-\/-\/-\/-\/-\/-\/-

这可以解释成``5不整除12''。但-\/-BZC-\/-\/-\/-\/-\/-\/-\/-\/-\/-\/-\/-不是个定理。要是你试图生成它，会出什么问题？

现在为了确定一个给定的数是素数，我们得建立某种关于它的不可整除性的知识。特别地，我们希望知道它不能被这样一些数整除：2、3、4、等等，如此下去，直到比那个数本身小1的数。但我们在形式系统里不能含糊不清地说``等等''之类的话。我们必须明白地说淸楚。我们大概会要有一种方式以这个形式系统的语言来说出``数Z在直到X的数内没有因子（记作zMYx）''，意思是在2和X之间没有能整除Z的数。这是能做到的，但要点技巧。如果你愿意，你可以自己想一想。

下面是解决方案：

规则：如果-\/-BZCz是个定理，则zMY-\/-是个定理。

规则：如果zMYx与x-BZCz都是定理，则zMYx-是个定理。

这两条规则把握住了``无因子性''这个概念。我们还需要做的仅仅就是说出：素数是那种在直到比它自己小1的数内没有因子的数：

规则：若z-MYz是个定理，则Pz-是个定理。

噢，我们可别忘了2是个素数！

公理：P-\/-。

到这里就算大功告成了。形式地表示素数性的要点，在于有一个对可整除性的测试，无需任何回溯就能工作。你稳步向前进展，先测试是否被2可除，然后被3，被4，如此下去。就是这种``单调性''或者说单向性------没有加长与缩短、增大与减小的交错出现------使我们把握住了素数性。然而正是形式系统的那种潜在复杂性，即它包含任意多的向前、向后的干扰，才导致了诸如\heititext{哥德尔定理}、图灵的停机问题、以及关于``并非所有递归可枚举集都是递归集''的事实等等这些限制性的结果。

\phantomsection\label{Dialog04.xhtml}{}

\chapter{对位藏头诗}

阿基里斯来拜访乌龟，看到乌龟家里有许多唱片。

\dialogname{阿基里斯：}\gebfont{ꪟ}赫有名的法国作曲家德彪西你知道吧？你这里有没有他的作品？

\dialogname{　　乌龟：}\heititext{赫}赫有名的得镖器？我可没听说过。我这儿倒收藏了不少``飞去来器''，就是那种扔出去能飞回来的飞镖。

\dialogname{阿基里斯：}\heititext{有}名极了，不过这没关系。你这里到底都有谁的作品？我想总该有巴赫的吧？

\dialogname{　　乌龟：}\heititext{名}曲我都有，巴赫的自然不用说了。你听过巴赫的曲子吧？那种美，那种技巧！我每次听他的曲子都有新的感受。

\dialogname{阿基里斯：}\heititext{的}确是技巧高超，百听不厌。哎，龟兄，你是从哪儿找来这些唱片的？

\dialogname{　　乌龟：}\heititext{德}意志，巴赫的故乡。到底还是德国人演奏的好。

\dialogname{阿基里斯：}\heititext{国}籍不重要，巴赫属于全世界。

\dialogname{　　乌龟：}\heititext{作}为一种个人嗜好，我还是宁愿听德国人的演奏。好了，咱们别谈这些了，有句拉丁格言，恐怕你也知道：口味无须争辩。哎，你知道我这些天在干什么吗？我在研究一种特殊的音乐。这种音乐能形成一个``家族''，我称之为``粉碎唱机的音乐''。

\dialogname{阿基里斯：}\heititext{曲}名就是这个？

\dialogname{　　乌龟：}\heititext{家}族音乐是许多相似但又不同的曲子的总合，并没有什么总的曲名。我说的``粉碎唱机的音乐''不是曲名，而是对这一大族曲子的描述。

\dialogname{阿基里斯：}\heititext{给}我讲讲是怎么回事吧，我现在一点也摸不着头脑。

\dialogname{　　乌龟：}\heititext{了}解这种音乐可不是一件简单的事情，这得从头说起。有一次螃蟹来我这里做客，当时我身体不太舒服------噢，对了，你见过螃蟹吗？

\dialogname{阿基里斯：}\heititext{我}没见过。不过已经久闻他的大名了，真希望能见到他。我本来应该有个机会能碰上他的，那是在一个叫凯瑟林的人家里。他家举办了一个音乐会，可我没能去成。我好像记得凯瑟林是个什么贵族，大概是个侯爵。

\dialogname{　　乌龟：}\heititext{侯}爵，没错，就是后知后觉的那个侯爵。巴赫的《哥德堡变奏曲》就是为他写的。

\dialogname{阿基里斯：}\heititext{世}上竟有这样不幸的事！这真太遗憾了，我不仅错过了螃蟹，还错过了那部著名的组曲！咳，我真是！

\dialogname{　　乌龟：}\heititext{达}观一些，阿基。你将来肯定会有机会听《哥德堡变奏曲》，见螃蟹也是迟早的事，说不定哪天在公园里我们就会碰上他。别这么垂头丧气的。你平日的那股灵气呢？

\dialogname{阿基里斯：}\heititext{灵}气？好吧，你接着说吧：螃蟹到你这儿来了，当时你不大对劲------

\dialogname{　　乌龟：}\heititext{感}到身体不太舒服，也正想找人聊聊，于是他跟我说起他新买的唱机------你知道，螃蟹这家伙特别容易对小玩艺着迷，这回轮到唱机了------他说他那唱机是一台完备的唱机，能重现任何音乐。我当然不信，可我没法说服他。他顽固地相信推销员对他说的每一个字！于是我回家后研究了一番，作了一支曲子，灌进唱片，带着去见螃蟹。这支曲子的名字是``我不能在唱机1上播放''，就是说\ldots\ldots{}

\dialogname{阿基里斯：}\heititext{\ldots\ldots{}}什么？那是个什么东西？

\dialogname{　　乌龟：}\heititext{在}唱机1上没法播放的曲子，就是这么回事。我建议他听听这支曲子，他很高兴地答应了，把唱片放上那个可爱的唱机。可是真不幸！刚刚播出几个音符，那唱机就开始剧烈地震颤起来，随后``砰''的一声巨响，唱机炸成了碎片，溅得满屋都是。唱片当然也跟着完了。

\dialogname{阿基里斯：}\heititext{此}刻螃蟹一定非常难受吧？你也太狠心了！

\dialogname{　　乌龟：}\heititext{我}不过是向他证明一下他的观念有多么荒唐罢了。那台唱机根本不可能重现我那支曲子，那种音乐会使它震碎自己。

\dialogname{阿基里斯：}\heititext{借}用一个术语，就是``自我震碎''了，对吗？

\dialogname{　　乌龟：}\heititext{用}得好！当然，也可以这么看这个问题：那个推销员多少有些吹牛，螃蟹买的那台唱机不是那么完备，它事实上不可能重现所有的音乐。

\dialogname{阿基里斯：}\heititext{他}后来明白这一点了吗？我想他不会就此罢休吧？

\dialogname{　　乌龟：}\heititext{的}确，他没认输，他后来\ldots\ldots 嗯，我想我该先讲讲我的曲子，是不是？

\dialogname{阿基里斯：}\heititext{对}对，你应该先跟我说说你的曲子是怎么回事。等一下，那种唱片你只有一张吗？我很想看看。你把它放在哪儿啦？

\dialogname{　　乌龟：}\heititext{位}子上。在那儿。\kaitismall{（阿基里斯找到了唱片。）}看唱片是看不出什么的，关键在于曲子的旋律。唱机播放这种旋律时会产生特殊的音响效果。我可以给你详细讲讲。在回访螃蟹之前，我去了螃蟹买唱机的那家商店，问清了制造那种唱机的工厂。然后我给制造商写了一封信，向他们索取设计图纸。得到图纸以后，我仔细分析了那种唱机的结构，发现只要附近有某种特殊的声音，就会引起唱机的共振，最终导致它破碎。

\dialogname{阿基里斯：}\heititext{技}术方面的问题你倒挺在行的。

\dialogname{　　乌龟：}\heititext{巧}妙地把这种声音组织成旋律，就是那首``我不能在唱机1上播放。''

\dialogname{阿基里斯：}\heititext{写}上``礼物''二字，送给螃蟹，这就是你干的好事！螃蟹真是太不幸了。后来呢？

\dialogname{　　乌龟：}\heititext{下}一次我见到他时，他说他又买了一台新的唱机。他不认输，他无论如何不相信那是唱机的问题。这台新唱机比上一台贵得多，推销员保证说如果发现有什么音乐那唱机无法重现，就付还他双倍的钱。螃蟹跟我说这些的时候兴高采烈，我答应了去看看他的唱机。

\dialogname{阿基里斯：}\heititext{一}定又是这么个过程：在你去之前，你写了一封信给制造商。然后你针对这个唱机的结构谱一支曲子，叫做``我不能在唱机2上播放''，到了螃蟹那儿以后，同样的事情发生了，唱机炸成了碎片，唱片也毁了。对不对？

\dialogname{　　乌龟：}\heititext{个}别时候你倒也挺聪明！你说的完全正确。看来你已经抓住了这件事的核心。

\dialogname{阿基里斯：}\heititext{对}我来说这不算什么。我要是螃蟹，早就看穿你这一套了，根本不会让你钻空子。我会设计一个能对付你这花招的唱机。唉！只可惜螃蟹想不到这一点。

\dialogname{　　乌龟：}\heititext{话}可不能这么说。我和螃蟹又较量了几个回合之后，他也明白了我的手段。于是他变得聪明了，准备\ldots\ldots{}

\dialogname{阿基里斯：}\heititext{并}不是螃蟹聪明。你那手段本来就没什么了不起的。螃蟹可以很容易就战胜你：比如说用一台低保真的唱机------无法严格重现那种能够毁坏自己的声音------你不就没辙了？

\dialogname{　　乌龟：}\heititext{嵌}上一块亮晶晶的小牌子，标明它是个``完备''的唱机，对吗？你怎么这么糊涂！一开始就确定了有些音乐不能重现，还算什么完备的唱机？你想想\ldots\ldots{}

\dialogname{阿基里斯：}\heititext{进}一步的解释让我来作：如果一个唱机------就说是唱机X吧------有足够的高保真度，那么当它试图播放曲子``我不能在唱机X上播放''时，恰恰就引起了那种粉碎它自己的共振\ldots\ldots 所以它不是完备的。然而唯一能绕开这种打击的方法------即让唱机X是低保真的------却更直接地表明了它的不完备。看来任何唱机都会有这样或那样的弱点，因此所有的唱机都是有缺陷的。可怜的螃蟹，无论是高保真唱机还是低保真唱机，他都是注定要失败的。

\dialogname{　　乌龟：}\heititext{他}的问题不在于使用了什么``有缺陷''的唱机。我不认为谈论唱机的``缺陷''有什么意义，一台唱机本来就不可能做你希望它做的所有事情。如果一定要说什么地方有缺陷，那也不是唱机，而是你对它们的期望------对于唱机能做些什么事情的期望。螃蟹恰恰就是充满了这类不现实的期望。

\dialogname{阿基里斯：}\heititext{的}确，不应该责怪唱机。我估计螃蟹买的也一定是好唱机，名牌的。

\dialogname{　　乌龟：}\heititext{名}牌不名牌无关紧要，反正是高保真的，因此注定要破碎。不过，螃蟹还是没认输。他知道了我谱曲的原理之后，想出了一个办法，他以为这样就能胜过我了。他写了封信给唱机制造商，描述了他自己发明的一种装置，让他们依法制造。螃蟹称之为``唱机欧米伽''。这是一种比一般唱机都复杂得多的唱机，复杂得以至于\ldots\ldots 怎么说呢？以至于\ldots\ldots{}

\dialogname{阿基里斯：}\heititext{字}斟句酌的干嘛？你就接着说吧，是什么样的？

\dialogname{　　乌龟：}\heititext{以}至于我似乎真的没办法了。唱机欧米伽和一般的唱机有很大区别，但表面上差不多。

\dialogname{阿基里斯：}\heititext{表}面上差不多但又有很大区别？那么就是内部构造很不一样了？给我详细讲讲吧，最好画个图。

\dialogname{　　乌龟：}\heititext{示}意图是这样。\kaitismall{（画了一个图，用笔指点着）}这是一部摄像机，连接着一台超微型计算机，都是嵌在唱机内部的，从外面看不出来。摄像机的任务是在播放一张唱片之前扫描其纹道，而计算机的任务是根据摄像机搜集的信息进行某种计算，精确地确定出这张唱片播放时会产生什么样的音响效果，判断它是否对欧米伽的唱机部分有害。

\dialogname{阿基里斯：}\heititext{我}想象不出判断出有害了又能怎么样。

\dialogname{　　乌龟：}\heititext{对}于那些有害的唱片，欧米伽自有一套聪明的办法。它有一个特殊的装置，能重新组织它的唱机部分，从而改变唱机的结构，使之播放那种唱片时不产生共振。这一切都在计算机的指挥下进行，精确可靠，因此任何一张唱片它都能播放------它总能选择一种不会产生共振的结构来播放一张唱片。

\dialogname{阿基里斯：}\heititext{他}可真有两下子！

\dialogname{　　乌龟：}\heititext{卓}绝的设计，是吧？

\dialogname{阿基里斯：}\heititext{越}来越邪乎了，这个老螃蟹！你彻底失败了吧，龟兄？

\dialogname{　　乌龟：}\heititext{才}没有呢！你恐怕不很了解哥德尔不完全性定理吧？按照这个定理，螃蟹再怎么耍花招也赢不了。

\dialogname{阿基里斯：}\heititext{能}赢！我不知道什么``疙瘩''不``疙瘩''，反正你肯定输了。嗯，这个关于``粉碎唱机的音乐''的小故事还真挺有意思。好了，你也不用再讲下去了，事情的结局是非常明显的：你老老实实认输，即便你是那么不愿意------这由不得你！

\dialogname{　　乌龟：}\heititext{由}不得我？\kaitismall{（打了一个呵欠）}算了算了，今天算了，你瞧，已经快半夜了，我得睡觉了，下次我再跟你说都是怎么回事吧。

\dialogname{阿基里斯：}\heititext{衷}心地感谢你。你的故事使我今天晚上很愉快。好，我回去了。\kaitismall{（刚走到门边，忽然停下来转回身）}我真糊涂！差点忘了。这是一个小小的礼物，我特地带给你的。\kaitismall{（递给乌龟一个整整齐齐的小包。乌龟打开一看，是一只高脚杯）}怎么样，不错吧？

\dialogname{　　乌龟：}\heititext{的}确好极了！多么奇特的一只高脚杯！哎，阿基，你知不知道，高脚杯这玩艺儿还是挺让我着迷的。只是近来景况不佳\ldots\ldots{}

\dialogname{阿基里斯：}\heititext{景}况不佳？我不明白。

\dialogname{　　乌龟：}\heititext{仰}头看看，那儿。\kaitismall{（阿基里斯仰起头，看见一个架子上摆了一排高脚杯。）}我收藏了这么多高脚杯，可它们总是有这样那样的缺陷。我可以告诉你我的心愿，不过你得为我保密。我一直试图找一只完备的高脚杯：它的形状必须没有任何缺陷。但我始终没能如愿。如果你送给我的这只高脚杯正是我要找的，那岂不是太妙了？告诉我，阿基，你怎么搞到它的？

\dialogname{阿基里斯：}\heititext{大}概不能告诉你，我很抱歉。这是我的一个小秘密。不过你应该知道这个高脚杯是属于谁的。

\dialogname{　　乌龟：}\heititext{家}里有这么多高脚杯，我不可能一一知道它们的来历，当然，这只高脚杯我就更不知道了。还是你来告诉我吧。

\dialogname{阿基里斯：}\heititext{也}好。它属于你最喜欢的作曲家巴赫。他有许多高脚杯，这一点不太为人所知------我给你的这个是他的最后一只高脚杯。

\dialogname{　　乌龟：}\heititext{许}多高脚杯中的最后一只，是吗？我的天！如果这真是巴赫的，那将是无价之宝了。可你怎么能肯定这是巴赫的？

\dialogname{阿基里斯：}\heititext{都}在上面写着呢。你看见刻在杯子上的那些精巧的字母了没有？

\dialogname{　　乌龟：}\heititext{记}载着巴赫的名字！这只高脚杯的确是巴赫的，多不寻常啊！\kaitismall{（小心翼翼地把这只高脚杯放在了架子上。）}哎，阿基，你知不知道巴赫的名字，也就是B-A-C-H，每个字母各是一个音名？

\dialogname{阿基里斯：}\heititext{得}排除H。你知道，音名只是从A到G。

\dialogname{　　乌龟：}\heititext{他}的情况有些特殊，你知道他是德国人，在德国，音名的记法和我们略有不同。我们用B的地方，他们用H，而我们用降B的地方，他们用B。比如说，我们说巴赫的``B小调弥撒曲''，他们却说``H小调弥撒曲''。明白了吗？

\dialogname{阿基里斯：}\heititext{曾}经有人告诉过我这些，我怎么给忘了。现在我想起来了。这么说，巴赫的名字实际上是一段旋律喽？

\dialogname{　　乌龟：}\heititext{把}那四个音符连续奏出，的的确确就是一段旋律。事实上，巴赫把这段旋律用意微妙地写进了他精雕细琢的一部作品------即《赋格的艺术》中最后一支四声部三重赋格里。这支曲子是巴赫所写的最后一支赋格。我第一次听它时，全然不知它会如何结尾。可突然，没有任何先兆，曲子戛然而止。然后\ldots\ldots 死一般的寂静。我立刻体会到巴赫就是在这儿死去的。那是一个难以名状的悲伤时刻，我的心都------碎了\ldots\ldots B-A-C-H是那部赋格的最后主题，它隐藏在这支曲子里。虽然巴赫没有明白地表明这一点，但如果你知道这支曲子是怎么回事，你就会很容易发现了。

\pandocbounded{\includegraphics[keepaspectratio]{./media/OEBPS/Images/img_019.png}}

图19．巴赫《赋格的艺术》最后一页

在原手稿上，巴赫的儿子卡尔·菲利普·埃玛努厄尔写道：``请注意，在这支赋格中，当名字B.A.C.H.作为对题被引入时，作者死去了。''{（方框标出的即是B-A-C-H。）}我把巴赫的这最后一支赋格的最后一页当作他的墓志铭。{[}乐谱由唐纳德·伯德的程序``斯马特''印出。该程序是在印第安纳大学开发的。{]}

\dialogname{阿基里斯：}\heititext{他}为什么要那样呢？

\dialogname{　　乌龟：}\heititext{的}确是个谜。无论如何，这是部杰作，而那段旋律也成为名句了。

\dialogname{阿基里斯：}\heititext{名}句？有这么说的吗？

\dialogname{　　乌龟：}\heititext{字}眼不重要，我相信你明白我的意思。

\dialogname{阿基里斯：}\heititext{写}下来的精彩句子才称作``名句''，那是文学领域的事，你可别乱用。

\dialogname{　　乌龟：}\heititext{进}入到文学领域你会发现更多这样的手法。比如说在刘易斯·卡罗尔的诗中，就常常把一些词或名字藏在每行的第一个字里。用这种方式藏有信息的诗称作``藏头诗''。

\dialogname{阿基里斯：}\heititext{一}提起卡罗尔，我就会想起一个一直困扰我的问题：他是靠什么搞出了那么多妙不可言的玩艺儿？小时候我读过他的``炸脖\gebfont{ꪡ}''，真是太绝了！

\dialogname{　　乌龟：}\heititext{首}先是靠他的天才，我看是这样。这个我们先不管它，总之卡罗尔写过不少藏头诗。巴赫偶尔也写一点藏头诗，这并不奇怪。毕竟，赋格与藏头诗都同时有好几层隐藏着的意思，在这一点上两者颇为相似。多数的藏头诗只有一层隐含的意思，但是没有理由不让你搞出个双层结构------一首藏头诗中的藏头诗，或者也可以搞一个``头尾反藏诗''------把一句话最后一个字和第一个字连起来时，是一个有意义的词。天哪！形式本身具有的可能性几乎是无限的！赋格里就充满了这种技巧。

\dialogname{阿基里斯：}\heititext{赋}格的情况更为复杂，是吧？

\dialogname{　　乌龟：}\heititext{格}式更为多样，技巧要求更高。其实即使是写对话也可以使用这种技巧\ldots\ldots 嗯，我忽然有了个想法。你设想一下，一位对话作者为了表示他对巴赫卓越才能由衷的景仰而写下一篇对话，对于这位作者来说这样做是不是更合适：即在对话中借用巴赫的对位技巧嵌进自己的名字------或者巴赫的名字？咳，何必操那份心呢，谁愿意写这样一篇对话就写去吧。我们还是来看看巴赫那个作为旋律的名字吧。你知不知道，B-A-C-H这个旋律如果上下颠倒地从后向前演奏，会和原先正着演奏时完全一样？

\dialogname{阿基里斯：}\heititext{曲}调怎么能上下颠倒演奏？从后向前我倒知道怎么回事------就是H-C-A-B------可上下颠倒会是什么样？你一定是在跟我开玩笑。

\dialogname{　　乌龟：}\heititext{的}确能颠倒着演奏，真的。用颠倒的方式演奏尾声正合适。

\dialogname{阿基里斯：}\heititext{尾}声怎么就合适？你能马上给我表演一下吗？

\dialogname{　　乌龟：}\gebfont{ꪞ}不得呢。你等一下，我去拿小提琴------\kaitismall{（走入另一个房间，不一会，拿着一把古雅的小提琴回来了）}------现在我来把这一段从前向后，从后向前以各种方式给你演奏一下，听好啊\ldots\ldots{}

\hfill\break

乌龟从架子上取下那只高脚杯，恭恭敬敬地把它倒置在谱架上，开始演奏B-A-C-，当他拉到最后一个音符H时，突然，没有任何先兆，高脚杯炸成了碎片，飞溅到乌龟跟前的地板上。然后\ldots\ldots 死一般的寂静。

\phantomsection\label{Chapter04.xhtml}{}

第四章

\chapter{一致性、完全性与几何学}

隐含意义与显明意义

在第二章里我们看到，当规则支配的符号与真实世界中的事物之间有个同构时，意义------至少在形式系统这种相对简单的语境里------是如何出现的。一般说来，这个同构越复杂，为从符号中抽出意义所需的``设备''------硬的和软的------就越多。若一个同构很简单（或很熟悉），我们倾向于说它相对于我们的意义是显明的。我们不必注意这个同构就见到了意义。最明显的例子是人类的语言，人们常常把意义归于词本身，丝毫没有意识到使词具有意义的那个非常复杂的同构。这是个很容易犯的错误：把全部意义归于对象（即词），而非对象与真实世界的联结。这可以与一种朴素信念作个比较：认为任何两个物体的碰撞都必然附带有声音。这是个错误的信念。当两个物体在真空中相撞时，就不会有任何声音。道理是一样的。错误来自把声音完全归于碰撞，而未看到把声音从对象带给耳朵的媒介的作用。

以前，我让``同构''这个词带上引号是表明必须要留有余地。支持着对人类语言进行理解的符号操作过程，比典型的形式系统中的符号操作过程要复杂得多。因此我们要想继续把意义当作以同构为媒介的，我们就得对于同构能是什么采取更为灵活的看法，事实上在我看来，回答``什么是意识？''这一问题的关键之点，是揭示出作为意义基础的``同构''的实质。

《对位藏头诗》的显明意义

以上都是在准备对《对位藏头诗》的讨论------一种对于意义层次的研究。这篇对话既有显明的意义也有隐含的意义。最显明的意义无非是它所讲述的故事。严格地说，这个``显明的''意义是极度隐含的，因为面对白纸上的那些黑色印记，要理解故事中的事件，需要复杂得难以想象的大脑过程。尽管如此，我们还是把故事中的事件当作对话的显明意义，并假定每位读者都多少是用同样的``同构''从纸上的印记中抽取出这层意义的。

即便如此，我还是想把那个故事的显明意义说得更明白些。首先我想谈谈唱机与唱片。主要的问题是唱片上的纹道有两层意义。第一层意义是音乐意义。那么什么是``音乐''？------是空气中一系列的震颤，还是大脑中相继的情感反应？两者都是。但在有情感反应之前，先得有震颤。震颤是由唱机从纹道里``抽''出来的：一个相对说来直接了当的机制。事实上你都可以用一根针来做到，只需把针沿着纹道划就行。这个阶段之后，耳朵把震颤转换成大脑中听觉神经元的发射。以后便是大脑中的一系列步骤，逐渐地把线性的震颤序列变换成相互作用着的情感反应的复杂模式------太复杂了，我们不能再深究，虽然我很想这样。所以我们还是满足于把空气中的声音当作纹道的``第一层''意义。

纹道的第二层意义是什么？是在唱机中引起的震颤序列。这层意义只能是在第一层意义从纹道上被抽出来之后才有，因为是空气中的震颤引起了唱机中的震颤。所以，第二层意义有赖于相继的两个同构：

⑴．任意的纹道模式与空气震颤之间的同构；

⑵．任意的空气震颤与唱机震颤之间的同构。

这两个相继的同构的示意图见图20。注意是同构1使第一层意义出现。第二层意义比第一层意义隐含一些，因为它的媒介是两个相继同构。就是这第二层意义``倒戈''，才使唱机破碎的。有趣的是第一层意义的产生迫使第二层意义也要同时产生------决不会只有第一层意义而没有第二层。所以，是唱片的隐含意义又转了回来，毁掉了唱片。

\pandocbounded{\includegraphics[keepaspectratio]{./media/OEBPS/Images/img_020.png}}

图20．哥德尔定理背后的原理的形象化表现

两个映射背靠背具有了出乎意料的飞去来器效果。第一个映射是从纹道模式到声音，由一个唱机实现。第二个------很平常，但却常常被忽视------是从声音到唱机的震颤。注意第二个映射独立于第一个映射而存在，因为附近的任何声音------不仅是唱机自己产生的声音------都会引起这种震颤。哥德尔定理转述在这里就是：对任何一个唱机，都有它不能播放的唱片，因为后者会导致前者的间接自摧毁。{[}作者绘{]}

高脚杯的情形是一样的，同样可以这么说。但有个区别：字母表中的字母到音符上的映射是另外一层同构，我们可以称之为``转录''。这个同构之后跟着的是``翻译''------音符到乐音的转换。再之后，震颤反作用到高脚杯上，恰如反作用到那一系列逐步升级的唱机上一样。

《对位藏头诗》的隐含意义

那篇对话的隐含意义又如何呢？（是的，有不止一层。）其中最简单的已在前面指出了------就是，对话前半部分与后半部分大体上是同构的：唱机变成了小提琴，乌龟变成了阿基里斯，螃蟹变成了乌龟，纹道变成了刻写的签名，等等。注意到了这个简单的同构，你还可以更进一步。注意在故事的前半部分乌龟是全部恶作剧的祸根，在后半部分他却成了受害者。于是你便看到正是他自己的方法掉转头倒戈回来打击他！你是不是想起了唱片上音乐的倒戈------或高脚杯铭文的倒戈------或者也许是乌龟收藏的那些飞去来器？的确如此。这故事讲了两个层次上的倒戈，如下所示：

第一层：倒戈的高脚杯与唱片；

第二层：乌龟那个利用隐含意义引起倒戈的险恶手段------它本身也倒戈。

于是我们甚至可以在故事的两个层次上做个同构，把唱片与高脚杯自食恶果地摧毁自身这一过程，和乌龟自己那残忍的手段在最后自食恶果地指向他自己这一过程等同起来。这样看来，故事本身就是它所讨论的倒戈的一个实例。因此我们可以把《对位藏头诗》当作一种间接自指来看待：它自身的结构同构于它所描述的事件（恰恰相当于高脚杯和唱片通过播放与引起震颤这两个背靠背的同构来隐含地指向自身）。当然，人们可以读着那篇对话而没有注意这一事实------但它总是在那里的。

《对位藏头诗》与哥德尔定理之间的映射

现在你是不是觉得有点眼花缭乱了？然而最好的还没来呢（实际上，某些隐含意义的层次甚至不会在这里讨论------留着让你去找出来）。写这篇对话的最深刻的原因在于阐释哥德尔定理。正如我在导言中所说的，哥德尔定理很大程度上是依赖于数论的陈述有两个不同层次上的意义。《对位藏头诗》的两部分分别都是一个哥德尔定理的``同构复本''。由于这个映射是那篇对话的中心思想，而且相当精致复杂，所以我仔细地把它列成了下表。

\begin{longtable}[]{@{}lll@{}}
\toprule\noalign{}
\endhead
\bottomrule\noalign{}
\endlastfoot
唱机 & \gebfont{⇔} & 数论的公理系统 \\
低保真唱机 & \gebfont{⇔} & ``弱''公理系统 \\
高保真唱机 & \gebfont{⇔} & ``强''公理系统 \\
``完备的''唱机 & \gebfont{⇔} & 数论的完全系统 \\
唱机的``图纸'' & \gebfont{⇔} & 形式系统的公理与规则 \\
唱片 & \gebfont{⇔} & 形式系统的符号串 \\
可以播放的唱片 & \gebfont{⇔} & 公理系统的定理 \\
不能播放的唱片 & \gebfont{⇔} & 公理系统的非定理 \\
声音 & \gebfont{⇔} & 数论的真陈述 \\
可重现的声音 & \gebfont{⇔} & 系统中经过解释了的定理 \\
不可重现的声音 & \gebfont{⇔} & 非定理的真陈述 \\
曲名：``我不能在唱机X上播放'' & \gebfont{⇔} & 哥德尔串的隐含意义：``我不能在形式系统X中导出'' \\
\end{longtable}

这还不是哥德尔定理与《对位藏头诗》之间那个同构的全部，而只是其核心部分。如果你到现在还没有完全把握哥德尔定理，那也不必着急------还要过几章我们才见到它呢！不过，读了那篇对话，你就已经尝到些哥德尔定理的味道了，但这不必就一定得理解它。现在我不再多说了，留给你去寻找《对位藏头诗》中其它类型的隐含意义。``觅之，自有所获''！

赋格的艺术

现在来说几句《赋格的艺术》。这是巴赫在他生命的最后一年写下的，由十八支基于同一主题的赋格曲组成。很明显，《音乐的奉献》的创作给了巴赫灵感。他决定在简单得多的主题上谱写另外一集赋格曲，以表明形式中所固有的全部可能性。在《赋格的艺术》中，巴赫以可能有的最为复杂的方式使用了一个非常简单的主题。整部作品都是同一个调。其中多数赋格曲有四个声部，逐渐地增加复杂程度及表现深度。接近最后，其错综与复杂上升到了这样的高度，以致于人们都怀疑他无法再驾驭了。然而他做到了------除了最后那首《对位曲》。

导致《赋格的艺术》的中断（也可以说，巴赫生命的中断）的情境是这样：几年来巴赫的视力给他带来不少麻烦，他希望做个手术。手术是做了，但没多大帮助。结果，在他生命最后一年的那些较好的时光里他失去了视觉。然而这并没有阻止他以充沛的精力去从事他那不朽的事业。他的目标是建立一个对于赋格曲创作的完全的展示，多重主题的使用是其中的一个重要的方面。在他计划里的倒数第二支赋格曲中，他把自己的名字做成音符插在里面作为第三主题。然而正在这么做的时候，他的健康情况恶化了，他被迫放弃了继续从事他的事业的抱负。病中，他想办法向他女婿口授了最后一首赞美诗前奏曲，对此巴赫的传记作者福凯尔写道：``每当我演奏一遍，都会被那虔诚的达观与奉献所深深打动，以致于我很难说我会更怀恋哪一个------这首赞美诗，还是最后那支赋格曲的结尾。''

有一天，没有任何先兆，巴赫恢复了视力。但几小时以后他的病又发作了。十天以后，他去世了，留下不完全的《赋格的艺术》让人们沉思。这会是由于巴赫达到了``自指''吗？

哥德尔的结果造成的问题

乌龟说没有一个足够强有力的唱机会在下述意义上是完备的：它能重现唱片上所有可能的声音。哥德尔说没有一个足够强有力的形式系统会在下述意义上是完备的：能够把每一个真陈述都作为定理而重现在该系统中。但正如乌龟关于唱机所指出的那样，只是当你对形式系统应该能做到什么抱有不现实的期望时，这件事才像是个缺陷。不过，数学家们正是带着这种不现实的期望进入二十世纪的，他们认为公理化推理是解决一切病症的良方。1931年他们发现了不是这样。对于任何一个形式系统，真理超出该系统所规定的定理资格这件事，被称作该系统的``不完全性''。

哥德尔的证明方法中最绕人之处是他所使用的种种推理方法看上去无法被``封住''------它们拒不卷入任何形式系统。于是，初看起来，哥德尔似乎是发掘出了以前不被人知，但却意味深长的人类推理与机械推理之间的区别。这种生命系统的能力与无生命系统的能力间的差异，在真理概念与定理资格概念之间的差异上反映了出来\ldots\ldots 或至少这是一种``浪漫''地看待这个问题的方式。

修改了的pq系统与不一致性

为了更现实地考察这个问题，有必要更深入地看一看在形式系统中，意义为什么以及怎样由同构来传送。但我相信这将导致更为浪漫地看待这个问题。所以现在我们来着手研究意义与形式间关系的某些进一步的方面。第一步是稍稍修改一下我们的老朋友pq系统，来做出一个新的形式系统。我们另外增加一个公理模式（保留原来那个，以及那唯一一条规则）：

公理模式：若x是一个短杠串，则xqxp-是一个公理。

于是很清楚，-\/-q-\/-p-在新系统中是个定理，-\/-\/-q-\/-p-\/-也是。像从前一样，它俩的解释分别为``2等于2加1''和``3等于2加2''。可以看出我们这个新系统将包括很多假陈述（如果你把符号串当作陈述的话）。这样，我们的新系统与外部世界是不一致的。

好像这还不够糟，我们的新系统还有内部问题。它含有彼此不协调的陈述，诸如-\/-q-p-（一条旧公理）和-q-p-（一条新公理）。因此我们的新系统在第二种意思上也是不一致的：内部的不一致。

那么，在这种关头，唯一合理的做法是不是就是把我们的新系统整个扔掉？很难说。我有意用障眼法来摆出这些``不一致性''：也就是，我尽可能地把一些糟糕的论证弄得像那么回事，目的在于使人误入岐途。事实上，你大概已经发现了我的错误。关键的错误发生在我不加思索地对新系统采用了同样的解释方法：和我对旧系统所使用的解释方法完全一样。回想一下，在上一章里用这种解释方法的时候，唯一的理由是这种解释使得符号操作同构于它们所对应的概念。如果你修改了支配系统的规则，你注定会无可挽救地破坏那个同构。这样看来，前面几段里的那些令人沉痛的问题是虚构的。只要重新解释系统中的某些符号，那些问题马上就消失了。注意我说的是``某些''：并不一定所有的符号都得映射到新的概念上。一部分符号可以完好地保持它们的``意义''，而其它的符号则改变了意义。

重新获得一致性

例如，假定我们只重新解释符号q，其它符号都保留原样。比如说，我们用短语``小于或等于''来解释q。现在，我们那两个``矛盾的''定理-q-p-和-\/-q-p-成了无害的两句话：``1小于或等于1加1''，``2小于或等于1加1''。我们同时摆脱了（1）与外部世界的不致性，和（2）内部的不一致性。我们的新解释是个有意义的解释。当然，原来那个是无意义的，就是说，它对于新系统来说没有意义。对于原来那个pq系统，它挺好。但现在看来，把它用在新的pq系统上过于随意，也没有道理，就像用``马-苹果-幸福''来解释旧的pq系统一样。

欧几里德几何的历史

虽然我力图使你在没有戒备的情况下，对怎样用语言来解释符号感到了意外，但看来似乎是一旦你掌握了其中的要领，学会这种技术并非难得不得了。事实上也的确不难。然而这却是整个十九世纪数学的最深刻的教训之一！那种数学起始于欧几里德。大约在公元前300年左右，他系统化地整理了他那个时代人们所知道的关于平面和立体几何的知识，其结果就是欧几里德的《几何原理》。这个几何学体系是如此稳固，几乎成了两千多年来几何学的圣经------所有时代维持最长久的著作之一。这是为什么呢？

主要原因在于欧几里德是数学中严格性的创始人。《几何原理》从非常简单的概念、定义等开始，逐渐地建立起许多结果，形成了一个庞大的体系。这个体系是精心组织的，其中给出的任何一个结果都仅仅依赖于出现在它之前的结果。这样，这种数学工作就有了一个确定的格式，其建构的艺术使之强壮并坚实。

\pandocbounded{\includegraphics[keepaspectratio]{./media/OEBPS/Images/img_021.jpg}}

图21．巴别塔

艾舍尔（木刻，1928）

不过，这种建构的艺术是另外一种类型，不同于------比如说------摩天大楼的建筑（见图21）。对于后者来说，能够矗立在那里，就足以证明它的各个结构要素是紧密组合起来的。但对于一本几何学的书来说，当其中的每个命题都声称是从前面的命题逻辑地推导出来时，若某个命题的证明是无效的，并不会显出整个建构会倒坍。这里横梁和立柱不是物质的，而是抽象的。事实上，在欧几里德的《几何原理》中，用来建构证明的东西是人类语言------一种充满隐患的复杂又难以捉摸的通讯媒介。那么，《几何原理》的建构聚合力是什么呢？能肯定它是由坚固的结构要素组合起来的吗？或者它也许有结构上的弱点？

我们用的每个词对于我们都有某种意义，在我们使用这个词的时候引导我们。越是普通的词，我们由它引起的联想就越多，其意义也就扎得越深。所以，如果某人对一个普通的词下个定义，指望我们遵守这个定义，我们肯定不这么做。相反，我们会在很大程度上无意识地按照另一些东西的引导去做，那些东西是我们的心智在我们与这个词相关联的储备中所发现的。我提及这些是因为这就是欧几里德的《几何原理》中所发生的问题，在那里他试图对诸如``点''、``直线''、``圆''等等普通平常的词下定义。你怎样才能给一个人人都已有了清晰概念的东西下定义？唯一的途径是假定你能清楚地表明你的词是当作术语的，不与同样拼写的日常用词相混淆。你必须得强调：与日常用词的联系仅仅是提示性的。可是，欧几里德没这么做，因为他觉得他的《几何原理》中的点和线的确就是现实世界中的那种点和线。所以，由于没有明确截断所有联系，欧几里德使得读者能自由地发挥他们的联想力了\ldots\ldots{}

这看起来几乎没有秩序可言了，而且对欧几里德也有点不公平。他确实制定了公理，或说公设，可用来证明命题。事实上，除了这些公理公设，别的东西一律不得使用。而就是在这里他失误了：他使用那些日常用词的一个不可避免的后果就是那些词生出的一些意象潜入了他给出的证明。不过，如果你读《几何原理》中的证明，无论如何别指望发现推理中扎眼的``跳跃''。正相反，它们非常精细，因为欧几里德是个思想敏锐的人，在他那里不会出现任何头脑简单所导致的错误。尽管如此，缝隙还是有的，造成了一部经典著作中轻微的瑕疵。但不应抱怨这一点，只应该体会到绝对严格与相对严格之间的差别。毕竟，在欧几里德写下他的著作两千多年以后，他缺乏绝对的严格性这一点导致了数学中一些最富成果的新分支。

欧几里德给出了五条公设作为几何学这个无穷级大厦的``地基''------他的《几何原理》只包括了这个大厦的前几百级。前四条公设相当优美简洁：

⑴．一条直线段可以联接两个点。

⑵．一条直线上任何一条直线段可以无限延伸。

⑶．给定一条直线段，可以以一个端点为圆心，以此线段为半径做一个圆。

⑷．一切直角都彼此相等。

然而第五条却不那么雅致：

⑸．如果两条直线与第三条直线相交时，在第三条直线的某一侧三条线所夹的内角之和小于两个直角的和，则那两条直线沿着这一侧延伸足够长之后必然相交。

虽然欧几里德从未明白地说过，但他确是把这条公设看得比其它几条地位低些，因为在前二十八个命题的证明里，他想法避免了使用它。这样，前二十八个命题属于所谓``四公设的几何学''------几何学中的一部分，可基于《几何原理》的前四条公设而导出，无需第五条公设的帮助（常常也称作``绝对几何学''）。无疑，欧几里德会发现：更为可取的是证明这只丑小鸭，而不必设定它。但他没找到证明，所以采纳它了。

但是欧几里德的后继者们在不得不设定这第五公设时，并不比他自在些。多少个世纪里，数不清的人花费了他们生命里数不清的岁月，试图去证明第五公设本身便是四公设几何学的一部分。到1763年，至少发表了二十八种不同的证明------都是错的！（都被一个名叫克吕格的人在一篇论文里批驳了）。所有这些错误的证明都涉及到了日常直观与严格的形式化属性的混淆。这么说是保险的：在今天几乎所有这些``证明''对我们都没有什么数学的或历史的价值了------除了几个例外。

非欧几里德几何面面观

济罗拉莫·萨彻利（1667-1733）生活在巴赫的时代。他有个抱负，要使欧几里德不带任何瑕疵。以他早期在逻辑方面的研究为基础，他决定试着从一个新颖的角度对那个著名的第五公设进行证明：假定你设定它的反面，然后以这样一条公设作为你的第五公设开始推演几何学\ldots\ldots 肯定不久之后你会制造出矛盾。因为没有任何数学系统能支持矛盾，你就表明了你自己的那个第五公设是不可靠的，于是表明了欧几里德的第五公设是可靠的。这里我们无需深究细节。指出下面这一点就够了：萨彻利以他高超的技巧推演出了``萨氏几何学''的一个又一个命题，最后终于变得厌倦了。有一个时候，他认定他达到了一个``与直线的本质相抵触''的命题。那正是他一直期望的------在他看来，那是个寻找已久的矛盾。就在那个时候他发表了他的研究结果，题为《不带任何瑕疵的欧几里德》，随后就不干了。

但就是因为这么做，他剥夺了自己身后的殊荣，因为他无意中发现了后来所谓的``双曲几何学''。五十年以后，兰伯特重复了这种``失之交臂''，这次，假如能这么说的话，甚至更接近了。最后，兰伯特之后四十年，也即萨彻利之后九十年，非欧几何学被认可了------一个真正的几何学新品种。这是自古以来单一的数学的一个分叉点。1823年，作为那种无从解释的巧合之一，非欧几何同时被两个人发现了。他们是匈牙利数学家雅诺（或约翰）·鲍耶，二十一岁，和俄国数学家尼古拉伊·罗巴切夫斯基，三十岁。具有讽刺意味的是，在同一年法国伟大的数学家阿德里安-马利·勒让德想出了他认定是对欧几里德第五公设的一个证明，他在相当程度上是沿着萨彻利的路子。

顺便指出，鲍耶的父亲法卡斯（或沃夫岗）·鲍耶，伟大的高斯的一个亲密朋友，投入了大量精力试图证明欧几里德的第五公设。在一封给他儿子雅诺的信中，他努力劝阻他不要去思考这种东西：

你一定不要去探究平行公设。我深知这条路通向哪里。我曾横穿于这无尽的黑夜，湮灭了我生命中所有的光明与欢乐。我恳求你放下平行公设的研究\ldots\ldots 我曾想为了真理而牺牲自己。我准备好成为一个殉教者，除去几何学的瑕疵，使之纯净而奉还给人类。我进行了巨大的、难以估量的努力，我得到的结果远比其他人好得多，却依然没达到完全的满意。当一个人彻底离开日常之琐碎，他就转向了崇高之最。当我看出没人能达到这黑夜的尽头时，我转回身了。我转回身来没有慰藉，却带着对自己以及全部人类的怜悯\ldots\ldots 我旅经了这地狱般的死海里的所有暗礁，总是帆破桅折地回来。就是从那时起，我开始衰老，性情也毁了。我不加考虑地用我的生命和幸福去冒险------或者光荣如凯撒，或者一无所有。\footnote{}

但是后来，当他确信他的儿子的确``有了东西''，他便催促他发表，正确地迎来了那个同时性。这种同时性在科学发现上真是屡屡发生：

当时机成熟时，那些东西在不同的地方出现，就像早春时的紫罗兰。\footnote{}

非欧几何学的情形真真确确就是这样！在德国，高斯自己和其他一些人都多少独立地触及了非欧几何观念。其中包括一个律师施外卡特，他于1818年寄给高斯一篇文章，描述了一种新的``星际''几何学；施外卡特的侄子陶利努斯做出了非欧几里德三角学；高斯的学生瓦赫特尔，他发现了非欧几何学的几个很深的结果，死于1817年，当时二十五岁。

理解非欧几何的线索是``率直对待''来自像萨彻利和兰伯特的几何学里的命题。只是当你没能摆脱先入为主的``直线''观念时，萨氏命题才``与直线的本质相抵触''。而若你能使自己摆脱那些先入的印象，只让``直线''是那种满足新命题的东西，那你就达到了一个崭新的视点。

未定义项

这听起来很熟悉。具体地说，让人想起pq系统以及它的变种，其中符号借助它们在定理中扮演的角色而获得被动意义。符号q尤其令人感兴趣，因为它的``意义''随着新公理的添加而改变了。人们可以以完全相同的方式让``点''、``线''等等的意义由它们出现于其中的定理（或命题）的集合来决定。这是非欧几何学发现者们的重大体会。他们以不同的方式否定欧几里德第五公设，并遵循其后果，从而发现了几种不同的非欧几何学。严格地说，他们（包括萨彻利）并不是直接否定第五公设，而是否定了一个等价的公设。它被称为平行公设，其叙述如下：

给定任一直线和不在直线上的一点，存在有一条，且仅仅存在一条通过那个点，且永不与前一条直线相交的直线，无论两直线延伸多远。

第二条直线就称为与前一条直线平行。如果你断言没有这样的直线存在，那么你得到的是椭圆几何学；如果你断言至少有两条这种直线存在，你得到的是双曲几何学。顺便指出，这样的变种还被称为``几何学''的理由是核心要素------绝对的或四公设几何学------嵌于其中。就是由于有这最小限度的核心，把它们当作对某种几何空间性质的描述才有意义，尽管这种几何空间不像平常的空间那么直观。

实际上，椭圆几何学是很容易视觉化的。所有的``点''、``线''之类东西都可认作是普通球面上的东西。让我们来规定一下：作为术语时，我们写``\heititext{点}''，作为日常用语时，我们写``点''。现在我们可以说一个\heititext{点}由一对球面上的对径点（球的一条直径的两端点）组成。一条\heititext{线}是球的一个大圆（一个圆，像赤道那样，以球心为其圆心）。在这种解释下，椭圆几何学的命题虽然也含有``\heititext{点}''、``\heititext{线}''之类的词，说的却是球面上发生的事情，而不是平面上的。注意：两条\heititext{线}总是相交于两个点，两个在球面上对称于球心的点------也就是恰好相交于一个\heititext{点}！而且正像两条线决定一个点，两个点也决定一条线。

这种对待``\heititext{点}''和``\heititext{线}''这些词的方法------即认为它们仅有的意义是那些它们出现于其中的命题灌注进去的------使我们向几何学的完全形式化前进了一步。这种半形式化的版本仍然使用了汉语中的许多词，带有它们通常的意义（像``一个''、``如果''、``和''、``连接''、``有''等词），虽然特殊的词，如``\heititext{点}''和``\heititext{线}''，其日常意义已被抽掉了。后者于是称作未定义项。未定义项，像pq系统中的p和q，在某种意义上是有定义的：它们被隐含地定义了------由它们出现于其中的那些命题全体定义，而不是在一个定义中显明地定义。

人们可以主张说未定义项的完整定义仅驻存于公设中，因为导出的命题已是隐含于公设中了。在这种观点下，公设隐含地定义了所有未定义项，这些定义体现于全部未定义项的相互关系之中。

多重解释的可能性

几何学的完全形式化需要一个激烈的措施：使每个词项都成为未定的------就是，把每个词项都变成形式系统中``没有意义''的符号。我在``没有意义''上加引号是因为，你也知道，那些符号自动地具有了从它们出现于其中的定理而来的被动意义。不过，人们是否发现了那些意义却是另外一个问题了，因为那将要求寻找这样一集概念，它们能通过同构联系于形式系统中的符号。如果一个人带着把几何学形式化的目的起步，他大概是对每个符号都有了一个意向中的解释，使被动意义铸入系统。这就是我开始构造pq系统时对p和q所做的。

但也许还有其它潜在地可接受的被动意义未被人注意。例如在原来那个pq系统中令人意外地把p解释成``等于''，把q解释成``减''。虽然这个例子很平凡，却包含了这样一个想法的核心：符号可以有许多有意义的解释------寻找它们是取决于观察者的。

我们可以用术语``一致性''来概括到目前为止我们的考察。我们的讨论开始于虚构一个显得不一致的形式系统------内部不一致，同样与外部世界也不一致。但过了一会我们就撤回了这种说法，我们意识到了我们的错误：我们为符号选择了不幸的解释。改换了解释以后，我们重新获得了一致性！现在很清楚了：一致性不单是形式系统的性质，还依赖于为之提出的解释。同理，不一致性也不是任何形式系统的固有性质。

各种各样的一致性

我们一直在谈论``一致性''和``不一致性''，却还没定义它们。我们只是依靠古老悠久的日常观念。现在让我们来严格地定义一个形式系统（加上一个解释）的一致性：这是指其中每个定理经解释后都成为一个真陈述。如果至少有一个经解释后的定理是假陈述，我们就说出现了不一致性。

这个定义似乎是在谈论系统与外部世界的不一致性。那么，内部的不一致性怎么样呢？大概是如果一个系统含有两个或更多的定理，其解释是彼此不相容的，那么它是内部不一致的；而若所有经解释后的定理彼此相容，则是内部一致的。作为一个例子，让我们考虑一个形式系统。它只有下述三个定理：WdZ、ZdG、GdW。如果W解释成``乌龟''，Z解释成``芝诺''，G解释成``吉世达''，而xdy解释成``x总是在下棋时打败y''，那么我们就有了下列经解释后的定理：

乌龟总是在下棋时打败芝诺。

芝诺总是在下棋时打败吉世达。

吉世达总是在下棋时打败乌龟。

陈述本身并非不相容，虽然它们描述了一个相当怪诞的棋手圈子。所以，在这个解释下，以那三个符号串为定理的形式系统是内部一致的，虽然实际看来那三个陈述没一个是真的！内部一致性不要求所有的定理都得成为真陈述，只需它们能成为彼此相容的陈述。

现在假定我们把xdy另外解释成``x是y笔下的人物''，那么我们就有：

乌龟是芝诺笔下的人物。

芝诺是吉世达笔下的人物。

吉世达是乌龟笔下的人物。

在这种情况下，单个的陈述是真是假没有关系------恐怕也没有什么办法能知道哪些是真哪些是假。不过，肯定不会是它们三个同时都真。这样，这个解释就使系统内部地不一致了。这个内部的不一致性不取决于对那三个大写字母的解释，而仅仅取决于对d的解释，同时还取决于那三个大写字母是循环地排列于d的各次出现周围这一事实。于是，人们可以不必解释形式系统的所有符号而找到内部不一致性（在这个例子里，单解释一个符号就够了）。到了足够多的符号都有了解释的时候，可能就会清楚地看出：没有什么解释其余符号的办法能使所有的定理成为真的。但这不仅是真假的问题------这是个可能性的问题。如果我们把大写字母都解释成真人的名字，三个定理就都成了假的------但那不是把系统称作内部不一致的理由。我们的根据是循环性，再加上对字母d的解释。（顺便提醒一句，你会在第二十章再见到这个``作者三角形''）。

假想的世界和一致性

我们给出了两种看待一致性的方式：前一个说，如果每个定理经解释后成为真的，则系统加上解释是与外部世界一致的；后一个说，如果所有的定理经解释后成为彼此相容的，则系统加上解释是内部一致的。这两种类型的一致性之间有密切的关系。为了决定几个陈述是否彼此相容，你得设法想象一个世界，在其中它们能同时都真。所以，内部的一致性有赖于与外部世界的一致性------只是现在，``外部世界''允许是任何想象的世界，而不必是我们生活于其中的那个世界。但这是一个极其模糊的、无法令人满意的结论。一个``想象''的世界都有些什么东西？是否可能想象一个世界，在其中三个人物是循环地相互创造出来的？是否可能想象一个世界，其中有方的圆？这样一个世界是不是可想象的：其中是牛顿定律而非相对论成立？是否可能想象一个世界，其中有绿的并且又不是绿的东西？或者这样一些世界：其中的动物不是细胞组成的、在其中巴赫在腓特烈大帝的主题上即兴演奏了八声部赋格曲、在其中蚊子比人更聪明、在其中乌龟能踢足球------或谈论足球？当然了，一只谈论足球的乌龟的确有点反常。

这些世界里，有一些似乎比另一些更可想象，因为有些世界似乎包含了逻辑矛盾------例如，绿的并且又不是绿的。而另一些似乎是------找不到更好的词------``说得通的''，例如巴赫即兴演奏八声部赋格曲、存在不是由细胞构成的动物。甚至可以设想一个有不同物理定律的世界\ldots\ldots 因此，大体上该是有可能建立不同品种的一致性。举个例子来说，最宽松的会是``逻辑一致性''：对事物不加任何制约，只对逻辑有所要求。更具体地说就是一个系统加上解释是``逻辑一致''的，仅当它其中没有任何两个定理在解释成陈述时直接地相互矛盾；``数学一致''的，仅当经解释后的定理不违背数学；``物理学一致''的，仅当所有经解释后的定理与物理法则相容；``生物学一致''的\ldots\ldots 如此等等。在一个生物学一致的系统里，可以有这样的定理，其解释是陈述``莎士比亚写了一部歌剧''，但不会有定理其解释是陈述``存在无细胞动物''。一般说来，这类玄想出来的不一致性没什么人去研究，原因是很难把它们一个个理清楚。举例来说，三个人物循环地互相创造这个问题涉及的不一致性应该说是哪一种？逻辑的？物理学的？生物学的？文学的？

通常，有意思和没有意思之间的分界线是画在数学一致性和物理学一致性之间的。（这当然是数学家和逻辑学家们画的了------很难说他们是一个不带偏见的群体\ldots\ldots）这意味着对于形式系统来说，只有数学和逻辑类型的不一致性才``算数''。若按照这个约定，那么我们就还没有找出一个解释，使那个三定理组WdZ、ZdG、GdW不一致。这一点可以通过把d解释成``大于''来做到。

W、Z和G怎么办？它们可以解释成自然数------比如，Z是2，W是3，G是6。注意这种解释下有两个定理成为真的，有一个假的。如果我们另外把Z解释成了8，就会有两个假的，仅有一个真的了。但两者都使我们得到了不一致性。事实上，指派给W、Z和G的值是无关紧要的，只要明白需限制在自然数范围内。我们又一次看到这种情形：为了认出内部不一致性，只需要某些解释。

形式系统中嵌入形式系统

前面那个例子，其中某些符号可以在其余符号没有解释的情况下具有解释，让人想起在自然语言里构造几何学时，使用的某些词是未定义项。在这种情况下，词就可以分成两类：一类词有固定不变的意义，另一类词的意义有待调整，直到系统成为一致的（这一类是未定义项）。用这种方式构造几何学，要求第一类词的意义已经在几何学之外制定好。这些词构成一个刚性的骨架，赋予系统一个基础结构；填充那个骨架的是其它材料，它们是可以改变的（欧氏或非欧几何学）。

形式系统常常就是以这种相继的、或者说分层的方式构造出来的。举例来说，形式系统Ⅰ设计好了，有一些规则和公理，它们赋予符号以某种意向的被动意义。然后形式系统Ⅰ完全合并到一个有更多符号的更大的系统里------形式系统Ⅱ。由于形式系统Ⅰ的公理和规则是形式系统Ⅱ的一部分，形式系统I的符号的被动意义仍然有效。它们构成一个不变的骨架，随后便在决定形式系统Ⅱ的新符号的被动意义时起很大的作用。然后，或许会轮到第二个系统对第三个什么系统起骨架的作用，如此下去。也可能------几何学就是个很好的例子------有一个系统（例如绝对几何学），它部分地固定住其中的未定义项的被动意义，并能对其添加另外的规则或公理，以进一步地限定那些未定义项的被动意义。这就是欧氏与非欧几何学的情形。

视知觉中的稳定性层次

我们获得新知识、新词汇，或感知不熟悉的事物，都是以与此类似的分层方式进行的。尤其有趣的是理解艾舍尔的画时的情形。比如，《相对性》（图22）中出现了简直是不可能的景象。你也许认为我们会一遍又一遍地力求重新解释这幅画，直到我们把它的各个部分解释成没有矛盾的------但我们根本不那么做。我们坐在那里，看着那些通向各处的楼梯和那些在同一楼梯上沿着不一致的方向走动的人们，又愉快又困惑。那些楼梯是些``确定性的岛屿''，我们对整幅画的解释就以它们为基础。一旦确定了它们，我们就试图扩展我们的理解，去探求它们彼此之间的关系。就是在这一步骤我们遇上了麻烦。但如果我们想要回溯------就是说，对``确定性的岛屿''质疑------我们也会遇上麻烦，只不过是另一种麻烦而已。没有什么通过回溯否认它们是楼梯的办法。它们不是鱼、鞭子或手------它们就是楼梯（实际上有另外一个办法------让画中的所有线条都完全没有解释，就像形式系统中的``没有意义的符号''。这种彻底的逃避途径是``U方式''反应的一个例子------一种对待符号表示的禅宗式态度）。

\pandocbounded{\includegraphics[keepaspectratio]{./media/OEBPS/Images/img_022.jpg}}

图22．相对性

艾舍尔作（版画，1953）

由于我们知觉过程的层级性质，我们被迫或者看到一个疯狂的世界，或者看到一簇毫无意思的线条。相似的分析可用于艾舍尔的许多画，它们都大量地依赖于对某些基本形状的识别，然后再以非标准的方式组合在一起。到观众在高层上看出悖谬的时候，就已经太晚了------他无法再回去，对怎样解释较低层次的对象改换想法。艾舍尔的画与非欧几何学的区别在于：对于后者，能找到对未定义项的可理解的解释，得出一个可理解的完整系统；而对于前者，无论盯着画看多长时间，最后的结果与人们关于世界的概念仍不可调和。当然，人们仍然可以虚构假想的世界，在其中艾舍尔的那类事件能够发生\ldots\ldots 但是在那样的世界里，生物学、物理学、数学甚至逻辑的法则会在一个层次上被违反而同时在另一个层次上被遵循，使之成为一个极其古怪的世界（例子之一是《瀑布》（图5），那里面正常的引力适用于流动的水，但空间的性质却违反了物理学原理）。

数学在每个可想象的世界里都是一样的吗？

上面我们强调过，形式系统（加上一个解释）的内部一致性要求存在某种想象的世界------对这类世界唯一的限制是其中的数学和逻辑与我们的世界一样------在其中所有经解释后的定理都成为真的；而外部的一致性------与外部世界的一致性------则要求所有的定理在现实世界里都成为真的。现在来看一个具体情形。一个人希望构造一个一致的形式系统，其定理要解释成数学陈述，那么看起来两种类型的一致性之间的区别就消失了，因为，如上所说，所有想象的世界都有与现实世界一样的数学。于是，在每个可想象的世界里，1加1必须是2；同样，必须存在无穷多的素数；进一步，在每个可想象的世界里所有直角都必须是相等的；而且当然还有：经过任一不在给定直线上的点，必是恰有一条与之平行的直线\ldots\ldots{}

但等一下！那可就是平行公设------断言其普遍性将是个错误，这我们刚刚说过。如果所有可想象的世界都遵循平行公设，那就是断言非欧几何是不可想象的，我们就又回到了萨彻利和兰伯特的心理水平了------肯定是不智之举。但什么车西，如果不是所有数学，必须是所有可想象的世界共有的？会是少到逻辑本身吗？或甚至逻辑也是可疑的：是否有那种世界：矛盾是存在的正常组成部分------矛盾在那里不是矛盾？

从某种意义上讲，仅仅作为发明出来的概念，我们已表明这种世界是可想象的；但在更深一层的意义上讲，它们是很不可想象的（这本身又是个小矛盾）。然而，认真地讲，如果我们想要有所交流，看来就得采纳某种共同基础，而逻辑几乎总是得包括进来的。（有些信仰系统是拒斥这种观点的------它太逻辑化了。具体说来，禅宗就是以同等的热情拥抱矛盾和非矛盾的。这似乎不太一致，但不一致也是禅宗的一部分，所以\ldots\ldots 你能说什么呢？）

数论在每个可想象的世界里都是一样的吗？

如果我们假定逻辑是每个可想象世界的一部分（注意，我们还没有给逻辑下定义，但我们会在以后的章节里给出来），是否就行了？真的能想象某些世界没有无穷多素数吗？所有可想象的世界里，数应该遵循同样的法则，这难道不是必要的吗？或者\ldots\ldots 把``自然数''这个概念也当作未定义项，就像``\heititext{点}''或``\heititext{线}''，是不是更好些？那样的话，数论就是个分叉的理论了，像几何学一样：就会有标准数论和非标准数论了。但要有与绝对几何学相应的某种东西：一个``核心''理论，一个所有数论的不变组成部分，以确认数论是数论，而不是关于------比如说，可可或橡皮或香蕉的理论。存在这样的一个核心数论，它应该与逻辑一起包含在我们认为是``可想象的世界''里，这似乎是大多数当代数学家和哲学家的共识。这个数论的核心------绝对几何学的对应物------被称作皮亚诺算术，我们会在第八章里形式化地给出来。另外，现在已经很明朗------实际上是哥德尔定理的直接后效------数论是分叉的理论，有标准的和非标准的版本。然而，与几何学的情形不同，数论的``品种''有无限多个。这使得数论的情形变得极其复杂。

从实际应用的角度看，所有的数论都是一样的。换句话说，如果桥梁的建造依赖于数论（在某种意义上确实如此），那么，不同的数论的存在不会造成什么影响，因为在与现实世界有关的那些方面，所有的数论都是重合的。对不同的几何学可不能这么说。例如，三角形的内角和只在欧几里德几何中是180度，在椭圆几何中要大于这个值，而在双曲几何中要小于这个值。据说高斯曾试图测量三座山峰所形成的三角形的内角和，以最终确定我们的宇宙到底服从哪种几何学。直到一百年后，爱因斯坦提出了一种理论（广义相对论），认为宇宙的几何性质是被其内含的质量所确定的，因此没有哪种几何是空间自身所固有的。所以，对``哪种几何学是正确的？''这个问题，自然界不仅在数学中，而且也在物理学中给出了一个模棱两可的答案。至于那个与此对应的问题``哪种数论是正确的？''，当我们详细地介绍了\heititext{哥德尔定理}后，还会有更多的话要说。

完全性

如果一致性是符号获得被动意义的最低条件，那么与之互补的概念，完全性，是那些被动意义的最高确认。一致性是说：``系统产生的每个东西都是真的''，完全性是倒过来：``每个真陈述都是由系统产生的''。现在稍稍修饰一下这个概念。我们不会是指世界上所有的真陈述------我们指的仅是这样的陈述：它们所属的领域是我们力图用该形式系统去表达的。所以，完全性的意思是：``每个能由系统中的概念表示出来的真陈述都是系统中的定理''。

{一致性：每个定理经解释后都成为真的（在某个想象的世界里）。}

{完全性：所有真的（在某个想象的世界里）且可表示成系统中的良构符号串的陈述都是定理。}

最初那个pq系统以及它原来那个解释是这种形式系统的一个例子：它在自己那谦逊的水平上是完全的。任何两个正整数的正确的加法式子都可表示为系统中的定理。我们也可以换一种方式来说：``任何两个正整数的加法式子都是在系统内可证的''。{（小心：当我们开始用``可证的陈述''而非``定理''这种说法时，表明我们开始混淆形式系统与其解释之间的区别。这是可以的，只要我们清楚地意识到这种混淆出现了，并且记得多重解释有时是可能的。）}pq系统加上它原来那个解释是完全的；同样也是一致的，因为没有假陈述------用我们的新说法------是在系统内可证的。

有人可能会争辩说这个系统是不完全的，根据是三个正整数的加法式子（如2+3+4=9）没有被pq系统的定理所体现，尽管能翻译成系统中的概念（例如，-\/-\/-\/-\/-\/-\/-\/-\/-q-\/-p-\/-\/-p-\/-\/-\/-）。然而，这个符号串不是良构的，因此我们认为它就像pqp-\/-\/-qpq一样缺乏意义。三个数的加式用系统中的概念根本不可表示------所以系统还是有完全性的。

尽管pq系统在这种解释下有完全性，离完全地捕捉数论中的真理当然还差得远。比如，pq系统就没有办法告诉我们有多少素数存在。\heititext{哥德尔不完全性定理}说的是任何``足够强有力''的系统，由于其能力较强，因而是不完全的。即：存在良构的符号串，表示了数论中的真陈述，但不是定理（有属于数论的真陈述在系统内不可证）。pq系统这样的系统是完全的，但能力不强。它们更像是低保真的唱机：太贫乏了，很明显做不了我们想要它们做的------即告诉我们数论中的每件事实。

一个解释怎样就能达到或破坏完全性？

那么所谓``完全性是被动意义的最高确认''------我在前面这么说过------是什么意思呢？意思是：如果一个系统是一致的，但不完全，符号和其解释之间就会错配。系统没有能力为那样的解释作辩护。有时，如果把解释稍稍``修剪''一下，系统就会变得完全。为说明这个思想，让我们来看看修改的pq系统（含有公理模式Ⅱ）及我们对它的解释。

修改了pq系统之后，我们把q的解释从``等于''修改为``小于或等于''。我们看到，修改的pq系统在这种解释下是一致的，但新解释却有些令人不太满意的地方。问题很简单：现在有许多可表示的真陈述不是定理。例如，``1小于或等于2加3''表示成非定理-q-\/-p-\/-\/-。这个解释实在是太不彻底了！它没有精确地反映出定理在系统中的作为。我们可以这样来纠正：或者（1）向系统添加新规则，使之增强能力，或者（2）紧缩解释。在这里，明智的选择看来是紧缩解释。不把q解释成``小于或等于''，而代之以``等于或加1后等于''。现在，修改后的pq系统变得既是一致的又是完全的，而完全性确保了这个解释是恰当的。

形式化数论的不完全性

在数论里我们会再次碰到不完全性。但在那里，为了挽救局面，我们被推向另一方向------增加新规则，使之增强能力。具有讽刺意味的是，每次我们增加一条新规则，我们就想，现在我们肯定使系统完全了！这种困境的性质可由下述类比来说明\ldots\ldots{}

我们有一台唱机，还有一张唱片，上面标着``基于B-A-C-H的卡农''。然而，当我们把唱片放在唱机上播放时，由反馈引起的震颤（与乌龟的唱片引起的震颤一样）干扰得太厉害，以致于我们甚至都辨不清曲调了。我们断定某个东西有缺陷------不是唱片就是唱机。为了检测我们的唱片，我们就得到朋友的唱机上去播放，听它的质量。为了检测我们的唱机，我们就得用朋友的唱片在其上播放，看看我们听到的音乐是否符合标签。如果我们的唱机通过了对它的检测，则我们就说是唱片有缺陷；反过来一样，如果唱片通过了对它的测试，则我们就说是唱机有缺陷。然而，若结果是两者分别通过了对它们的测试，我们作什么结论呢？现在正是时候：回忆一下那两个相继的同构（图20），仔细想想！

\begin{enumerate}
\item
  \phantomsection\label{Chapter04.xhtml_B_1}
\item
  \phantomsection\label{Chapter04.xhtml_B_2}
\end{enumerate}

\phantomsection\label{Dialog05.xhtml}{}

\chapter{和声小迷宫}

乌龟和阿基里斯在游乐场玩了一天。他们先买了一对棉花糖，然后决定去乘``大风车''玩。

\dialogname{乌龟：}这是我最喜欢的游艺项目。坐在上面你感觉似乎移动了很远，其实仍在原地。

\dialogname{阿基里斯：}我知道它为什么叫你着迷。你系好安全带了吗？

\dialogname{乌龟：}好啦，我觉得我把这扣子扣好了。哎，开始吧。喔！

\dialogname{阿基里斯：}你今天很兴奋吧。

\dialogname{乌龟：}这是有缘故的。我姑妈会算命，她对我说今天我会遇到好运。正是因为期待着那个好运我才这么兴奋的。

\dialogname{阿基里斯：}你居然相信算命！

\dialogname{乌龟：}不\ldots\ldots 但是人们说即使你不信，它也一样灵。

\dialogname{阿基里斯：}哼，要真是这样倒不错。

\dialogname{乌龟：}啊，风景多美啊。你看这片海滩、海边上游泳的人们、蓝色的大海、还有这座城市\ldots\ldots{}

\dialogname{阿基里斯：}对，的确美不胜收。喂，瞧那边那架直升飞机，似乎朝我们这儿飞来了。现在它实际上几乎就在我们头顶上了。

\dialogname{乌龟：}怪事------飞机上还垂下一条缆绳来，离我们非常近。我们只要伸手就能抓得到。

\dialogname{阿基里斯：}瞧！那条缆绳的头儿上有个大钩子，上面挂着一张纸条。

（他伸手抓下纸条。旋转着的大风车带着他们转了下去。）

\dialogname{乌龟：}你能认出条上写的什么吗？

\dialogname{阿基里斯：}能------写的是，``你们好，朋友们。再转上来的时候抓住钩子，会叫你们喜出望外的。''

\dialogname{乌龟：}条子上写的跟游艺场广告上的陈词滥调一个味儿，不过谁知道这意味着什么。也许跟我要遇到的好运有关。不管怎么着，我们得试试！

\dialogname{阿基里斯：}得试试！

（``大风车''把他们转上来的时候，他们两个都解开了安全带的系扣。当转到最高点时，他们抓住了那个巨大的钩子，缆绳马上把他们悠起来，一直绞向悬在空中的直升飞机。一只粗大有力的手帮他们进了飞机。）

\dialogname{声音：}欢迎登机------两位冒失鬼。

\dialogname{阿基里斯：}你------你是谁？

\dialogname{声音：}请允许我自我介绍一下：我是郝晕，逍遥自在的拐子，最优秀的噬龟者，听候二位的吩咐。

\dialogname{乌龟：}哎呀！

\dialogname{阿基里斯}唔------哦------我看这位``郝晕''真不是我们期待的那个。\kaitismall{（对郝暈说）}哎------如果我可以如此大胆的话------你这是要把我们拐到哪儿去啊？

\dialogname{郝晕：}哈哈！去我的全电气化空中厨房，在那儿我要把这位做成有滋有味的小点心------\kaitismall{（他说这话的时候，瞟了乌龟一眼）}------做成美味可口的天上馅饼！没错儿------拿来满足我的饕餮之欲！哈哈哈！

\dialogname{阿基里斯：}我只能说你的笑声像恶魔一样。

\dialogname{郝晕}哈哈哈！你竟敢这么放肆！朋友，有你好瞧的。哈哈！

\dialogname{阿基里斯：}天哪！他这话是什么意思！

\dialogname{郝晕：}很简单------我给你们二位准备了恶运！你们等着吧！哈哈哈！哈哈哈！

\dialogname{阿基里斯：}糟啦！

\dialogname{郝晕：}好，我们到了。下机吧，朋友们，到我那绝妙的全电气化空中厨房里来吧。

（他们走进去。）

在你们倒霉之前，让我带着你们先参观参观。这是我的卧室，这是我的书房。请先在这儿等我一会儿，我得去磨磨我的刀子。你们可以边等边吃点``弹出锅酥''。哈哈哈！乌龟馅！乌龟馅！我最喜欢的一种馅！\kaitismall{（走开了。）}

\dialogname{阿基里斯：}喔，天哪------弹出锅酥！我可得甩开腮帮子足吃！

\dialogname{乌龟：}阿基！你可是刚刚大嚼了一通棉花糖啊！况且，这种时候你怎么还想吃东西？

\dialogname{阿基里斯：}糟糕------哦，真对不起------我不应该用带有``糕''这种食品名称的词，我是想说在这样不幸的时候\ldots\ldots{}

\dialogname{乌龟：}我想恐怕我们都得引颈就戮了。

\dialogname{阿基里斯：}是啊------你是应该翘首引颈好好瞧瞧郝晕那家伙书房里的那些书。那是些什么书啊------哦，原来是孤本书大全：《我所认识的呆瓜》、《象棋和转伞安然相处》、《为踢踏舞和管弦乐队写的协奏曲》\ldots\ldots 哼！

\dialogname{乌龟：}书桌上放着一本什么书啊？紧挨着那个十二面体和打开的速写本？

\dialogname{阿基里斯：}这本吗？喔，书名是《阿基里斯和乌龟在全球瞎逛时引人入胜的历险》。

\dialogname{乌龟：}名字有点吸引人呢。

\dialogname{阿基里斯：}的确------而且这页上的故事看上去也很吸引人。故事的名字是《神怪和煮调饮》。

\dialogname{乌龟：}冷饮？

\dialogname{阿基里斯：}不是，是``煮调饮''，大概是一种煮着喝加调料的饮料吧？我也不懂。

\dialogname{乌龟：}哼\ldots\ldots 怪事。我们读读它好吗？我读乌龟的台词，你当阿基里斯。

\dialogname{阿基里斯：}行啊。干嘛不呢\ldots\ldots{}

（他们开始阅读《神怪和煮调饮》。）

\begin{adjustwidth}{2em}{0em}

（阿基里斯邀请乌龟来欣赏他收集的他最喜欢的画家艾舍尔的作品。）

\dialogname{乌龟：}这些画好极了，阿基。

\dialogname{阿基里斯：}我知道你会喜欢的。你最喜欢哪一幅？

\dialogname{乌龟：}我最喜欢的一幅是《凹与凸》，在那幅画中有两个在各自的内部完全统一的世界，然而把它们合在一起，却形成一个完全不统一的世界。不统一的世界里总是很好玩的，不过，我不想住在那儿。

\dialogname{阿基里斯：}你说什么？去哪儿玩？不统一的世界根本不存在，你怎么能去玩？

\dialogname{乌龟：}哎，可我们不是都认为艾舍尔的这幅画画出了一个不统一的世界吗？

\dialogname{阿基里斯：}对，可这只是个二维世界------是个虚构的世界------一幅画而已。你无法去这样的世界里玩。

\dialogname{乌龟：}我自有办法\ldots\ldots{}

\dialogname{阿基里斯：}你怎么能把自己弄进一幅平面上的画里去？

\dialogname{乌龟：}喝上一小杯``推入露''，这个戏法就变成了。

\dialogname{阿基里斯：}推入露？这是种什么鬼东西？

\dialogname{乌龟：}是种装在陶瓷小瓶里的液体，一个人若是在看着一幅画时喝它，它就能把那个人``推入''到那幅画的世界里。那些不了解推入露魔力的人经常会因为他们周围环境的突变而惊讶不已。

\dialogname{阿基里斯：}没有解药吗？一旦被推入就无可挽回了吗？

\dialogname{乌龟：}在某种情况下事情没有那么糟。事实上还有另外一种露------嗯，实际不是一种露，而是一种制剂------不，不是什么制剂，是一种、一种------

\end{adjustwidth}

\dialogname{乌龟：}他可能是想说``煮调饮''。

\begin{adjustwidth}{2em}{0em}

\dialogname{阿基里斯：}煮调饮？

\dialogname{乌龟：}正是我要找的那个词！它就是被称作``弹出煮调饮''的那种东西，如果你喝推入露时，想着在你右手里拿上这么一瓶弹出煮调饮，那它也可以被推入图画。这样一来，无论你何时想``弹出''，回到现实生活中来，你只需喝一口弹出煮调饮，接着，疾！就回到现实世界了，恰好回到你把自己推入的那个地方。

\dialogname{阿基里斯：}听起来挺有趣。要是你在此之前没把你自己推入一幅画里，却喝了弹出煮调饮，会怎么样呢？

\dialogname{乌龟：}这我不太清楚，阿基。不过摆弄这些奇异的推入弹出药水一定要小心。我曾经有个亲戚，一只鳖，就是因为像你所说的那样做了，从那儿以后再也没有人见过它。

\dialogname{阿基里斯：}真不幸，你能把一瓶推入露也带在身上吗？

\dialogname{乌龟：}哦，当然啦。左手拿着它，它就会跟你一起被推入你在看的画里。

\dialogname{阿基里斯：}要是你在你已经进入的那幅画里再发现一幅画，那会怎么样呢？再喝一大口推入露吗？

\dialogname{乌龟：}正像你所设想的那样：你会进入画中的画。

\dialogname{阿基里斯：}我想你得弹出两次，这样你才能从你所在的两幅叠套的画中解脱出来，重新浮现在现实世界。

\dialogname{乌龟：}对，你推入一次也得弹出一次，因为一次推入使你进入一幅画，而一次弹出则会解除它。

\dialogname{阿基里斯：}我觉得所有这一切都不太对劲儿。你不是在测验我多容易上当受骗吧？

\dialogname{乌龟：}我发誓不是！瞧------这是两只药水瓶，就在我兜里。

（把手伸进翻领口袋，拿出两只挺大的无标签的药水瓶，晃动它们时，可以听出一种药水是红的，另一种是蓝的。）

如果你愿意，我们可以试试，怎么样？

\dialogname{阿基里斯：}嗯\ldots\ldots 我认为\ldots\ldots 啊哼\ldots\ldots 也许\ldots\ldots 啊哼\ldots\ldots{}

\dialogname{乌龟：}我就知道你想试试。我们把我们自己推入艾舍尔的《凹与凸》的世界里去好吗？

\dialogname{阿基里斯：}嗯，啊\ldots\ldots{}

\dialogname{乌龟：}就这么定了。我可别忘了拿着这瓶煮调饮一起进去，不然我们就没办法弹出来了。你愿意肩负这一重任吗，阿基？

\dialogname{阿基里斯：}如果不太麻烦你的话，我，我有点紧张，（啊哼）我更愿意让你，你------你有经验------你来负责这件事吧。

\dialogname{乌龟：}那好吧。

（乌龟一边说着，一边倒出两份推入露。然后他拿起那瓶煮调饮，右手紧紧的地抓住它，他和阿基里斯都把杯子放到唇边。）

\dialogname{乌龟：}干杯！

（他们一饮而尽。）

\begin{adjustwidth}{2em}{0em}

\dialogname{阿基里斯：}味道十分特别。

\dialogname{乌龟：}会习惯的。

\dialogname{阿基里斯：}煮调饮的滋味也这么特别吗？

\dialogname{乌龟：}哦，感觉很不一样，无论你什么时候尝那种煮调饮，你都会有一种深深的满足感，就好像你整个一生都在等着尝它似的。

\dialogname{阿基里斯：}哦，我真想喝到它。

\dialogname{乌龟：}哎，阿基，我们在哪儿呢？

\dialogname{阿基里斯}我们是在一条小平底船上，沿一条运河顺流而下！我想下去，船夫先生，请让我们下去吧。

（船夫对这个请求毫不理会。）

\dialogname{乌龟：}他不懂汉语。如果我们想离开这儿，我们最好还是在他进入前面那条罪恶的``爱之隧''以前，赶紧从河道爬出去。

（阿基里斯的脸色有些苍白，霎时就爬到了河岸上，然后把他那位动作缓慢的朋友也拽了出来。）

\dialogname{阿基里斯：}我有点不喜欢那隧洞的名字。我们能逃出来我真高兴。哎，你怎么会对这些地方这么熟悉？你以前来过这儿吗？

\dialogname{乌龟：}来过很多次，虽然我总是从艾舍尔其他的画中进来。你知道，它们在画框背后彼此相连。你一旦进了其中的一幅画，你也可以进入别的里面。

\dialogname{阿基里斯：}真的？！要是我不是在这儿亲眼看到了这些事，我肯定不会相信。\kaitismall{（他们穿过一座小拱门，走了出来。）}哦，瞧这两只逗人喜爱的蜥蜴！

\dialogname{乌龟：}逗人喜爱？它们可不逗人喜爱------只要一想到它们我就不寒而栗！它们是两个邪恶的卫士，负责守护悬挂在那边天花板上那个有魔力的黄铜灯。只要被它们的舌头碰一下，任何生物都会变成肉酱。

\dialogname{阿基里斯：}甜的还是咸的？

\dialogname{乌龟：}都不是，是辣的。

\dialogname{阿基里斯：}太毒辣了！不过如果这盏灯有魔力，我倒真想得到它。

\end{adjustwidth}

\end{adjustwidth}

\pandocbounded{\includegraphics[keepaspectratio]{./media/OEBPS/Images/img_023.jpg}}

图23．凹与凸

艾舍尔作（石板画，1955）

\begin{adjustwidth}{2em}{0em}

\begin{adjustwidth}{2em}{0em}

\dialogname{乌龟：}这很危险，我的朋友。我不愿意冒这个险。

\dialogname{阿基里斯：}我只想试一次。

（他偷偷地靠近那盏灯，满以为不会弄醒睡在附近的那个小伙子。可是突然，他在一块奇特的贝壳形凹状地板上滑倒，跌出来掉进空中，东歪西倒的，他想随便够着什么，还想用一只手抓住那盏灯。由于他这么猛烈地摇晃，那两只蜥蜴``咝咝''叫着迅猛地朝他伸出了舌头，而他却吊在空中无法脱身。）

\dialogname{阿基里斯：}救------命------啊！

（一个女人正走下楼梯，呼喊声引起了她的注意，还惊醒了睡着的小伙子。后者弄清了是怎么回事之后，脸上露出了和善的微笑，向阿基里斯打手势，告诉他一切都会好的。他用一种奇怪的带有喉音的语言向高处窗户上的一对号手喊了句什么，然后，马上有奇异的声音奏响了，而且还彼此相呼应。那个睡意惺忪的小伙子指着蜥蜴，阿基里斯看到这音乐对它们有一种强烈的催眠作用。不一会儿，它们完全丧失了知觉。随后那个救命的小伙子对两个正沿着梯子向上爬的人喊话。他们俩把各自的梯子拿上来，伸到在空中被困的阿基里斯那里，构成一座桥，然后打手势让阿基里斯赶快爬下来。而阿基里斯却先小心地把挂着灯的链子从屋顶上的挂钩中脱出来，取下那盏灯，然后才从梯桥上爬下来，那三个小伙子让他脱了险。阿基里斯搂住他们，十分感激地同他们紧紧拥抱。）

\dialogname{阿基里斯：}嘿，龟兄，我怎么报答他们呢？

\dialogname{乌龟：}我正好知道这三位勇敢的孩子喜欢喝咖啡，而就在这座城里有个地方，咖啡煮得没比了。请他们去喝杯咖啡吧！

\dialogname{阿基里斯：}好极了！

（这样，通过一系列可笑的手势、微笑和个别字词，阿基里斯对这些小伙子表达了他的邀请。随后，他们五个一起走出来，走下一段陡楼梯，来到城里。他们来到一家吸引人的咖啡馆，在外面坐下来，要了五杯咖啡。在他们品呷时，阿基里斯记起了他带着的那盏灯。）

\dialogname{阿基里斯：}我都忘了，龟兄------我弄到了这盏魔灯！但是------它有什么魔力？

\dialogname{乌龟：}哦，很普通------一种怪物。

\dialogname{阿基里斯：}什么？你是想说你擦它时，就会出来一个怪物听你使唤？

\dialogname{乌龟：}没错。你以为呢？天上掉馅饼吗？

\dialogname{阿基里斯：}嗯，这是挺神的！我可以随心所欲，嗯？我总是希望我碰见这种事\ldots\ldots{}

（阿基里斯轻轻地擦着刻在那盏黄铜灯表面上的大字``灯''\ldots\ldots 突然，冒出了一大股烟，从烟的形状中，五个朋友都辨认出那是一个精怪，鬼一般地立在他们上面。）

\dialogname{怪物：}你们好，我的朋友们------非常感谢把我的灯从邪恶的双蜥蜴手中救出来。

（怪物一边说着，一边把那盏灯拿起来，塞进隐藏在他那飘摆于灯外的长袍皱褶中的口袋里。）

作为对你这一英雄行为的感谢，我愿意代表灯向你提供实现你任意三个愿望的机会。

\dialogname{阿基里斯：}真是活见鬼！你不觉得是这样吗，龟兄？

\dialogname{乌龟：}我的确也这么想。来吧，阿基，说第一个愿望。

\dialogname{阿基里斯：}唔！可我应该提什么愿望呢？哦，我知道了！这是我在第一次读《天方夜谭》时想过的（那是本由许多愚蠢的------也是叠套的------故事组成的集子）------我愿我有一百个愿望，而不是三个！真聪明，是吧，龟兄？我敢说你从没有想到过这种把戏，我总是想为什么故事中的那些大脑迟钝的人从没有自己这么试试。

\dialogname{乌龟：}也许现在你会找到答案的。

\dialogname{怪物：}很抱歉，阿基里斯，我不能满足元愿望。

\dialogname{阿基里斯：}我愿你告诉我什么是``元愿望''！

\dialogname{怪物：}可这是元元愿望，阿基里斯------我也不能满足你。

\dialogname{阿基里斯：}什什什么？我一点也听不懂。

\dialogname{乌龟：}你干嘛不把你最后那个请求换种说法，阿基？

\dialogname{阿基里斯：}你想说什么？我干嘛要换？

\dialogname{乌龟：}嗯，你刚才说``我愿\ldots\ldots''，你既然只是在问一件事，干嘛不用问句？

\dialogname{阿基里斯：}好吧，虽然我不明白为什么。告诉我，怪物------什么是元愿望？

\dialogname{怪物：}就是关于愿望的愿望。我不能满足元愿望，我只能满足简单的普通愿望，比如想要十箱啤酒啦、月里嫦娥啦、两吨黄金啦、环球旅行啦。你明白吗------类似这样的简单事物。而我无法满足元愿望。\heititext{造物神}不允许我这么做。

\dialogname{阿基里斯：}\heititext{造物神}？谁是\heititext{造物神}？他为什么不让你满足元愿望？这看起来比你刚才提到的那些事要容易得多嘛。

\dialogname{怪物：}嗯，这事挺复杂的，你为什么不来讲讲你那三个愿望呢？或者至少说一个，我不能总呆在这个世界里\ldots\ldots{}

\dialogname{阿基里斯：}哦，真糟糕。我真的希望有一百个愿望\ldots\ldots{}

\dialogname{怪物：}唉，我不愿看到任何人这么失望。而且，元愿望也是我最喜欢的一种愿望。让我想想我是否可以做点什么。请等一分钟。

\hfill\break

（怪物从他袍子上的一簇皱褶里拿出一样东西，那东西看上去就像他拿走的那盏黄铜灯一样，只是这盏是用白银做的。以前的那盏上刻有``灯''字，而在这盏灯上同样的地方却刻有小一些的字``元灯''。）

\dialogname{阿基里斯：}那是什么？

\dialogname{怪物：}是我的元灯\ldots\ldots{}

（他擦擦元灯，它便冒出一股烟来。在翻滚着的烟云中间，他们全都看到了他们上面立着一个魔鬼般的形体。）

\end{adjustwidth}

\begin{adjustwidth}{2em}{0em}

\begin{adjustwidth}{2em}{0em}

\dialogname{元怪物：}我是元怪物。你叫我吗，怪物阁下？你有什么愿望？

\end{adjustwidth}

\dialogname{怪物：}我对你，元怪物阁下，以及\heititext{造物神}，有个特殊的愿望。我想请求暂时解除一切对愿望类型的限制，允许一个无类型愿望存在一段时间。你能俯准我这个愿望吗？

\begin{adjustwidth}{2em}{0em}

\dialogname{元怪物：}我会把它通过层层机构发送出去的，这没有问题。请等半分钟。

（这位元怪物以怪物两倍的速度从她袍子上的一簇襞褶里拿出一样东西，那东西看去就像那盏白银元灯一样，只是这盏是用黄金做的。以前的那盏上刻有``元灯''字样，而在这盏灯上同样的地方却刻有更小一些的字``元元灯''。）

\end{adjustwidth}

\begin{adjustwidth}{2em}{0em}

\dialogname{阿基里斯：}\kaitismall{（他的声音比以前高了八度）}：那是什么？

\dialogname{元怪物：}这是我的元元灯\ldots\ldots{}

（她擦擦元元灯，它便冒出一股烟来。在翻滚着的烟云中间，他们全都看到了他们上面立着一个魔鬼般的形体。）

\begin{adjustwidth}{2em}{0em}

\dialogname{元元怪物：}我是元元怪物。你叫我吗，元怪物阁下？你有什么愿望？

\end{adjustwidth}

\dialogname{元怪物：}我对你，元元怪物阁下，以及\heititext{造物神}，有个特殊的愿望。我想请求暂时解除一切对愿望类型的限制，请求允许一个无类型愿望存在一段时间。你能俯准我这个愿望吗？

\begin{adjustwidth}{2em}{0em}

\dialogname{元元怪物：}我会把它通过层层机构发送出去的，这没有问题。请等十五秒钟。

（这位元元怪物以元怪物两倍的速度从他袍子上的一簇皱褶里拿出一样东西，那东西看去就像那盏黄金元元灯一样，只是这盏是用\ldots\ldots）

\begin{adjustwidth}{2em}{0em}

\begin{longtable}[]{@{}lllll@{}}
\toprule\noalign{}
\endhead
\bottomrule\noalign{}
\endlastfoot
• & & & & \\
　 & • & & & \\
　 & 　 & • & & \\
　 & 　 & 　 & • & \\
　 & 　 & 　 & 　 & •\heitismall{\{造物神\}} \\
　 & 　 & 　 & • & \\
　 & 　 & • & & \\
　 & • & & & \\
• & & & & \\
\end{longtable}

\end{adjustwidth}

（\ldots\ldots 旋回到元元元灯中来，元元怪物把它卷回他自己的袍子里，这一动作只有元元元怪物的一半那么快。）

你的愿望被批准了，元怪物阁下。

\end{adjustwidth}

\dialogname{元怪物：}谢谢你，元元怪物阁下，以及\heititext{造物神}。

（这位元元怪物像所有比他更高的那些神怪一样，旋回到元元灯中去，元怪物把它卷回她自己的袍子里，这一动作只有元元怪物的一半那么快。）

你的愿望被批准了，怪物阁下。

\end{adjustwidth}

\dialogname{怪物：}谢谢你，元怪物阁下，以及\heititext{造物神。}

（这位元怪物像所有比她更高的那些神怪一样，旋回到元灯中去，怪物把它卷回他自己的袍子里，这一动作只有元怪物的一半那么快。）

你的愿望被批准了，阿基里斯。

（从他说``请等一分钟''，到这时正好过了一分钟。）

\dialogname{阿基里斯：}谢谢你，怪物阁下，以及\heititext{造物神}。

\dialogname{怪物：}我很高兴通知你，阿基里斯，你可以有一个无类型愿望------这就是说一个愿望、一个元愿望、或一个元元愿望，你愿意有多少``元''都可以------甚至可以有无数多个（如果你愿意）。

\dialogname{阿基里斯：}啊，非常感谢，怪物阁下，不过，我的好奇心给逗起来了。在我说出我的愿望前，你能不能告诉我谁------或者说什么------是``\heititext{造物神}''？

\dialogname{怪物：}行啊。``\heititext{造物神}''是代表``\heititext{造物神}-\heititext{物}色的-\heititext{神}怪''这个短语中几个词的词首字组合。``神怪''这个词一般指怪物、元怪物、元元怪物等等。这是个无类型词。

\dialogname{阿基里斯：}可是------可是------这个由几个词的词首字组合而成的``\heititext{造物神}''怎么会是一个词呢？这不会有任何意思！

\dialogname{怪物：}噢，你没见过递归的词首字组合吗？我以为人人都知道呢。你瞧，``\heititext{造物神}''代表``\heititext{造物神}-\heititext{物}色的-\heititext{神}怪''------而它又可以被展成``\heititext{造物神}-\heititext{物}色的-\heititext{神}怪-\heititext{物}色的-\heititext{神}怪''------它还可以依次被展成``\heititext{造物神}-\heititext{物}色的-\heititext{神}怪-\heititext{物}色的-\heititext{神}怪-\heititext{物}色的-\heititext{神}怪''------它还可以依次被进一步展开\ldots\ldots 你可以想展多开就展多开。因为在每种情况下最前面的三个字总是``\heititext{造物神}''，而这最前面的三个字总可以被展开。

\dialogname{阿基里斯：}可我将永远不会展完！

\dialogname{怪物：}当然不会。你永远也不会把``\heititext{造物神}''全部展开。

\dialogname{阿基里斯：}哼\ldots\ldots 真叫人费解。你对元怪物说``我对你，元怪物阁下，以及\heititext{造物神}，有个特殊的愿望''，这是什么意思呢？

\dialogname{怪物：}我是想不仅向元怪物发出请求，而且也向她之上的所有神怪发出请求。递归词首字组合很自然地完成了这一任务。你瞧，当元怪物接受了我的请求时，她就不得不把它向上传给她的\heititext{造物神}。所以她向元元怪物送去一个同样的信息，后者也照样对元元元怪物这么做\ldots\ldots 沿着这条链一直向上，就把信息传给了\heititext{造物神}。

\dialogname{阿基里斯：}我明白了。你的意思是说\heititext{造物神}端坐在神怪之梯的顶端？

\dialogname{怪物：}不，不，不！``顶端''一无所有，因为根本就没有顶端。这就是为什么\heititext{造物神}是个递归的词首字组合。\heititext{造物神}不是某个终极神怪。\heititext{造物神}是个位于任何已知神怪之上的众神怪之塔。

\dialogname{乌龟：}在我看来，每一个神怪都对\heititext{造物神}是什么有彼此不同的看法，这样一来，由于对任何神怪来说，\heititext{造物神}都是位于他或她之上的一组神怪，所以没有两个神怪具有同样的\heititext{造物神}。

\dialogname{怪物：}绝对正确------因为我是一切神怪中最低的，我可怜那些高层的神怪们，它们幻想着自己多少离\heititext{造物神}更近些，多亵渎啊！

\dialogname{阿基里斯：}乖乖，发明\heititext{造物神}可真需要点儿神思。

\dialogname{乌龟：}你真的相信关于\heititext{造物神}的这些玩意儿吗，阿基？

\dialogname{阿基里斯：}那当然啦。难道你是无神论者吗，龟兄？要不你就是个不可知论者。

\dialogname{乌龟：}我不认为我是不可知论者。我也许是个元不可知论者。

\dialogname{阿基里斯：}什什什么？我一点也听不懂。

\dialogname{乌龟：}是这样\ldots\ldots 如果我是个元不可知论者，我就会为我是否是一个不可知论者而困惑------但是我并不清楚我是否是个元不可知论者；因此我一定是个元元不可知论者（我这么猜想）。啊，好啦。告诉我，怪物，神怪们犯过什么错误吗？在这条链上来回传递消息的时候，有过什么歪曲吗？

\dialogname{怪物：}这种事确实有过。这是无类型愿望不被批准的最根本的原因。注意，这种歪曲在这条链上的任何一个个别环节发生的可能性都是无穷小的------但是当你把这无穷多种可能列出来，就能够清楚这样的歪曲一定会在某些地方发生。事实上，奇怪的是，经常会有无穷个歪曲，虽然它们在这条链上分布得很分散。

\dialogname{阿基里斯：}这么说，能有什么无类型愿望被通过，那简直就是个奇迹了？

\dialogname{怪物：}不。大多数歪曲不会造成严重后果，许多歪曲都互相抵消了。不过偶尔------事实上很少见------某种无类型愿望不能满足的原因可以追溯到某一个不幸的神怪所做的歪曲。一旦发生这种事，那个有罪的神怪就得被迫接受责罚，\heititext{造物神}会鞭打他或她的屁股。对于执鞭者来说，这很有趣；对受笞者来说，也没多大害处。看到这一场面你会觉得挺逗的。

\dialogname{阿基里斯：}我真想看看！不过只有在一个无类型愿望不被批准时才会有这种事，对吗？

\dialogname{怪物：}正是。

\dialogname{阿基里斯：}嗯\ldots\ldots 这给了我一个启发，是关于我的愿望的。

\dialogname{乌龟：}噢，真的吗？什么启发？

\dialogname{阿基里斯：}我希望我的愿望不被批准！

（就在这时，发生了一个无法描述的事件------用``事件''这个词行吗？------因此我们将不去描述它了。）

\end{adjustwidth}

\dialogname{阿基里斯：}这个难以捉摸的解释究竟是什么意思？

\dialogname{乌龟：}它是指阿基里斯说出的无类型愿望。

\dialogname{阿基里斯：}可他还没说出来呢。

\dialogname{乌龟：}他说了。他说：``我希望我的愿望不被批准''，而怪物把这当做是他的愿望了。

（就在这时，他们听到他们那一边的厅廊里有脚步声走来。）

\dialogname{阿基里斯：}哦，天呐！听起来可真不祥。

（脚步声停了，随后又转变了方向，并逐渐消失了。）

\dialogname{乌龟}唉！

\dialogname{阿基里斯：}故事说到哪儿了？这就是结尾吗？我们翻过这页看看。

（乌龟把《神怪和煮调饮》这页翻过去，他们发现故事在继续进行着\ldots\ldots）

\begin{adjustwidth}{2em}{0em}

\dialogname{阿基里基：}嗨！怎么回事？我的怪物哪儿去了？我的灯呢？我那杯咖啡呢？我那些从凹与凸的世界里来的朋友们怎么样了？那些小蜥蜴在这儿干什么呀？

\dialogname{乌龟：}恐怕我们的环境恢复错了。

\dialogname{阿基里基：}你这个难以捉摸的解释究竟是什么意思？

\dialogname{乌龟：}我是指你说出的无类型愿望。

\dialogname{阿基里基：}可我还没说出来呢。

\dialogname{乌龟：}你说了。你说：``我希望我的愿望不被批准''，而怪物把这当做是你的愿望了。

\dialogname{阿基里基：}哦，天哪！听起来可真不祥。

\dialogname{乌龟：}这其实是个悖论。因为批准这个无类型愿望，就是否定了它------而不批准它，就是批准了它。

\dialogname{阿基里基：}这会怎么样呢？地球会停止转动吗？宇宙会塌下来吗？

\dialogname{乌龟：}不会，但是系统塌了。

\dialogname{阿基里基：}什么意思？

\dialogname{乌龟：}就是说你我，阿基，会忽然地、同时地被驱逐到堕界去。

\dialogname{阿基里基：}到哪儿去？

\dialogname{乌龟：}堕界：没有打出来的嗝和已熄灭的灯光所在的地方。它是一种候室，在这儿处于休眠的软件等着宿主硬件回来。无法知道这个系统要瘫痪多久，我们会一直呆在堕界里，也许几分钟，也许要几小时、几天------甚至几年。

\dialogname{阿基里基：}我不知道软件是什么，我也不知道硬件是什么。不过我可知道我没有说出我的愿望！我想要怪物回来！

\dialogname{乌龟：}很抱歉，阿基------是你闯的祸。是你搞垮了那个系统。我们能回来真得谢天谢地，不然后事情会变得更糟。不过我一点也不知道我们现在在哪儿。

\dialogname{阿基里基：}我现在认出来了------我们在艾舍尔的另一幅画里。这回是《爬虫》。

\end{adjustwidth}

\end{adjustwidth}

\pandocbounded{\includegraphics[keepaspectratio]{./media/OEBPS/Images/img_024.jpg}}

图24．爬虫

艾舍尔作（石板画，1943）

\begin{adjustwidth}{2em}{0em}

\begin{adjustwidth}{2em}{0em}

\dialogname{乌龟：}啊哈！系统在坍塌之前尽可能地努力保存我们的环境，它仅仅来得及记下我们是在艾舍尔的一幅有蜥蜴的画里，就完蛋了。也真为难它了。

\dialogname{阿基里基：}瞧------那边桌子上不是我们的那瓶弹出煮调饮吗？就是在蜥蜴圈旁边！

\dialogname{乌龟：}没错儿，阿基。我得说我们的确很幸运。系统把弹出煮调饮还给了我们，这对我们真是太好了------它对我们来说真是太珍贵了！

\dialogname{阿基里基：}我也这么说！现在我们可以从艾舍尔的世界里弹回我的房子里了。

\dialogname{乌龟：}煮调饮旁边还有两本书。什么书啊？\kaitismall{（拿起小的那本，随便翻开一页。）}这书看起来挺吸引人的。

\dialogname{阿基里基：}哦，真的？书名是什么？

\dialogname{乌龟：}《乌龟和阿基里斯在全球各地转悠时的历险》。看样子读起来一定很有趣。我得好好看看------怪物和阿基里斯旅行的故事在哪儿？原来这本书有好几个层次，我得仔细找找。

\dialogname{阿基里基：}好吧，你要想看就看吧，但是我可不想让那瓶弹出煮调饮冒险------没准儿哪只蜥蜴会把它碰下桌子，我现在就想把它拿到手！

（他冲到桌子前，伸手去拿那瓶弹出煮调饮，可是忙乱中，他大概是撞了那瓶子一下，瓶子翻倒了，掉下桌子开始滚起来。）

哦，哦！龟兄------你看！我不小心把那瓶子碰到了地板上，它朝------朝------楼梯滚去了！快------别让它掉下去！

（可乌龟这时完全沉浸在他手中的那本小薄书里去了。）

\dialogname{乌龟}嗯嗯，看来在这一层里是找不到那种意思了。

\dialogname{阿基里斯：}龟兄，龟兄，帮帮忙！帮忙抓住那个药瓶子！

\dialogname{乌龟：}干嘛这么大惊小怪的？

\dialogname{阿基里斯：}那只装煮调饮的瓶子------我把它从桌子上碰掉了，现在它正滚呢，它------

（这时它滾到了楼梯井的边缘，然后掉下去了。）

啊，可别！这可怎么办？龟兄------你怎么还跟没事一样？我们的药水没了！掉到楼梯下面去了！只有一个办法了，我们得下一层。

\dialogname{乌龟：}下一层？没问题，请跟我来！

（他开始大声朗读，一时间阿基里斯进退两难，最后终于安定下来，充当了乌龟的角色。）

\begin{adjustwidth}{2em}{0em}

\dialogname{阿基里斯：}这儿真黑，龟兄。我什么都看不见。吁！我撞在墙上了。当心！

\dialogname{乌龟：}来------我这里有两只俄国手杖。拿上一根吧。你可以用它探路，这样你就不会碰东撞西了。

\dialogname{阿基里斯：}好主意。\kaitismall{（他拿了手杖。）}你觉出我们走的这条路有点往左拐吗？

\dialogname{乌龟：}对，有一点点。

\dialogname{阿基里斯：}我不知道咱们到哪儿了。我们还能重见天光吗？我当时真不该听你的，喝下了贴着``喝我''的那瓶东西。

\dialogname{乌龟：}我向你保证这没什么不好。我干过无数次了，一次也没有后悔过。放松点，享受一下变小后的感觉。

\dialogname{阿基里斯：}变小？龟兄，你都对我干了些什么呀？

\dialogname{乌龟：}别怪我，你这样做是你自愿的啊。

\dialogname{阿基里斯：}你把我缩小了吗？那么我们呆在其中的这个迷宫实际上是个小得可以被一脚踏过的玩意啦？

\dialogname{乌龟：}迷宫？迷宫？可能吗？难道我们是在那个著名的由可怕的大雕所占据的和声小迷宫里吗？

\end{adjustwidth}

\end{adjustwidth}

\end{adjustwidth}

\pandocbounded{\includegraphics[keepaspectratio]{./media/OEBPS/Images/img_025.jpg}}

图25．克里特迷宫

意大利蚀刻画；费尼盖拉派{[}选自马修斯{[}W.H.Mathews{]}所著《迷津和迷宫------其历史及发展》{[}Mazes and Labyrinths: Their History and Development{]}纽约：Dover Publications，1970年版{]}

\begin{adjustwidth}{2em}{0em}

\begin{adjustwidth}{2em}{0em}

\begin{adjustwidth}{2em}{0em}

\dialogname{阿基里斯：}天哪！那是什么？

\dialogname{乌龟：}据说------我本人从不相信------一只邪恶的大雕造了一个小迷宫，自己居于它中间的一个小窝里，等待无辜的牺牲品在它那可怕的繁复通道之间转向。然后，当他们茫然失措地转悠到迷宫中心时，它就大笑不止，冲他们大笑------使劲地笑，直到把他们笑死！

\dialogname{阿基里斯：}太可怕了！

\dialogname{乌龟：}不过这只是个神话。勇敢些，阿基。

（这坚强的一对艰难地走着。）

\dialogname{阿基里斯：}摸摸这些墙。他们好像是有褶皱的锡板什么的。不过褶皱的大小不同。

（为了强调他的看法，他一边走，一边把那只俄国手杖伸出来抵着墙壁。当手杖在褶皱上来回弹跳时，在他们所在那条弯曲的长长走廊里，回响了奇特的声音。）

\dialogname{乌龟}这是什么声音？

\dialogname{阿基里斯：}噢，是我在墙上蹭我的手杖呢。

\dialogname{乌龟：}哦！我以为是凶猛的大雕在叫唤呢！

\dialogname{阿基里斯：}当然啦，没有什么可怕的。

（阿基里斯又把他的手杖蹭着墙，继续走着。这时可以听到他的手杖刮墙的地方发出了音乐声。）

\dialogname{乌龟：}嗯，我觉得不妙，阿基。那个迷宫也许根本就不是个神话。

\dialogname{阿基里斯：}等等，是什么让你突然改主意了？

\dialogname{乌龟：}你听到音乐了吗？

（为了听得更清楚些，阿基里斯放下了抬起的手杖，而那美妙的旋律也随之停止了。）

嗨！把手杖放回去！我想听听这段曲子的结尾！

（阿基里斯稀里糊涂地听从了乌龟的指令，音乐又继续下去了。）

谢谢，我刚才要说话的时候，突然明白我们在哪儿。

\dialogname{阿基里斯：}真的？我们在哪儿？

\dialogname{乌龟：}我们正沿着一张放在套子里的唱片上的纹沟走呢。你的手杖刮着墙上的那些奇形怪状的东西，就像是唱针沿着唱片的音纹转，就这样我们听到了音乐。

\dialogname{阿基里斯：}哎呀，这可怎么办？这\ldots\ldots{}

\dialogname{乌龟：}怎么啦？你怎么不高兴？你以前有过这么密切地同音乐接触的机会吗？

\dialogname{阿基里斯：}龟兄，我比个跳蚤还小，怎么还能有劲跑过身高正常的人呢？

\dialogname{乌龟：}噢，就这点让你不安吗？根本不必为此发愁，阿基。

\dialogname{阿基里斯：}看你说话的样子，我觉得你一点也不担心。

\dialogname{乌龟：}我不知道，不过有一点可以肯定，我不担心没劲儿。尤其是在面临遇见大雕这种危险的时候。

\dialogname{阿基里斯：}太可怕了！你是说------

\dialogname{乌龟：}恐怕是的，阿基。我一听那段音乐就知道了\ldots\ldots{}

\dialogname{阿基里斯：}怎么回事？

\dialogname{乌龟：}很简单。当我听到最高声部的B-A-C-H旋律时，我立刻听出我们正在穿行而过的这些音纹只能是《和声小迷宫》------巴赫的一首不太为人所知的管风琴作品。取这个名字是因为它那令人心醉神迷的频繁转调。

\dialogname{阿基里斯：}什------什么玩意儿？

\dialogname{乌龟：}嗯，你知道，大多数乐曲都用一个调或调性写成的。比如C大调，这只曲子就是C大调。

\dialogname{阿基里斯：}我以前听说过这个词。它的意思是不是说你希望它结束在音符C上？

\dialogname{乌龟：}对。在某种意义上说，C起一个``基底''的作用。实际上，通常叫它``主调音''。

\dialogname{阿基里斯：}是不是离开了主调音就是为了最后回到主调音？

\dialogname{乌龟：}正是这样。随着乐曲的发展，运用了模棱两可的和弦旋律，它们都离开了主调音。渐渐地，期待形成了------你越来越渴望回复，听到主调音。

\dialogname{阿基里斯：}这就是为什么我在一支曲子结束时总是感到那么满足，就像我全身心都在等待着听到那个主调音似的，是吗？

\dialogname{乌龟：}没错儿。作曲家运用他的和声进行来控制你的感情，使你产生渴望听到主调音的愿望。

\dialogname{阿基里斯：}可是你该给我讲讲转调。

\dialogname{乌龟：}哦，对。作曲家可以做的一件非常重要的事就是在一支曲子中``转调''，其目的在于建立一个暂时的终结，而不是一种主调音的解决。

\dialogname{阿基里斯：}我想\ldots\ldots 我明白了。你是说某组用不同音调反复演奏的和弦改换了和声的紧张点，这样一来就使我实际上期待着在一个新的调上把不谐和音转变为谐和音？

\dialogname{乌龟：}对极了。这就使情况更复杂了，因为虽然一时你期望在一个新调上把不谐和音转变为谐和音，但在你内心深处一直存在着一个渴望，即要达到原来那个终结------在这支曲子里就是大调的C。而当次要终结到达时，就会有------

\dialogname{阿基里斯}哦，听听这华美的突然上升的和弦，这是《和声小迷宫》结尾的标志。

\dialogname{乌龟：}不，阿基，这不是结尾，这只是------

\dialogname{阿基里斯：}它肯定是！喔！多么强有力的结尾啊！真让人有如释重负之感！这才是个解决呢！妙极了！

（确实，音乐就在这里停止了，他们走出来，到了一块没有墙壁的空地上。）

你瞧，它确实完了。我说什么来着！

\dialogname{乌龟：}真邪了！这张唱片是音乐界的耻辱。

\dialogname{阿基里斯：}什么意思？

\dialogname{乌龟：}就是我刚告诉你的意思。这里巴赫从C调转到G调，建立起听G的第二个目的。这意味着你同时体验到两种张力------等待着向G转变，可同时心里还埋藏着那个最大的渴望------成功地转到C大调。

\dialogname{阿基里斯：}你为什么听一支乐曲时一定要苦思冥想呢？音乐只是一种智力游戏吗？

\dialogname{乌龟：}不，当然不是。某些音乐是极富于智慧的，不过大多数音乐不是。多数的时候你的耳朵或大脑为你``估算''，使你的情感知道它们想听什么，你不必有意识地去想它。但是在这支曲子里，巴赫在耍手腕，想把你引入歧途。而在你身上，阿基，他成功了。

\dialogname{阿基里斯：}你是指我对副调解决的反应？

\dialogname{乌龟：}正是。

\dialogname{阿基里斯：}可我还是觉得它像个结尾。

\dialogname{乌龟：}巴赫有意那么安排的。你恰好落入他的圈套了。它故意被设计得像个结尾，但是如果你仔细品味那个和声的进行，你会发现它的调性不对。显然，不只是你，连这个蹩脚的唱片公司也落入了同样的圈套------他们过早截断了这支曲子！

\dialogname{阿基里斯：}巴赫可把我给涮了！

\dialogname{乌龟：}这就是他的全部花招------让你在迷宫里迷路！大雕跟巴赫是同谋，你瞧，如果你不小心，他现在就会把你给笑死的------没准儿连我也给捎上呢。

\dialogname{阿基里斯：}哦，让我们赶快离开这儿吧！快，咱们沿着音纹往回跑，在大雕发现我们之前跑出这张唱片！

\dialogname{乌龟：}天哪，不行！我的感觉太脆弱，受不了倒序演奏时那古怪的和声进行。

\dialogname{阿基里斯：}哦，龟兄，要是我们不能顺原路回去，我们怎么才能离开这里呢？

\dialogname{乌龟：}是啊。

（阿基里斯有点绝望，开始在黑暗里乱跑。突然出现了一声轻微的喘息，接着``砰''地一声。）

阿基------你怎么了？

\dialogname{阿基里斯：}有点错位，不过还好。我掉进了一个大洞里。

\dialogname{乌龟：}你掉到邪恶大雕的窝里去啦！我在这儿呢，我来帮你出来。我们得赶紧走！

\dialogname{阿基里斯：}小心，龟兄------我不想让你也掉进来\ldots\ldots{}

\dialogname{乌龟：}别怕，阿基，一切都会------

（突然间出现了一声轻微的喘息，接着``砰''地一声。）

\dialogname{阿基里斯：}龟兄------你也掉进来了！你怎么样？

\dialogname{乌龟：}只是伤着了我的自尊心，不过还好。

\dialogname{阿基里斯：}我们崴泥了，是吧？

（突然，一阵隆隆的笑声响起，并吓人地靠近了。）

\dialogname{乌龟：}当心，阿基，这可不是玩笑的事！

\dialogname{大雕：}嘿嘿嘿！呵呵呵！哈哈哈！

\dialogname{阿基里斯：}我开始顶不住了，龟兄\ldots\ldots{}

\dialogname{乌龟：}千万别理它的笑，阿基。这是你唯一的希望。

\dialogname{阿基里斯：}我尽力而为吧。只要我的肚子不那么饿！

\dialogname{乌龟：}哎，我闻到什么啦？附近有一碗弹出锅酥吧？

\dialogname{阿基里斯：}我也闻到了。从哪儿飘来的？

\dialogname{乌龟：}我想就在附近。哦！我碰在一个大碗一样的东西上了。对，没错儿------像是一大碗弹出锅酥。

\dialogname{阿基里斯：}哦，天哪------弹出锅酥！我可得甩开腮帮子足吃！

\end{adjustwidth}

\end{adjustwidth}

\end{adjustwidth}

\dialogname{乌龟：}可别是推入锅酥！推入锅酥和弹出锅酥可真不好分清呀。

\begin{adjustwidth}{2em}{0em}

\begin{adjustwidth}{2em}{0em}

\begin{adjustwidth}{2em}{0em}

\dialogname{阿基里斯：}焚琴？你说什么呢？

\dialogname{乌龟：}我什么也没说呀。你大概是有点风声鹤唳了吧。

\dialogname{阿基里斯：}还煮鹤！我真的不喜欢鹤肉。好啦，我们就吃这个吧！

（两个朋友开始大嚼弹出锅酥（或者推入锅酥？）------可突然------``弹''地一声！我猜他们吃的到底还是弹出锅酥。）

\end{adjustwidth}

\dialogname{乌龟：}多有趣的故事啊。你喜欢它吗？

\dialogname{阿基里斯：}一般。我想知道的只是他们离开那邪恶大雕的窝没有。可怜的阿基里斯------他想恢复正常身高。

\dialogname{乌龟：}别担心------他们出去了，他也恢复了正常身高。这就是刚才为什么有一声``弹''。

\dialogname{阿基里斯：}哦，我没听出来。嗯，我现在的确想找那瓶煮调饮了。不知为什么，我的嘴唇都要裂了。没有什么比弹出煮调饮更好喝的了。

\dialogname{乌龟：}这种饮料以消燥解渴著名。在有些地方，人们对它着了魔。在十九世纪末、二十世纪初的维也纳，勋伯格食品工厂停止生产煮调饮，转产十二瓶菌氯，你无法想象由此而引起的骚乱。

\dialogname{阿基里斯：}我略有所闻。不过还是让我们找找那瓶煮调饮吧，嘿------等一会儿。桌上这些蜥蜴------你发现它们有什么特别的吗？

\dialogname{乌龟：}嗯------没什么特别的。你看见了什么这么有趣？

\dialogname{阿基里斯：}你没看见？它们没喝什么弹出煮调饮就从那张平面上的画里跑出来了。它们怎么做到的？

\dialogname{乌龟：}我没告诉你吗？要是你没有弹出煮调饮的话，你以同画面成直角的方向移动，也能从里面出来。这些小蜥蜴知道，要想从两维的画中世界里出来，就得向上爬。

\dialogname{阿基里斯：}我们要从艾舍尔的这幅画中出来，也得做同样的事吗？

\dialogname{乌龟：}当然！我们只需要上一层。你想试试吗？

\dialogname{阿基里斯：}只要能回到我的房子里！这次吸引人的历险已经使我很累了。

\dialogname{乌龟：}跟着我，从这条路上来。

（他们走上一层。）

\end{adjustwidth}

\dialogname{阿基里斯：}可回来了。真好。不过有点不对劲儿。这不是我的房子！这是你的，龟兄。

\dialogname{乌龟：}唔，是的------我倒愿意这样呢！我一点也不想从你那里再走那么长的路回来了。我累坏了，简直怀疑我是否走了那么远。

\dialogname{阿基里斯：}我不在乎走回家，所以我想我们在这儿结束也挺好。

\dialogname{乌龟：}说得对！这的确是次好运！

\end{adjustwidth}

\phantomsection\label{Chapter05.xhtml}{}

第五章


\end{document}
